\chapter{Data Structures}

\section{Linked Lists}
\setcounter{section}{1}
\setcounter{subsection}{0}
\subsection{Singly Linked List Operations}
Implements common singly linked list operations using raw nodes. Linked lists are rarely the best contest choice compared with arrays, vectors, and \inlinecode{std::list}, but the pointer patterns are useful for interview-style problems and for understanding constant-time splicing.

For production C++, prefer containers such as \inlinecode{std::list}, \inlinecode{std::forward\_list}, or \inlinecode{std::vector} when they fit the problem. Use the raw-node style below when the problem statement already gives nodes, or when manual pointer manipulation is the point of the exercise.

\begin{compactitem}
  \item \inlinecode{reverse\_list(head)} reverses a singly linked list and returns the new head.
  \item \inlinecode{merge\_sorted\_lists(a, b)} merges two sorted lists using a dummy node and returns the merged head.
  \item \inlinecode{split\_half(head, \&second)} cuts a list into two halves, returning the first half and storing the second half in \inlinecode{second}.
\end{compactitem}

For linked-list cycle detection, use Floyd or Brent cycle detection from section 1.5.

\Needspace*{3\baselineskip}
\textbf{Time Complexity}:
\begin{compactitem}
  \item $O(n)$ per call to \inlinecode{reverse\_list(head)} and \inlinecode{split\_half(head, \&second)}.
  \item $O(n + m)$ per call to \inlinecode{merge\_sorted\_lists(a, b)}.
\end{compactitem}

\textbf{Space Complexity}: $O(1)$ auxiliary for all operations.

\begin{lstlisting}[language=C++]
#include <cstddef>

struct ListNode {
  int value;
  ListNode *next;

  explicit ListNode(int value = 0) : value(value), next(nullptr) {}
};

ListNode *reverse_list(ListNode *head) {
  ListNode *prev = nullptr;
  while (head != nullptr) {
    ListNode *next = head->next;
    head->next = prev;
    prev = head;
    head = next;
  }
  return prev;
}

ListNode *merge_sorted_lists(ListNode *a, ListNode *b) {
  ListNode dummy;
  ListNode *tail = &dummy;
  while (a != nullptr && b != nullptr) {
    if (b->value < a->value) {
      tail->next = b;
      b = b->next;
    } else {
      tail->next = a;
      a = a->next;
    }
    tail = tail->next;
  }
  tail->next = a != nullptr ? a : b;
  return dummy.next;
}

ListNode *split_half(ListNode *head, ListNode **second) {
  if (head == nullptr || head->next == nullptr) {
    *second = nullptr;
    return head;
  }
  ListNode *slow = head, *fast = head->next;
  while (fast != nullptr && fast->next != nullptr) {
    slow = slow->next;
    fast = fast->next->next;
  }
  *second = slow->next;
  slow->next = nullptr;
  return head;
}
\end{lstlisting}
\Needspace*{4\baselineskip}
\textbf{Example Usage}\par
\begin{lstlisting}[language=C++]
#include <cassert>
#include <vector>
using namespace std;

vector<int> values(ListNode *head) {
  vector<int> res;
  for (; head != nullptr; head = head->next) {
    res.push_back(head->value);
  }
  return res;
}

int main() {
  vector<ListNode> a{ListNode(1), ListNode(3), ListNode(5)};
  vector<ListNode> b{ListNode(2), ListNode(4), ListNode(6)};
  for (int i = 0; i + 1 < static_cast<int>(a.size()); i++) {
    a[i].next = &a[i + 1];
    b[i].next = &b[i + 1];
  }
  ListNode *merged = merge_sorted_lists(&a[0], &b[0]);
  vector<int> merged_values = values(merged);
  for (int i = 0; i < 6; i++) {
    assert(merged_values[i] == i + 1);
  }

  ListNode *second = nullptr;
  ListNode *first = split_half(merged, &second);
  assert(values(first).size() == 3);
  assert(values(second).size() == 3);
  assert(values(reverse_list(first))[0] == 3);
  return 0;
}
\end{lstlisting}
\subsection{Linked List Splicing}
Moves nodes or ranges between singly linked lists in $O(1)$ pointer changes once the neighboring nodes are known. Splicing is the main operation that makes linked lists useful: it moves existing nodes without copying values or invalidating pointers to the moved nodes.

The functions below use a dummy head node for each list. This makes insertions and removals at the beginning of a list match all other positions. In the C++ standard library, \inlinecode{std::list::splice()} provides analogous operations for doubly linked lists, while \inlinecode{std::forward\_list::splice\_after()} provides analogous operations for singly linked lists.

\begin{compactitem}
  \item \inlinecode{splice\_after(pos, before)} moves the node after \inlinecode{before} so it appears after \inlinecode{pos}.
  \item \inlinecode{splice\_range\_after(pos, before\_first, last)} moves the half-open range \inlinecode{(before\_first, last)} so it appears after \inlinecode{pos}.
\end{compactitem}

\Needspace*{3\baselineskip}
\textbf{Time Complexity}:
\begin{compactitem}
  \item $O(1)$ per call to \inlinecode{splice\_after(pos, before)}.
  \item $O(k)$ per call to \inlinecode{splice\_range\_after(pos, before\_first, last)}, where $k$ is the number of moved nodes. The actual pointer splicing is $O(1)$, but this compact version walks to the end of the moved range.
\end{compactitem}

\textbf{Space Complexity}: $O(1)$ auxiliary.

\begin{lstlisting}[language=C++]
#include <cstddef>

struct ListNode {
  int value;
  ListNode *next;

  explicit ListNode(int value = 0) : value(value), next(nullptr) {}
};

void splice_after(ListNode *pos, ListNode *before) {
  ListNode *node = before->next;
  if (node == nullptr || pos == before || pos == node) {
    return;
  }
  before->next = node->next;
  node->next = pos->next;
  pos->next = node;
}

void splice_range_after(ListNode *pos, ListNode *before_first, ListNode *last) {
  ListNode *first = before_first->next;
  if (first == last) {
    return;
  }
  ListNode *tail = first;
  while (tail->next != last) {
    tail = tail->next;
  }
  before_first->next = last;
  tail->next = pos->next;
  pos->next = first;
}
\end{lstlisting}
\Needspace*{4\baselineskip}
\textbf{Example Usage}\par
\begin{lstlisting}[language=C++]
#include <cassert>
#include <vector>
using namespace std;

vector<int> values(ListNode *dummy) {
  vector<int> res;
  for (ListNode *p = dummy->next; p != nullptr; p = p->next) {
    res.push_back(p->value);
  }
  return res;
}

int main() {
  ListNode dummy(0), a(1), b(2), c(3), d(4), e(5);
  dummy.next = &a;
  a.next = &b;
  b.next = &c;
  c.next = &d;
  d.next = &e;

  splice_after(&dummy, &c);  // Move 4 to the front: 4, 1, 2, 3, 5.
  vector<int> v = values(&dummy);
  vector<int> expected1{4, 1, 2, 3, 5};
  for (int i = 0; i < static_cast<int>(expected1.size()); i++) {
    assert(v[i] == expected1[i]);
  }

  splice_range_after(&d, &a, &e);  // Move 2, 3 after 4: 4, 2, 3, 1, 5.
  v = values(&dummy);
  vector<int> expected2{4, 2, 3, 1, 5};
  for (int i = 0; i < static_cast<int>(expected2.size()); i++) {
    assert(v[i] == expected2[i]);
  }
  return 0;
}
\end{lstlisting}
\subsection{LRU Cache}
Maintains a least-recently-used cache with fixed capacity. The cache supports lookups and updates by key, evicting the least recently used entry whenever an insertion would exceed capacity.

This implementation combines a doubly linked list (\inlinecode{std::list}) with a map from keys to list iterators. The front of the list stores the most recently used item. The operation \inlinecode{items.splice(items.begin(), items, it)} is the standard-library way to move an existing list node to the front in $O(1)$ time.

\begin{compactitem}
  \item \inlinecode{LRUCache(capacity)} constructs an empty cache holding at most \inlinecode{capacity} keys.
  \item \inlinecode{get(key, \&value)} returns whether \inlinecode{key} is present. If present, it stores the value in \inlinecode{value} and marks the key as most recently used.
  \item \inlinecode{put(key, value)} inserts or updates \inlinecode{key}, marking it as most recently used and evicting the least recently used key if needed.
  \item \inlinecode{size()} returns the number of keys currently stored.
\end{compactitem}

\textbf{Time Complexity}: $O(1)$ amortized per call to \inlinecode{get(key, \&value)} and \inlinecode{put(key, value)}.

\textbf{Space Complexity}: $O(n)$, where $n$ is the cache capacity.

\begin{lstlisting}[language=C++]
#include <list>
#include <unordered_map>
#include <utility>

template<class Key, class Value>
class LRUCache {
  using ListIter = typename std::list<std::pair<Key, Value>>::iterator;

  int cap;
  std::list<std::pair<Key, Value>> items;
  std::unordered_map<Key, ListIter> where;

 public:
  explicit LRUCache(int capacity) : cap(capacity) {}

  int size() const { return items.size(); }

  bool get(const Key &key, Value *value) {
    auto it = where.find(key);
    if (it == where.end()) {
      return false;
    }
    items.splice(items.begin(), items, it->second);
    it->second = items.begin();
    *value = it->second->second;
    return true;
  }

  void put(const Key &key, const Value &value) {
    auto it = where.find(key);
    if (it != where.end()) {
      it->second->second = value;
      items.splice(items.begin(), items, it->second);
      it->second = items.begin();
      return;
    }
    if (cap == 0) {
      return;
    }
    if (static_cast<int>(items.size()) == cap) {
      where.erase(items.back().first);
      items.pop_back();
    }
    items.push_front({key, value});
    where[key] = items.begin();
  }
};
\end{lstlisting}
\Needspace*{4\baselineskip}
\textbf{Example Usage}\par
\begin{lstlisting}[language=C++]
#include <cassert>
using namespace std;

int main() {
  LRUCache<int, int> cache(2);
  int value = 0;
  cache.put(1, 10);
  cache.put(2, 20);
  assert(cache.get(1, &value) && value == 10);
  cache.put(3, 30);  // Evicts key 2.
  assert(!cache.get(2, &value));
  assert(cache.get(3, &value) && value == 30);
  cache.put(1, 11);
  assert(cache.get(1, &value) && value == 11);
  return 0;
}
\end{lstlisting}

\section{Heaps and Priority Queues}
\setcounter{section}{2}
\setcounter{subsection}{0}
\subsection{Binary Heap}
Maintain a min-priority queue, that is, a collection of elements with support for querying and extraction of the minimum. This implementation requires an ordering on the set of possible elements defined by \inlinecode{operator\textless{}}. A binary min-heap implements a priority queue by inserting and deleting nodes into a binary tree such that the parent of any node is always less than its children.

\begin{compactitem}
  \item \inlinecode{BinaryHeap()} constructs an empty priority queue.
  \item \inlinecode{BinaryHeap(lo, hi)} constructs a priority queue from two ForwardIterators, consisting of elements in the range \inlinecode{[lo, hi)}.
  \item \inlinecode{size()} returns the size of the priority queue.
  \item \inlinecode{empty()} returns whether the priority queue is empty.
  \item \inlinecode{push(v)} inserts the value \inlinecode{v} into the priority queue.
  \item \inlinecode{pop()} removes the minimum element from the priority queue.
  \item \inlinecode{top()} returns the minimum element in the priority queue.
\end{compactitem}

\Needspace*{3\baselineskip}
\textbf{Time Complexity}:
\begin{compactitem}
  \item $O(1)$ per call to the first constructor, \inlinecode{size()}, \inlinecode{empty()}, and \inlinecode{top()}.
  \item $O(\log n)$ per call to \inlinecode{push()} and \inlinecode{pop()}, where $n$ is the number of elements in the priority queue.
  \item $O(n)$ per call to the second constructor on the distance between \inlinecode{lo} and \inlinecode{hi} (bottom-up heapify).
\end{compactitem}

\Needspace*{3\baselineskip}
\textbf{Space Complexity}:
\begin{compactitem}
  \item $O(n)$ for storage of the priority queue elements.
  \item $O(1)$ auxiliary for all operations.
\end{compactitem}

\begin{lstlisting}[language=C++]
#include <algorithm>
#include <stdexcept>
#include <vector>

template<class T>
class BinaryHeap {
  std::vector<T> heap;

  void sift_down(int i) {
    for (;;) {
      int child = 2 * i + 1;
      if (child >= static_cast<int>(heap.size())) {
        break;
      }
      if (child + 1 < static_cast<int>(heap.size()) && heap[child + 1] < heap[child]) {
        child++;
      }
      if (heap[child] < heap[i]) {
        std::swap(heap[i], heap[child]);
        i = child;
      } else {
        break;
      }
    }
  }

 public:
  BinaryHeap() {}

  template<class It>
  BinaryHeap(It lo, It hi) : heap(lo, hi) {
    for (int i = static_cast<int>(heap.size()) / 2 - 1; i >= 0; i--) {
      sift_down(i);
    }
  }

  int size() const { return heap.size(); }
  bool empty() const { return heap.empty(); }

  void push(const T &v) {
    heap.push_back(v);
    int i = heap.size() - 1;
    while (i > 0) {
      int parent = (i - 1) / 2;
      if (!(heap[i] < heap[parent])) {
        break;
      }
      std::swap(heap[i], heap[parent]);
      i = parent;
    }
  }

  void pop() {
    if (heap.empty()) {
      throw std::runtime_error("Cannot pop from empty heap.");
    }
    heap[0] = heap.back();
    heap.pop_back();
    if (!heap.empty()) {
      sift_down(0);
    }
  }

  T top() const {
    if (heap.empty()) {
      throw std::runtime_error("Cannot get top of empty heap.");
    }
    return heap[0];
  }
};
\end{lstlisting}
\Needspace*{4\baselineskip}
\textbf{Example Usage}\par
\begin{lstlisting}[language=C++]
#include <iostream>
using namespace std;

int main() {
  vector<int> a{0, 5, -1, 12};
  BinaryHeap<int> h(a.begin(), a.end());
  h.push(10);
  while (!h.empty()) {
    cout << h.top() << endl;
    h.pop();
  }
  return 0;
}
\end{lstlisting}
\begin{exampleoutput}
-1
0
5
10
12
\end{exampleoutput}
\subsection{Skew Heap}
Maintain a mergeable min-priority queue, that is, a collection of elements with support for querying and extraction of the minimum as well as efficient merging with other instances. This implementation requires an ordering on the set of possible elements defined by \inlinecode{operator\textless{}}. A skew heap attempts to maintain balance by unconditionally swapping all nodes in the merge path when merging.

\begin{compactitem}
  \item \inlinecode{SkewHeap()} constructs an empty priority queue.
  \item \inlinecode{SkewHeap(lo, hi)} constructs a priority queue from two ForwardIterators, consisting of elements in the range \inlinecode{[lo, hi)}.
  \item \inlinecode{size()} returns the size of the priority queue.
  \item \inlinecode{empty()} returns whether the priority queue is empty.
  \item \inlinecode{push(v)} inserts the value \inlinecode{v} into the priority queue.
  \item \inlinecode{pop()} removes the minimum element from the priority queue.
  \item \inlinecode{top()} returns the minimum element in the priority queue.
  \item \inlinecode{absorb(h)} inserts every value from \inlinecode{h} and sets \inlinecode{h} to the empty priority queue.
\end{compactitem}

\Needspace*{3\baselineskip}
\textbf{Time Complexity}:
\begin{compactitem}
  \item $O(1)$ per call to the first constructor, \inlinecode{size()}, \inlinecode{empty()}, and \inlinecode{top()}.
  \item $O(\log n)$ amortized auxiliary per call to \inlinecode{push()}, \inlinecode{pop()}, and \inlinecode{absorb()}, where $n$ is the number of elements in the priority queue.
  \item $O(n)$ per call to the second constructor, where $n$ is the distance between \inlinecode{lo} and \inlinecode{hi}.
\end{compactitem}

\Needspace*{3\baselineskip}
\textbf{Space Complexity}:
\begin{compactitem}
  \item $O(n)$ for storage of the priority queue elements.
  \item $O(\log n)$ amortized auxiliary stack space for \inlinecode{push()}, \inlinecode{pop()}, and \inlinecode{absorb()}.
  \item $O(1)$ auxiliary for all other operations.
\end{compactitem}

\begin{lstlisting}[language=C++]
#include <algorithm>
#include <cstddef>
#include <stdexcept>

template<class T>
class SkewHeap {
  struct Node {
    T value;
    Node *left, *right;

    explicit Node(const T &v) : value(v), left(nullptr), right(nullptr) {}
  } *root;

  int num_nodes;

  static Node *merge(Node *a, Node *b) {
    if (a == nullptr) {
      return b;
    }
    if (b == nullptr) {
      return a;
    }
    if (b->value < a->value) {
      std::swap(a, b);
    }
    std::swap(a->left, a->right);
    a->left = merge(b, a->left);
    return a;
  }

  static void clean_up(Node *n) {
    if (n != nullptr) {
      clean_up(n->left);
      clean_up(n->right);
      delete n;
    }
  }

 public:
  SkewHeap() : root(nullptr), num_nodes(0) {}

  template<class It>
  SkewHeap(It lo, It hi) : root(nullptr), num_nodes(0) {
    while (lo != hi) {
      push(*(lo++));
    }
  }

  ~SkewHeap() { clean_up(root); }
  int size() const { return num_nodes; }
  bool empty() const { return root == nullptr; }

  void push(const T &v) {
    root = merge(root, new Node(v));
    num_nodes++;
  }

  void pop() {
    if (empty()) {
      throw std::runtime_error("Cannot pop from empty heap.");
    }
    Node *tmp = root;
    root = merge(root->left, root->right);
    delete tmp;
    num_nodes--;
  }

  T top() const {
    if (empty()) {
      throw std::runtime_error("Cannot get top of empty heap.");
    }
    return root->value;
  }

  void absorb(SkewHeap &h) {
    root = merge(root, h.root);
    num_nodes += h.num_nodes;
    h.root = nullptr;
    h.num_nodes = 0;
  }
};
\end{lstlisting}
\Needspace*{4\baselineskip}
\textbf{Example Usage}\par
\begin{lstlisting}[language=C++]
#include <iostream>
using namespace std;

int main() {
  SkewHeap<int> h, h2;
  h.push(12);
  h.push(10);
  h2.push(5);
  h2.push(-1);
  h2.push(0);
  h.absorb(h2);
  while (!h.empty()) {
    cout << h.top() << endl;
    h.pop();
  }
  return 0;
}
\end{lstlisting}
\begin{exampleoutput}
-1
0
5
10
12
\end{exampleoutput}
\subsection{Pairing Heap}
Maintain a mergeable min-priority queue, that is, a collection of elements with support for querying and extraction of the minimum as well as efficient merging with other instances. This implementation requires an ordering on the set of possible elements defined by \inlinecode{operator\textless{}}. A pairing heap is a heap-ordered multi-way tree, using a two-pass merge to self-adjust during each deletion.

\begin{compactitem}
  \item \inlinecode{PairingHeap()} constructs an empty priority queue.
  \item \inlinecode{PairingHeap(lo, hi)} constructs a priority queue from two ForwardIterators, consisting of elements in the range \inlinecode{[lo, hi)}.
  \item \inlinecode{size()} returns the size of the priority queue.
  \item \inlinecode{empty()} returns whether the priority queue is empty.
  \item \inlinecode{push(v)} inserts the value \inlinecode{v} into the priority queue.
  \item \inlinecode{pop()} removes the minimum element from the priority queue.
  \item \inlinecode{top()} returns the minimum element in the priority queue.
  \item \inlinecode{absorb(h)} inserts every value from \inlinecode{h} and sets \inlinecode{h} to the empty priority queue.
\end{compactitem}

\Needspace*{3\baselineskip}
\textbf{Time Complexity}:
\begin{compactitem}
  \item $O(1)$ per call to the first constructor, \inlinecode{size()}, \inlinecode{empty()}, \inlinecode{top()}, \inlinecode{push()}, and \inlinecode{absorb()}.
  \item $O(\log n)$ amortized per call to \inlinecode{pop()}.
  \item $O(n)$ per call to the second constructor on the distance between \inlinecode{lo} and \inlinecode{hi}.
\end{compactitem}

\Needspace*{3\baselineskip}
\textbf{Space Complexity}:
\begin{compactitem}
  \item $O(n)$ for storage of the priority queue elements.
  \item $O(\log n)$ amortized auxiliary stack space for \inlinecode{pop()}.
  \item $O(1)$ auxiliary for all other operations.
\end{compactitem}

\begin{lstlisting}[language=C++]
#include <cstddef>
#include <stdexcept>

template<class T>
class PairingHeap {
  struct Node {
    T value;
    Node *left, *next;

    explicit Node(const T &v) : value(v), left(nullptr), next(nullptr) {}

    void add_child(Node *n) {
      if (left == nullptr) {
        left = n;
      } else {
        n->next = left;
        left = n;
      }
    }
  } *root;

  int num_nodes;

  static Node *merge(Node *a, Node *b) {
    if (a == nullptr) {
      return b;
    }
    if (b == nullptr) {
      return a;
    }
    if (a->value < b->value) {
      a->add_child(b);
      return a;
    }
    b->add_child(a);
    return b;
  }

  static Node *merge_pairs(Node *n) {
    if (n == nullptr || n->next == nullptr) {
      return n;
    }
    Node *a = n, *b = n->next, *c = n->next->next;
    a->next = b->next = nullptr;
    return merge(merge(a, b), merge_pairs(c));
  }

  static void clean_up(Node *n) {
    if (n != nullptr) {
      clean_up(n->left);
      clean_up(n->next);
      delete n;
    }
  }

 public:
  PairingHeap() : root(nullptr), num_nodes(0) {}

  template<class It>
  PairingHeap(It lo, It hi) : root(nullptr), num_nodes(0) {
    while (lo != hi) {
      push(*(lo++));
    }
  }

  ~PairingHeap() { clean_up(root); }
  int size() const { return num_nodes; }
  bool empty() const { return root == nullptr; }

  void push(const T &v) {
    root = merge(root, new Node(v));
    num_nodes++;
  }

  void pop() {
    if (empty()) {
      throw std::runtime_error("Cannot pop from empty heap.");
    }
    Node *tmp = root;
    root = merge_pairs(root->left);
    delete tmp;
    num_nodes--;
  }

  T top() const {
    if (empty()) {
      throw std::runtime_error("Cannot get top of empty heap.");
    }
    return root->value;
  }

  void absorb(PairingHeap &h) {
    root = merge(root, h.root);
    num_nodes += h.num_nodes;
    h.root = nullptr;
    h.num_nodes = 0;
  }
};
\end{lstlisting}
\Needspace*{4\baselineskip}
\textbf{Example Usage}\par
\begin{lstlisting}[language=C++]
#include <iostream>
using namespace std;

int main() {
  PairingHeap<int> h, h2;
  h.push(12);
  h.push(10);
  h2.push(5);
  h2.push(-1);
  h2.push(0);
  h.absorb(h2);
  while (!h.empty()) {
    cout << h.top() << endl;
    h.pop();
  }
  return 0;
}
\end{lstlisting}
\begin{exampleoutput}
-1
0
5
10
12
\end{exampleoutput}

\section{Dictionaries and Ordered Sets}
\setcounter{section}{3}
\setcounter{subsection}{0}
\subsection{Binary Search Tree}
Maintain a map, that is, a collection of key-value pairs such that each possible key appears at most once in the collection. This implementations requires an ordering on the set of possible keys defined by \inlinecode{operator\textless{}} on the key type. A binary search tree implements this map by inserting and deleting keys into a binary tree such that every node's left child has a lesser key and every node's right child has a greater key.

\begin{compactitem}
  \item \inlinecode{BST()} constructs an empty map.
  \item \inlinecode{size()} returns the size of the map.
  \item \inlinecode{empty()} returns whether the map is empty.
  \item \inlinecode{insert(k, v)} adds an entry with key \inlinecode{k} and value \inlinecode{v} to the map, returning \inlinecode{true} if an new entry was added or \inlinecode{false} if the key already exists (in which case the map is unchanged and the old value associated with the key is preserved).
  \item \inlinecode{erase(k)} removes the entry with key \inlinecode{k} from the map, returning \inlinecode{true} if the removal was successful or \inlinecode{false} if the key to be removed was not found.
  \item \inlinecode{find(k)} returns a pointer to a const value associated with key \inlinecode{k}, or \inlinecode{nullptr} if the key was not found.
  \item \inlinecode{walk(f)} calls the function \inlinecode{f(k, v)} on each entry of the map, in ascending order of keys.
\end{compactitem}

\Needspace*{3\baselineskip}
\textbf{Time Complexity}:
\begin{compactitem}
  \item $O(1)$ per call to the constructor, \inlinecode{size()}, and \inlinecode{empty()}.
  \item $O(n)$ per call to \inlinecode{insert()}, \inlinecode{erase()}, \inlinecode{find()}, and \inlinecode{walk()}, where $n$ is the number of nodes currently in the map.
\end{compactitem}

\Needspace*{3\baselineskip}
\textbf{Space Complexity}:
\begin{compactitem}
  \item $O(n)$ for storage of the map elements.
  \item $O(n)$ auxiliary stack space for \inlinecode{insert()}, \inlinecode{erase()}, and \inlinecode{walk()}.
  \item $O(1)$ auxiliary for all other operations.
\end{compactitem}

\begin{lstlisting}[language=C++]
#include <cstddef>

template<class K, class V>
class BST {
  struct Node {
    K key;
    V value;
    Node *left, *right;

    Node(const K &k, const V &v) : key(k), value(v), left(nullptr), right(nullptr) {}
  } *root;

  int num_nodes;

  static bool insert(Node *&n, const K &k, const V &v) {
    if (n == nullptr) {
      n = new Node(k, v);
      return true;
    }
    if (k < n->key) {
      return insert(n->left, k, v);
    } else if (n->key < k) {
      return insert(n->right, k, v);
    }
    return false;
  }

  static bool erase(Node *&n, const K &k) {
    if (n == nullptr) {
      return false;
    }
    if (k < n->key) {
      return erase(n->left, k);
    } else if (n->key < k) {
      return erase(n->right, k);
    }
    if (n->left != nullptr && n->right != nullptr) {
      Node *tmp = n->right, *parent = nullptr;
      while (tmp->left != nullptr) {
        parent = tmp;
        tmp = tmp->left;
      }
      n->key = tmp->key;
      n->value = tmp->value;
      if (parent != nullptr) {
        return erase(parent->left, parent->left->key);
      }
      return erase(n->right, n->right->key);
    }
    Node *tmp = (n->left != nullptr) ? n->left : n->right;
    delete n;
    n = tmp;
    return true;
  }

  template<class Fn>
  static void walk(Node *n, Fn f) {
    if (n != nullptr) {
      walk(n->left, f);
      f(n->key, n->value);
      walk(n->right, f);
    }
  }

  static void clean_up(Node *n) {
    if (n != nullptr) {
      clean_up(n->left);
      clean_up(n->right);
      delete n;
    }
  }

 public:
  BST() : root(nullptr), num_nodes(0) {}

  ~BST() { clean_up(root); }
  int size() const { return num_nodes; }
  bool empty() const { return root == nullptr; }

  bool insert(const K &k, const V &v) {
    if (insert(root, k, v)) {
      num_nodes++;
      return true;
    }
    return false;
  }

  bool erase(const K &k) {
    if (erase(root, k)) {
      num_nodes--;
      return true;
    }
    return false;
  }

  const V *find(const K &k) const {
    Node *n = root;
    while (n != nullptr) {
      if (k < n->key) {
        n = n->left;
      } else if (n->key < k) {
        n = n->right;
      } else {
        return &(n->value);
      }
    }
    return nullptr;
  }

  template<class Fn>
  void walk(Fn f) const {
    walk(root, f);
  }
};
\end{lstlisting}
\Needspace*{4\baselineskip}
\textbf{Example Usage}\par
\begin{lstlisting}[language=C++]
#include <cassert>
#include <iostream>
using namespace std;

int main() {
  BST<int, char> t;
  t.insert(2, 'b');
  t.insert(1, 'a');
  t.insert(3, 'c');
  t.insert(5, 'e');
  assert(t.insert(4, 'd'));
  assert(*t.find(4) == 'd');
  assert(!t.insert(4, 'd'));
  auto printch = [](int k, char v) { cout << v; };
  t.walk(printch);
  cout << endl;
  assert(t.erase(1));
  assert(!t.erase(1));
  assert(t.find(1) == nullptr);
  t.walk(printch);
  cout << endl;
  return 0;
}
\end{lstlisting}
\begin{exampleoutput}
abcde
bcde
\end{exampleoutput}
\subsection{Treap}
Maintain a map, that is, a collection of key-value pairs such that each possible key appears at most once in the collection. This implementations requires an ordering on the set of possible keys defined by \inlinecode{operator\textless{}} on the key type. A Treap is a binary search tree where each node holds a randomly assigned priority. The tree satisfies the BST property on keys and the min-heap property on priorities (lower priority value is closer to root). Since priorities are random, this keeps the tree balanced with high probability, making insertions and deletions run in $O(\log n)$.

\begin{compactitem}
  \item \inlinecode{Treap()} constructs an empty map.
  \item \inlinecode{size()} returns the size of the map.
  \item \inlinecode{empty()} returns whether the map is empty.
  \item \inlinecode{insert(k, v)} adds an entry with key \inlinecode{k} and value \inlinecode{v} to the map, returning \inlinecode{true} if an new entry was added or \inlinecode{false} if the key already exists (in which case the map is unchanged and the old value associated with the key is preserved).
  \item \inlinecode{erase(k)} removes the entry with key \inlinecode{k} from the map, returning \inlinecode{true} if the removal was successful or \inlinecode{false} if the key to be removed was not found.
  \item \inlinecode{find(k)} returns a pointer to a const value associated with key \inlinecode{k}, or \inlinecode{nullptr} if the key was not found.
  \item \inlinecode{walk(f)} calls the function \inlinecode{f(k, v)} on each entry of the map, in ascending order of keys.
\end{compactitem}

\Needspace*{3\baselineskip}
\textbf{Time Complexity}:
\begin{compactitem}
  \item $O(1)$ per call to the constructor, \inlinecode{size()}, and \inlinecode{empty()}.
  \item $O(\log n)$ on average per call to \inlinecode{insert()}, \inlinecode{erase()}, and \inlinecode{find()}, where $n$ is the number of entries currently in the map.
  \item $O(n)$ per call to \inlinecode{walk()}.
\end{compactitem}

\Needspace*{3\baselineskip}
\textbf{Space Complexity}:
\begin{compactitem}
  \item $O(n)$ for storage of the map elements.
  \item $O(\log n)$ auxiliary stack space on average for \inlinecode{insert()}, \inlinecode{erase()}, and \inlinecode{walk()}.
  \item $O(1)$ auxiliary for all other operations.
\end{compactitem}

\begin{lstlisting}[language=C++]
#include <random>

template<class K, class V>
class Treap {
  struct Node {
    static inline int rand32() {
      static std::mt19937 rng(std::random_device{}());
      return static_cast<int>(rng());
    }

    K key;
    V value;
    int priority;
    Node *left, *right;

    Node(const K &k, const V &v)
        : key(k), value(v), priority(rand32()), left(nullptr), right(nullptr) {}
  } *root;

  int num_nodes;

  static void rotate_left(Node *&n) {
    Node *tmp = n;
    n = n->right;
    tmp->right = n->left;
    n->left = tmp;
  }

  static void rotate_right(Node *&n) {
    Node *tmp = n;
    n = n->left;
    tmp->left = n->right;
    n->right = tmp;
  }

  static bool insert(Node *&n, const K &k, const V &v) {
    if (n == nullptr) {
      n = new Node(k, v);
      return true;
    }
    if (k < n->key && insert(n->left, k, v)) {
      if (n->left->priority < n->priority) {
        rotate_right(n);
      }
      return true;
    }
    if (n->key < k && insert(n->right, k, v)) {
      if (n->right->priority < n->priority) {
        rotate_left(n);
      }
      return true;
    }
    return false;
  }

  static bool erase(Node *&n, const K &k) {
    if (n == nullptr) {
      return false;
    }
    if (k < n->key) {
      return erase(n->left, k);
    } else if (n->key < k) {
      return erase(n->right, k);
    }
    if (n->left != nullptr && n->right != nullptr) {
      if (n->left->priority < n->right->priority) {
        rotate_right(n);
        return erase(n->right, k);
      }
      rotate_left(n);
      return erase(n->left, k);
    }
    Node *tmp = (n->left != nullptr) ? n->left : n->right;
    delete n;
    n = tmp;
    return true;
  }

  template<class Fn>
  static void walk(Node *n, Fn f) {
    if (n != nullptr) {
      walk(n->left, f);
      f(n->key, n->value);
      walk(n->right, f);
    }
  }

  static void clean_up(Node *n) {
    if (n != nullptr) {
      clean_up(n->left);
      clean_up(n->right);
      delete n;
    }
  }

 public:
  Treap() : root(nullptr), num_nodes(0) {}

  ~Treap() { clean_up(root); }
  int size() const { return num_nodes; }
  bool empty() const { return root == nullptr; }

  bool insert(const K &k, const V &v) {
    if (insert(root, k, v)) {
      num_nodes++;
      return true;
    }
    return false;
  }

  bool erase(const K &k) {
    if (erase(root, k)) {
      num_nodes--;
      return true;
    }
    return false;
  }

  const V *find(const K &k) const {
    Node *n = root;
    while (n != nullptr) {
      if (k < n->key) {
        n = n->left;
      } else if (n->key < k) {
        n = n->right;
      } else {
        return &(n->value);
      }
    }
    return nullptr;
  }

  template<class Fn>
  void walk(Fn f) const {
    walk(root, f);
  }
};
\end{lstlisting}
\Needspace*{4\baselineskip}
\textbf{Example Usage}\par
\begin{lstlisting}[language=C++]
#include <cassert>
#include <iostream>
using namespace std;

int main() {
  Treap<int, char> t;
  t.insert(2, 'b');
  t.insert(1, 'a');
  t.insert(3, 'c');
  t.insert(5, 'e');
  assert(t.insert(4, 'd'));
  assert(*t.find(4) == 'd');
  assert(!t.insert(4, 'd'));
  auto printch = [](int k, char v) { cout << v; };
  t.walk(printch);
  cout << endl;
  assert(t.erase(1));
  assert(!t.erase(1));
  assert(t.find(1) == nullptr);
  t.walk(printch);
  cout << endl;
  return 0;
}
\end{lstlisting}
\begin{exampleoutput}
abcde
bcde
\end{exampleoutput}
\subsection{AVL Tree}
Maintain a map, that is, a collection of key-value pairs such that each possible key appears at most once in the collection. This implementation requires an ordering on the set of possible keys defined by \inlinecode{operator\textless{}} on the key type. An AVL tree is a binary search tree balanced by height, guaranteeing $O(\log n)$ worst-case running time in insertions and deletions by making sure that the heights of the left and right subtrees at every node differ by at most 1.

\begin{compactitem}
  \item \inlinecode{AVLTree()} constructs an empty map.
  \item \inlinecode{size()} returns the size of the map.
  \item \inlinecode{empty()} returns whether the map is empty.
  \item \inlinecode{insert(k, v)} adds an entry with key \inlinecode{k} and value \inlinecode{v} to the map, returning \inlinecode{true} if an new entry was added or \inlinecode{false} if the key already exists (in which case the map is unchanged and the old value associated with the key is preserved).
  \item \inlinecode{erase(k)} removes the entry with key \inlinecode{k} from the map, returning \inlinecode{true} if the removal was successful or \inlinecode{false} if the key to be removed was not found.
  \item \inlinecode{find(k)} returns a pointer to a const value associated with key \inlinecode{k}, or \inlinecode{nullptr} if the key was not found.
  \item \inlinecode{walk(f)} calls the function \inlinecode{f(k, v)} on each entry of the map, in ascending order of keys.
\end{compactitem}

\Needspace*{3\baselineskip}
\textbf{Time Complexity}:
\begin{compactitem}
  \item $O(1)$ per call to the constructor, \inlinecode{size()}, and \inlinecode{empty()}.
  \item $O(\log n)$ per call to \inlinecode{insert()}, \inlinecode{erase()}, and \inlinecode{find()}, where $n$ is the number of entries currently in the map.
  \item $O(n)$ per call to \inlinecode{walk()}.
\end{compactitem}

\Needspace*{3\baselineskip}
\textbf{Space Complexity}:
\begin{compactitem}
  \item $O(n)$ for storage of the map elements.
  \item $O(\log n)$ auxiliary stack space for \inlinecode{insert()}, \inlinecode{erase()}, and \inlinecode{walk()}.
  \item $O(1)$ auxiliary for all other operations.
\end{compactitem}

\begin{lstlisting}[language=C++]
#include <algorithm>
#include <cstddef>

template<class K, class V>
class AVLTree {
  struct Node {
    K key;
    V value;
    int height;
    Node *left, *right;

    Node(const K &k, const V &v) : key(k), value(v), height(1), left(nullptr), right(nullptr) {}
  } *root;

  int num_nodes;

  static int height(Node *n) { return (n != nullptr) ? n->height : 0; }

  static void update_height(Node *n) {
    if (n != nullptr) {
      n->height = 1 + std::max(height(n->left), height(n->right));
    }
  }

  static void rotate_left(Node *&n) {
    Node *tmp = n;
    n = n->right;
    tmp->right = n->left;
    n->left = tmp;
    update_height(tmp);
    update_height(n);
  }

  static void rotate_right(Node *&n) {
    Node *tmp = n;
    n = n->left;
    tmp->left = n->right;
    n->right = tmp;
    update_height(tmp);
    update_height(n);
  }

  static int balance_factor(Node *n) {
    return (n != nullptr) ? (height(n->left) - height(n->right)) : 0;
  }

  static void rebalance(Node *&n) {
    if (n == nullptr) {
      return;
    }
    update_height(n);
    int bf = balance_factor(n);
    if (bf > 1 && balance_factor(n->left) >= 0) {
      rotate_right(n);
    } else if (bf > 1 && balance_factor(n->left) < 0) {
      rotate_left(n->left);
      rotate_right(n);
    } else if (bf < -1 && balance_factor(n->right) <= 0) {
      rotate_left(n);
    } else if (bf < -1 && balance_factor(n->right) > 0) {
      rotate_right(n->right);
      rotate_left(n);
    }
  }

  static bool insert(Node *&n, const K &k, const V &v) {
    if (n == nullptr) {
      n = new Node(k, v);
      return true;
    }
    if ((k < n->key && insert(n->left, k, v)) || (n->key < k && insert(n->right, k, v))) {
      rebalance(n);
      return true;
    }
    return false;
  }

  static bool erase(Node *&n, const K &k) {
    if (n == nullptr) {
      return false;
    }
    if (!(k < n->key || n->key < k)) {
      if (n->left != nullptr && n->right != nullptr) {
        Node *tmp = n->right, *parent = nullptr;
        while (tmp->left != nullptr) {
          parent = tmp;
          tmp = tmp->left;
        }
        n->key = tmp->key;
        n->value = tmp->value;
        if (parent != nullptr) {
          if (!erase(parent->left, parent->left->key)) {
            return false;
          }
        } else if (!erase(n->right, n->right->key)) {
          return false;
        }
      } else {
        Node *tmp = (n->left != nullptr) ? n->left : n->right;
        delete n;
        n = tmp;
      }
      rebalance(n);
      return true;
    }
    if ((k < n->key && erase(n->left, k)) || (n->key < k && erase(n->right, k))) {
      rebalance(n);
      return true;
    }
    return false;
  }

  template<class Fn>
  static void walk(Node *n, Fn f) {
    if (n != nullptr) {
      walk(n->left, f);
      f(n->key, n->value);
      walk(n->right, f);
    }
  }

  static void clean_up(Node *n) {
    if (n != nullptr) {
      clean_up(n->left);
      clean_up(n->right);
      delete n;
    }
  }

 public:
  AVLTree() : root(nullptr), num_nodes(0) {}

  ~AVLTree() { clean_up(root); }
  int size() const { return num_nodes; }
  bool empty() const { return root == nullptr; }

  bool insert(const K &k, const V &v) {
    if (insert(root, k, v)) {
      num_nodes++;
      return true;
    }
    return false;
  }

  bool erase(const K &k) {
    if (erase(root, k)) {
      num_nodes--;
      return true;
    }
    return false;
  }

  const V *find(const K &k) const {
    Node *n = root;
    while (n != nullptr) {
      if (k < n->key) {
        n = n->left;
      } else if (n->key < k) {
        n = n->right;
      } else {
        return &(n->value);
      }
    }
    return nullptr;
  }

  template<class Fn>
  void walk(Fn f) const {
    walk(root, f);
  }
};
\end{lstlisting}
\Needspace*{4\baselineskip}
\textbf{Example Usage}\par
\begin{lstlisting}[language=C++]
#include <cassert>
#include <iostream>
using namespace std;

int main() {
  AVLTree<int, char> t;
  t.insert(2, 'b');
  t.insert(1, 'a');
  t.insert(3, 'c');
  t.insert(5, 'e');
  assert(t.insert(4, 'd'));
  assert(*t.find(4) == 'd');
  assert(!t.insert(4, 'd'));
  auto printch = [](int k, char v) { cout << v; };
  t.walk(printch);
  cout << endl;
  assert(t.erase(1));
  assert(!t.erase(1));
  assert(t.find(1) == nullptr);
  t.walk(printch);
  cout << endl;
  return 0;
}
\end{lstlisting}
\begin{exampleoutput}
abcde
bcde
\end{exampleoutput}
\subsection{Red-Black Tree}
Maintain a map, that is, a collection of key-value pairs such that each possible key appears at most once in the collection. This implementation requires an ordering on the set of possible keys defined by \inlinecode{operator\textless{}} on the key type. A red black tree is a binary search tree balanced by coloring its nodes red or black, then constraining node colors on any simple path from the root to a leaf.

\begin{compactitem}
  \item \inlinecode{RedBlackTree()} constructs an empty map.
  \item \inlinecode{size()} returns the size of the map.
  \item \inlinecode{empty()} returns whether the map is empty.
  \item \inlinecode{insert(k, v)} adds an entry with key \inlinecode{k} and value \inlinecode{v} to the map, returning \inlinecode{true} if an new entry was added or \inlinecode{false} if the key already exists (in which case the map is unchanged and the old value associated with the key is preserved).
  \item \inlinecode{erase(k)} removes the entry with key \inlinecode{k} from the map, returning \inlinecode{true} if the removal was successful or \inlinecode{false} if the key to be removed was not found.
  \item \inlinecode{find(k)} returns a pointer to a const value associated with key \inlinecode{k}, or \inlinecode{nullptr} if the key was not found.
  \item \inlinecode{walk(f)} calls the function \inlinecode{f(k, v)} on each entry of the map, in ascending order of keys.
\end{compactitem}

\Needspace*{3\baselineskip}
\textbf{Time Complexity}:
\begin{compactitem}
  \item $O(1)$ per call to the constructor, \inlinecode{size()}, and \inlinecode{empty()}.
  \item $O(\log n)$ per call to \inlinecode{insert()}, \inlinecode{erase()}, and \inlinecode{find()}, where $n$ is the number of entries currently in the map.
  \item $O(n)$ per call to \inlinecode{walk()}.
\end{compactitem}

\Needspace*{3\baselineskip}
\textbf{Space Complexity}:
\begin{compactitem}
  \item $O(n)$ for storage of the map elements.
  \item $O(\log n)$ auxiliary stack space for \inlinecode{walk()}.
  \item $O(1)$ auxiliary for all other operations.
\end{compactitem}

\begin{lstlisting}[language=C++]
#include <algorithm>
#include <cstddef>

template<class K, class V>
class RedBlackTree {
  enum Color { RED, BLACK };
  struct Node {
    K key;
    V value;
    Color color;
    Node *left, *right, *parent;

    Node(const K &k, const V &v, Color c)
        : key(k), value(v), color(c), left(nullptr), right(nullptr), parent(nullptr) {}
  } *root, *LEAF_NIL;

  int num_nodes;

  void rotate_left(Node *n) {
    Node *tmp = n->right;
    if ((n->right = tmp->left) != LEAF_NIL) {
      n->right->parent = n;
    }
    if ((tmp->parent = n->parent) == LEAF_NIL) {
      root = tmp;
    } else if (n->parent->left == n) {
      n->parent->left = tmp;
    } else {
      n->parent->right = tmp;
    }
    tmp->left = n;
    n->parent = tmp;
  }

  void rotate_right(Node *n) {
    Node *tmp = n->left;
    if ((n->left = tmp->right) != LEAF_NIL) {
      n->left->parent = n;
    }
    if ((tmp->parent = n->parent) == LEAF_NIL) {
      root = tmp;
    } else if (n->parent->right == n) {
      n->parent->right = tmp;
    } else {
      n->parent->left = tmp;
    }
    tmp->right = n;
    n->parent = tmp;
  }

  void insert_fix(Node *n) {
    while (n->parent->color == RED) {
      Node *parent = n->parent;
      Node *grandparent = n->parent->parent;
      if (parent == grandparent->left) {
        Node *uncle = grandparent->right;
        if (uncle->color == RED) {
          grandparent->color = RED;
          parent->color = BLACK;
          uncle->color = BLACK;
          n = grandparent;
        } else {
          if (n == parent->right) {
            rotate_left(parent);
            n = parent;
            parent = n->parent;
          }
          rotate_right(grandparent);
          std::swap(parent->color, grandparent->color);
          n = parent;
        }
      } else if (parent == grandparent->right) {
        Node *uncle = grandparent->left;
        if (uncle->color == RED) {
          grandparent->color = RED;
          parent->color = BLACK;
          uncle->color = BLACK;
          n = grandparent;
        } else {
          if (n == parent->left) {
            rotate_right(parent);
            n = parent;
            parent = n->parent;
          }
          rotate_left(grandparent);
          std::swap(parent->color, grandparent->color);
          n = parent;
        }
      }
    }
    root->color = BLACK;
  }

  void replace(Node *n, Node *replacement) {
    if (n->parent == LEAF_NIL) {
      root = replacement;
    } else if (n == n->parent->left) {
      n->parent->left = replacement;
    } else {
      n->parent->right = replacement;
    }
    replacement->parent = n->parent;
  }

  void erase_fix(Node *n) {
    while (n != root && n->color == BLACK) {
      Node *parent = n->parent;
      if (n == parent->left) {
        Node *sibling = parent->right;
        if (sibling->color == RED) {
          sibling->color = BLACK;
          parent->color = RED;
          rotate_left(parent);
          sibling = parent->right;
        }
        if (sibling->left->color == BLACK && sibling->right->color == BLACK) {
          sibling->color = RED;
          n = parent;
        } else {
          if (sibling->right->color == BLACK) {
            sibling->left->color = BLACK;
            sibling->color = RED;
            rotate_right(sibling);
            sibling = parent->right;
          }
          sibling->color = parent->color;
          parent->color = BLACK;
          sibling->right->color = BLACK;
          rotate_left(parent);
          n = root;
        }
      } else {
        Node *sibling = parent->left;
        if (sibling->color == RED) {
          sibling->color = BLACK;
          parent->color = RED;
          rotate_right(parent);
          sibling = parent->left;
        }
        if (sibling->left->color == BLACK && sibling->right->color == BLACK) {
          sibling->color = RED;
          n = parent;
        } else {
          if (sibling->left->color == BLACK) {
            sibling->right->color = BLACK;
            sibling->color = RED;
            rotate_left(sibling);
            sibling = parent->left;
          }
          sibling->color = parent->color;
          parent->color = BLACK;
          sibling->left->color = BLACK;
          rotate_right(parent);
          n = root;
        }
      }
    }
    n->color = BLACK;
  }

  template<class Fn>
  void walk(Node *n, Fn f) const {
    if (n != LEAF_NIL) {
      walk(n->left, f);
      f(n->key, n->value);
      walk(n->right, f);
    }
  }

  void clean_up(Node *n) {
    if (n != LEAF_NIL) {
      clean_up(n->left);
      clean_up(n->right);
      delete n;
    }
  }

 public:
  RedBlackTree() : num_nodes(0) { root = LEAF_NIL = new Node(K(), V(), BLACK); }

  ~RedBlackTree() {
    clean_up(root);
    delete LEAF_NIL;
  }

  int size() const { return num_nodes; }
  bool empty() const { return num_nodes == 0; }

  bool insert(const K &k, const V &v) {
    Node *curr = root, *prev = LEAF_NIL;
    while (curr != LEAF_NIL) {
      prev = curr;
      if (k < curr->key) {
        curr = curr->left;
      } else if (curr->key < k) {
        curr = curr->right;
      } else {
        return false;
      }
    }
    Node *n = new Node(k, v, RED);
    n->parent = prev;
    if (prev == LEAF_NIL) {
      root = n;
    } else if (k < prev->key) {
      prev->left = n;
    } else {
      prev->right = n;
    }
    n->left = n->right = LEAF_NIL;
    insert_fix(n);
    num_nodes++;
    return true;
  }

  bool erase(const K &k) {
    Node *n = root;
    while (n != LEAF_NIL) {
      if (k < n->key) {
        n = n->left;
      } else if (n->key < k) {
        n = n->right;
      } else {
        break;
      }
    }
    if (n == LEAF_NIL) {
      return false;
    }
    Color color = n->color;
    Node *replacement;
    if (n->left == LEAF_NIL) {
      replacement = n->right;
      replace(n, n->right);
    } else if (n->right == LEAF_NIL) {
      replacement = n->left;
      replace(n, n->left);
    } else {
      Node *tmp = n->right;
      while (tmp->left != LEAF_NIL) {
        tmp = tmp->left;
      }
      color = tmp->color;
      replacement = tmp->right;
      if (tmp->parent == n) {
        replacement->parent = tmp;
      } else {
        replace(tmp, tmp->right);
        tmp->right = n->right;
        tmp->right->parent = tmp;
      }
      replace(n, tmp);
      tmp->left = n->left;
      tmp->left->parent = tmp;
      tmp->color = n->color;
    }
    delete n;
    if (color == BLACK) {
      erase_fix(replacement);
    }
    return true;
  }

  const V *find(const K &k) const {
    Node *n = root;
    while (n != LEAF_NIL) {
      if (k < n->key) {
        n = n->left;
      } else if (n->key < k) {
        n = n->right;
      } else {
        return &(n->value);
      }
    }
    return nullptr;
  }

  template<class Fn>
  void walk(Fn f) const {
    walk(root, f);
  }
};
\end{lstlisting}
\Needspace*{4\baselineskip}
\textbf{Example Usage}\par
\begin{lstlisting}[language=C++]
#include <cassert>
#include <iostream>
using namespace std;

int main() {
  RedBlackTree<int, char> t;
  t.insert(2, 'b');
  t.insert(1, 'a');
  t.insert(3, 'c');
  t.insert(5, 'e');
  assert(t.insert(4, 'd'));
  assert(*t.find(4) == 'd');
  assert(!t.insert(4, 'd'));
  auto printch = [](int k, char v) { cout << v; };
  t.walk(printch);
  cout << endl;
  assert(t.erase(1));
  assert(!t.erase(1));
  assert(t.find(1) == nullptr);
  t.walk(printch);
  cout << endl;
  return 0;
}
\end{lstlisting}
\begin{exampleoutput}
abcde
bcde
\end{exampleoutput}
\subsection{Splay Tree}
Maintain a map, that is, a collection of key-value pairs such that each possible key appears at most once in the collection. This implementation requires an ordering on the set of possible keys defined by \inlinecode{operator\textless{}} on the key type. A splay tree is a balanced binary search tree with the additional property that recently accessed elements are quick to access again.

\begin{compactitem}
  \item \inlinecode{SplayTree()} constructs an empty map.
  \item \inlinecode{size()} returns the size of the map.
  \item \inlinecode{empty()} returns whether the map is empty.
  \item \inlinecode{insert(k, v)} adds an entry with key \inlinecode{k} and value \inlinecode{v} to the map, returning \inlinecode{true} if an new entry was added or \inlinecode{false} if the key already exists (in which case the map is unchanged and the old value associated with the key is preserved).
  \item \inlinecode{erase(k)} removes the entry with key \inlinecode{k} from the map, returning \inlinecode{true} if the removal was successful or \inlinecode{false} if the key to be removed was not found.
  \item \inlinecode{find(k)} returns a pointer to a const value associated with key \inlinecode{k}, or \inlinecode{nullptr} if the key was not found.
  \item \inlinecode{walk(f)} calls the function \inlinecode{f(k, v)} on each entry of the map, in ascending order of keys.
\end{compactitem}

\Needspace*{3\baselineskip}
\textbf{Time Complexity}:
\begin{compactitem}
  \item $O(1)$ per call to the constructor, \inlinecode{size()}, and \inlinecode{empty()}.
  \item $O(\log n)$ per call to \inlinecode{insert()}, \inlinecode{erase()}, and \inlinecode{find()}, where $n$ is the number of entries currently in the map.
  \item $O(n)$ per call to \inlinecode{walk()}.
\end{compactitem}

\Needspace*{3\baselineskip}
\textbf{Space Complexity}:
\begin{compactitem}
  \item $O(n)$ for storage of the map elements.
  \item $O(\log n)$ auxiliary stack space for \inlinecode{insert()}, \inlinecode{erase()}, and \inlinecode{walk()}.
  \item $O(1)$ auxiliary for all other operations.
\end{compactitem}

\begin{lstlisting}[language=C++]
#include <cstddef>

template<class K, class V>
class SplayTree {
  struct Node {
    K key;
    V value;
    Node *left, *right;

    Node(const K &k, const V &v) : key(k), value(v), left(nullptr), right(nullptr) {}
  } *root;

  int num_nodes;

  static void rotate_left(Node *&n) {
    Node *tmp = n;
    n = n->right;
    tmp->right = n->left;
    n->left = tmp;
  }

  static void rotate_right(Node *&n) {
    Node *tmp = n;
    n = n->left;
    tmp->left = n->right;
    n->right = tmp;
  }

  static void splay(Node *&n, const K &k) {
    if (n == nullptr) {
      return;
    }
    if (k < n->key && n->left != nullptr) {
      if (k < n->left->key) {
        splay(n->left->left, k);
        rotate_right(n);
      } else if (n->left->key < k) {
        splay(n->left->right, k);
        if (n->left->right != nullptr) {
          rotate_left(n->left);
        }
      }
      if (n->left != nullptr) {
        rotate_right(n);
      }
    } else if (n->key < k && n->right != nullptr) {
      if (k < n->right->key) {
        splay(n->right->left, k);
        if (n->right->left != nullptr) {
          rotate_right(n->right);
        }
      } else if (n->right->key < k) {
        splay(n->right->right, k);
        rotate_left(n);
      }
      if (n->right != nullptr) {
        rotate_left(n);
      }
    }
  }

  static bool insert(Node *&n, const K &k, const V &v) {
    if (n == nullptr) {
      n = new Node(k, v);
      return true;
    }
    splay(n, k);
    if (k < n->key) {
      Node *tmp = new Node(k, v);
      tmp->left = n->left;
      tmp->right = n;
      n->left = nullptr;
      n = tmp;
    } else if (n->key < k) {
      Node *tmp = new Node(k, v);
      tmp->left = n;
      tmp->right = n->right;
      n->right = nullptr;
      n = tmp;
    } else {
      return false;
    }
    return true;
  }

  static bool erase(Node *&n, const K &k) {
    if (n == nullptr) {
      return false;
    }
    splay(n, k);
    if (k < n->key || n->key < k) {
      return false;
    }
    Node *tmp = n;
    if (n->left == nullptr) {
      n = n->right;
    } else {
      splay(n->left, k);
      n = n->left;
      n->right = tmp->right;
    }
    delete tmp;
    return true;
  }

  template<class Fn>
  static void walk(Node *n, Fn f) {
    if (n != nullptr) {
      walk(n->left, f);
      f(n->key, n->value);
      walk(n->right, f);
    }
  }

  static void clean_up(Node *n) {
    if (n != nullptr) {
      clean_up(n->left);
      clean_up(n->right);
      delete n;
    }
  }

 public:
  SplayTree() : root(nullptr), num_nodes(0) {}

  ~SplayTree() { clean_up(root); }
  int size() const { return num_nodes; }
  bool empty() const { return root == nullptr; }

  bool insert(const K &k, const V &v) {
    if (insert(root, k, v)) {
      num_nodes++;
      return true;
    }
    return false;
  }

  bool erase(const K &k) {
    if (erase(root, k)) {
      num_nodes--;
      return true;
    }
    return false;
  }

  const V *find(const K &k) {
    splay(root, k);
    return (k < root->key || root->key < k) ? nullptr : &(root->value);
  }

  template<class Fn>
  void walk(Fn f) const {
    walk(root, f);
  }
};
\end{lstlisting}
\Needspace*{4\baselineskip}
\textbf{Example Usage}\par
\begin{lstlisting}[language=C++]
#include <cassert>
#include <iostream>
using namespace std;

int main() {
  SplayTree<int, char> t;
  t.insert(2, 'b');
  t.insert(1, 'a');
  t.insert(3, 'c');
  t.insert(5, 'e');
  assert(t.insert(4, 'd'));
  assert(*t.find(4) == 'd');
  assert(!t.insert(4, 'd'));
  auto printch = [](int k, char v) { cout << v; };
  t.walk(printch);
  cout << endl;
  assert(t.erase(1));
  assert(!t.erase(1));
  assert(t.find(1) == nullptr);
  t.walk(printch);
  cout << endl;
  return 0;
}
\end{lstlisting}
\begin{exampleoutput}
abcde
bcde
\end{exampleoutput}
\subsection{Size Balanced Tree}
Maintain a map, that is, a collection of key-value pairs such that each possible key appears at most once in the collection. In addition, support queries for keys given their ranks as well as queries for the ranks of given keys. This implementation requires an ordering on the set of possible keys defined by \inlinecode{operator\textless{}} on the key type. A size balanced tree augments each nodes with the size of its subtree, using it to maintain balance and compute order statistics.

\begin{compactitem}
  \item \inlinecode{SBTree()} constructs an empty map.
  \item \inlinecode{size()} returns the size of the map.
  \item \inlinecode{empty()} returns whether the map is empty.
  \item \inlinecode{insert(k, v)} adds an entry with key \inlinecode{k} and value \inlinecode{v} to the map, returning \inlinecode{true} if an new entry was added or \inlinecode{false} if the key already exists (in which case the map is unchanged and the old value associated with the key is preserved).
  \item \inlinecode{erase(k)} removes the entry with key \inlinecode{k} from the map, returning \inlinecode{true} if the removal was successful or \inlinecode{false} if the key to be removed was not found.
  \item \inlinecode{find(k)} returns a pointer to a const value associated with key \inlinecode{k}, or \inlinecode{nullptr} if the key was not found.
  \item \inlinecode{select(r)} returns a key-value pair of the node with a key of zero-based rank \inlinecode{r} in the map, throwing an exception if the rank is not between 0 and \inlinecode{size() - 1}.
  \item \inlinecode{rank(k)} returns the zero-based rank of key \inlinecode{k} in the map, throwing an exception if the key was not found in the map.
  \item \inlinecode{walk(f)} calls the function \inlinecode{f(k, v)} on each entry of the map, in ascending order of keys.
\end{compactitem}

\Needspace*{3\baselineskip}
\textbf{Time Complexity}:
\begin{compactitem}
  \item $O(1)$ per call to the constructor, \inlinecode{size()}, and \inlinecode{empty()}.
  \item $O(\log n)$ per call to \inlinecode{insert()}, \inlinecode{erase()}, \inlinecode{find()}, \inlinecode{select()}, and \inlinecode{rank()}, where $n$ is the number of entries currently in the map.
  \item $O(n)$ per call to \inlinecode{walk()}.
\end{compactitem}

\Needspace*{3\baselineskip}
\textbf{Space Complexity}:
\begin{compactitem}
  \item $O(n)$ for storage of the map elements.
  \item $O(\log n)$ auxiliary stack space for \inlinecode{insert()}, \inlinecode{erase()}, and \inlinecode{walk()}.
  \item $O(1)$ auxiliary for all other operations.
\end{compactitem}

\begin{lstlisting}[language=C++]
#include <cstddef>
#include <stdexcept>
#include <utility>

template<class K, class V>
class SBTree {
  struct Node {
    K key;
    V value;
    int size;
    Node *left, *right;

    Node(const K &k, const V &v) : key(k), value(v), size(1), left(nullptr), right(nullptr) {}

    inline Node *&child(int c) { return (c == 0) ? left : right; }

    void update() {
      size = 1;
      if (left != nullptr) {
        size += left->size;
      }
      if (right != nullptr) {
        size += right->size;
      }
    }
  } *root;

  static inline int size(Node *n) { return (n == nullptr) ? 0 : n->size; }

  static void rotate(Node *&n, int c) {
    Node *tmp = n->child(c);
    n->child(c) = tmp->child(!c);
    tmp->child(!c) = n;
    n->update();
    tmp->update();
    n = tmp;
  }

  static void maintain(Node *&n, int c) {
    if (n == nullptr || n->child(c) == nullptr) {
      return;
    }
    Node *&tmp = n->child(c);
    if (size(tmp->child(c)) > size(n->child(!c))) {
      rotate(n, c);
    } else if (size(tmp->child(!c)) > size(n->child(!c))) {
      rotate(tmp, !c);
      rotate(n, c);
    } else {
      return;
    }
    maintain(n->left, 0);
    maintain(n->right, 1);
    maintain(n, 0);
    maintain(n, 1);
  }

  static bool insert(Node *&n, const K &k, const V &v) {
    if (n == nullptr) {
      n = new Node(k, v);
      return true;
    }
    bool found;
    if (k < n->key) {
      found = insert(n->left, k, v);
      maintain(n, 0);
    } else if (n->key < k) {
      found = insert(n->right, k, v);
      maintain(n, 1);
    } else {
      return false;
    }
    n->update();
    return found;
  }

  static bool erase(Node *&n, const K &k) {
    if (n == nullptr) {
      return false;
    }
    bool found;
    int c = (k < n->key);
    if (k < n->key) {
      found = erase(n->left, k);
    } else if (n->key < k) {
      found = erase(n->right, k);
    } else {
      if (n->right == nullptr || n->left == nullptr) {
        Node *tmp = n;
        n = (n->right == nullptr) ? n->left : n->right;
        delete tmp;
        return true;
      }
      Node *p = n->right;
      while (p->left != nullptr) {
        p = p->left;
      }
      n->key = p->key;
      found = erase(n->right, p->key);
    }
    maintain(n, c);
    n->update();
    return found;
  }

  static std::pair<K, V> select(Node *n, int r) {
    int rank = size(n->left);
    if (r < rank) {
      return select(n->left, r);
    } else if (r > rank) {
      return select(n->right, r - rank - 1);
    }
    return {n->key, n->value};
  }

  static int rank(Node *n, const K &k) {
    if (n == nullptr) {
      throw std::runtime_error("Cannot rank key that's not in tree.");
    }
    int r = size(n->left);
    if (k < n->key) {
      return rank(n->left, k);
    } else if (n->key < k) {
      return rank(n->right, k) + r + 1;
    }
    return r;
  }

  template<class Fn>
  static void walk(Node *n, Fn f) {
    if (n != nullptr) {
      walk(n->left, f);
      f(n->key, n->value);
      walk(n->right, f);
    }
  }

  static void clean_up(Node *n) {
    if (n != nullptr) {
      clean_up(n->left);
      clean_up(n->right);
      delete n;
    }
  }

 public:
  SBTree() : root(nullptr) {}

  ~SBTree() { clean_up(root); }
  int size() const { return size(root); }
  bool empty() const { return root == nullptr; }
  bool insert(const K &k, const V &v) { return insert(root, k, v); }
  bool erase(const K &k) { return erase(root, k); }
  int rank(const K &k) const { return rank(root, k); }

  const V *find(const K &k) const {
    Node *n = root;
    while (n != nullptr) {
      if (k < n->key) {
        n = n->left;
      } else if (n->key < k) {
        n = n->right;
      } else {
        return &(n->value);
      }
    }
    return nullptr;
  }

  std::pair<K, V> select(int r) const {
    if (r < 0 || r >= size(root)) {
      throw std::runtime_error("Select rank must be between 0 and size() - 1.");
    }
    return select(root, r);
  }

  template<class Fn>
  void walk(Fn f) const {
    walk(root, f);
  }
};
\end{lstlisting}
\Needspace*{4\baselineskip}
\textbf{Example Usage}\par
\begin{lstlisting}[language=C++]
#include <cassert>
#include <iostream>
using namespace std;

int main() {
  SBTree<int, char> t;
  t.insert(2, 'b');
  t.insert(1, 'a');
  t.insert(3, 'c');
  t.insert(5, 'e');
  assert(t.insert(4, 'd'));
  assert(*t.find(4) == 'd');
  assert(!t.insert(4, 'd'));
  auto printch = [](int k, char v) { cout << v; };
  t.walk(printch);
  cout << endl;
  assert(t.erase(1));
  assert(!t.erase(1));
  assert(t.find(1) == nullptr);
  t.walk(printch);
  cout << endl;
  assert(t.rank(2) == 0);
  assert(t.rank(3) == 1);
  assert(t.rank(5) == 3);
  assert(t.select(0).first == 2);
  assert(t.select(1).first == 3);
  assert(t.select(2).first == 4);
  return 0;
}
\end{lstlisting}
\begin{exampleoutput}
abcde
bcde
\end{exampleoutput}
\subsection{Interval Treap}
Maintain a map from closed, one-dimensional intervals to values while supporting efficient reporting of any or all entries that intersect with a given query interval. This implementation uses \inlinecode{std::pair} to represent intervals, requiring operators \inlinecode{\textless{}} and \inlinecode{==} to be defined on the numeric key type. A treap is used to process the entries, where keys are compared lexicographically as pairs.

\begin{compactitem}
  \item \inlinecode{IntervalTreap()} constructs an empty map.
  \item \inlinecode{size()} returns the size of the map.
  \item \inlinecode{empty()} returns whether the map is empty.
  \item \inlinecode{insert(lo, hi, v)} adds an entry with key \inlinecode{[lo, hi]} and value \inlinecode{v} to the map, returning \inlinecode{true} if a new interval was added or \inlinecode{false} if the interval already exists (in which case the map is unchanged and the old value associated with the key is preserved).
  \item \inlinecode{erase(lo, hi)} removes the entry with key \inlinecode{[lo, hi]} from the map, returning \inlinecode{true} if the removal was successful or \inlinecode{false} if the interval was not found.
  \item \inlinecode{find\_key(lo, hi)} returns a pointer to a const \inlinecode{std::pair} representing the key of some interval in the map which intersects with \inlinecode{[lo, hi]}, or \inlinecode{nullptr} if no such entry was found.
  \item \inlinecode{find\_value(lo, hi)} returns a pointer to a const value of some entry in the map with a key that intersects with \inlinecode{[lo, hi]}, or \inlinecode{nullptr} if no such entry was found.
  \item \inlinecode{find\_all(lo, hi, f)} calls the function \inlinecode{f(lo, hi, v)} on each entry in the map that overlaps with \inlinecode{[lo, hi]}, in lexicographically ascending order of intervals.
  \item \inlinecode{walk(f)} calls the function \inlinecode{f(lo, hi, v)} on each interval in the map, in lexicographically ascending order of intervals.
\end{compactitem}

\Needspace*{3\baselineskip}
\textbf{Time Complexity}:
\begin{compactitem}
  \item $O(1)$ per call to the constructor, \inlinecode{size()}, and \inlinecode{empty()}.
  \item $O(\log n)$ on average per call to \inlinecode{insert()}, \inlinecode{erase()}, and \inlinecode{find\_any()}, where $n$ is the number of intervals currently in the set.
  \item $O(\log n + m)$ on average per call to \inlinecode{find\_all()}, where $m$ is the number of intersecting intervals that are reported.
  \item $O(n)$ per call to \inlinecode{walk()}.
\end{compactitem}

\Needspace*{3\baselineskip}
\textbf{Space Complexity}:
\begin{compactitem}
  \item $O(n)$ for storage of the map elements.
  \item $O(1)$ auxiliary for \inlinecode{size()} and \inlinecode{empty()}.
  \item $O(\log n)$ auxiliary stack space on average for all other operations.
\end{compactitem}

\begin{lstlisting}[language=C++]
#include <cstdlib>
#include <utility>

template<class K, class V>
class IntervalTreap {
  using Interval = std::pair<K, K>;

  struct Node {
    static inline int rand32() { return (rand() & 0x7fff) | ((rand() & 0x7fff) << 15); }

    Interval interval;
    V value;
    K max;
    int priority;
    Node *left, *right;

    Node(const Interval &i, const V &v)
        : interval(i), value(v), max(i.second), priority(rand32()), left(nullptr), right(nullptr) {}

    void update() {
      max = interval.second;
      if (left != nullptr && left->max > max) {
        max = left->max;
      }
      if (right != nullptr && right->max > max) {
        max = right->max;
      }
    }
  } *root;

  int num_nodes;

  static void rotate_left(Node *&n) {
    Node *tmp = n;
    n = n->right;
    tmp->right = n->left;
    n->left = tmp;
    tmp->update();
  }

  static void rotate_right(Node *&n) {
    Node *tmp = n;
    n = n->left;
    tmp->left = n->right;
    n->right = tmp;
    tmp->update();
  }

  static bool insert(Node *&n, const Interval &i, const V &v) {
    if (n == nullptr) {
      n = new Node(i, v);
      return true;
    }
    if (i < n->interval && insert(n->left, i, v)) {
      if (n->left->priority < n->priority) {
        rotate_right(n);
      }
      n->update();
      return true;
    }
    if (i > n->interval && insert(n->right, i, v)) {
      if (n->right->priority < n->priority) {
        rotate_left(n);
      }
      n->update();
      return true;
    }
    return false;
  }

  static bool erase(Node *&n, const Interval &i) {
    if (n == nullptr) {
      return false;
    }
    if (i < n->interval) {
      return erase(n->left, i);
    }
    if (i > n->interval) {
      return erase(n->right, i);
    }
    if (n->left != nullptr && n->right != nullptr) {
      bool res;
      if (n->left->priority < n->right->priority) {
        rotate_right(n);
        res = erase(n->right, i);
      } else {
        rotate_left(n);
        res = erase(n->left, i);
      }
      n->update();
      return res;
    }
    Node *tmp = (n->left != nullptr) ? n->left : n->right;
    delete n;
    n = tmp;
    return true;
  }

  static Node *find_any(Node *n, const Interval &i) {
    if (n == nullptr) {
      return nullptr;
    }
    if (n->interval.first <= i.second && i.first <= n->interval.second) {
      return n;
    }
    if (n->left != nullptr && i.first <= n->left->max) {
      return find_any(n->left, i);
    }
    return find_any(n->right, i);
  }

  template<class Fn>
  static void find_all(Node *n, const Interval &i, Fn f) {
    if (n == nullptr || n->max < i.first) {
      return;
    }
    if (n->interval.first <= i.second && i.first <= n->interval.second) {
      f(n->interval.first, n->interval.second, n->value);
    }
    find_all(n->left, i, f);
    find_all(n->right, i, f);
  }

  template<class Fn>
  static void walk(Node *n, Fn f) {
    if (n != nullptr) {
      walk(n->left, f);
      f(n->interval.first, n->interval.second, n->value);
      walk(n->right, f);
    }
  }

  static void clean_up(Node *n) {
    if (n != nullptr) {
      clean_up(n->left);
      clean_up(n->right);
      delete n;
    }
  }

 public:
  IntervalTreap() : root(nullptr), num_nodes(0) {}

  ~IntervalTreap() { clean_up(root); }
  int size() const { return num_nodes; }
  bool empty() const { return root == nullptr; }

  bool insert(const K &lo, const K &hi, const V &v) {
    if (insert(root, Interval{lo, hi}, v)) {
      num_nodes++;
      return true;
    }
    return false;
  }

  bool erase(const K &lo, const K &hi) {
    if (erase(root, Interval{lo, hi})) {
      num_nodes--;
      return true;
    }
    return false;
  }

  const Interval *find_key(const K &lo, const K &hi) const {
    Node *n = find_any(root, Interval{lo, hi});
    return (n == nullptr) ? nullptr : &(n->interval);
  }

  const V *find_value(const K &lo, const K &hi) const {
    Node *n = find_any(root, Interval{lo, hi});
    return (n == nullptr) ? nullptr : &(n->value);
  }

  template<class Fn>
  void find_all(const K &lo, const K &hi, Fn f) const {
    find_all(root, Interval{lo, hi}, f);
  }

  template<class Fn>
  void walk(Fn f) const {
    walk(root, f);
  }
};
\end{lstlisting}
\Needspace*{4\baselineskip}
\textbf{Example Usage}\par
\begin{lstlisting}[language=C++]
#include <cassert>
#include <iostream>
using namespace std;

int main() {
  IntervalTreap<int, char> t;
  t.insert(15, 20, 'a');
  t.insert(10, 30, 'b');
  t.insert(17, 19, 'c');
  t.insert(5, 20, 'd');
  t.insert(12, 15, 'e');
  t.insert(10, 40, 'f');
  assert(t.size() == 6);
  assert(!t.insert(5, 20, 'x'));
  t.erase(17, 19);
  assert(t.size() == 5);
  assert(*t.find_key(3, 9) == (pair<int, int>{5, 20}));
  assert(*t.find_value(3, 9) == 'd');
  cout << "Intervals intersecting [16, 20]:";
  auto print = [](int lo, int hi, char v) { cout << " [" << lo << ", " << hi << "]"; };
  t.find_all(16, 20, print);
  cout << "\nAll intervals:";
  t.walk(print);
  cout << endl;
  return 0;
}
\end{lstlisting}
\begin{exampleoutput}
Intervals intersecting [16, 20]: [15, 20] [10, 30] [5, 20] [10, 40]
All intervals: [5, 20] [10, 30] [10, 40] [12, 15] [15, 20]
\end{exampleoutput}
\subsection{Hash Map}
Maintain a map, that is, a collection of key-value pairs such that each possible key appears at most once in the collection. This implementation requires \inlinecode{operator==} to be defined on the key type. A hash map implements a map by hashing keys into buckets using a hash function. This implementation resolves collisions by chaining entries hashed to the same bucket into a linked list.

\begin{compactitem}
  \item \inlinecode{HashMap()} constructs an empty map.
  \item \inlinecode{size()} returns the size of the map.
  \item \inlinecode{empty()} returns whether the map is empty.
  \item \inlinecode{insert(k, v)} adds an entry with key \inlinecode{k} and value \inlinecode{v} to the map, returning \inlinecode{true} if an new entry was added or \inlinecode{false} if the key already exists (in which case the map is unchanged and the old value associated with the key is preserved).
  \item \inlinecode{erase(k)} removes the entry with key \inlinecode{k} from the map, returning \inlinecode{true} if the removal was successful or \inlinecode{false} if the key to be removed was not found.
  \item \inlinecode{find(k)} returns a pointer to the value associated with key \inlinecode{k}, or \inlinecode{nullptr} if the key was not found.
  \item \inlinecode{operator[]} returns a reference to key \inlinecode{k}'s associated value (which may be modified), or if necessary, inserts and returns a new entry with the default constructed value if key \inlinecode{k} was not originally found.
  \item \inlinecode{walk(f)} calls the function \inlinecode{f(k, v)} on each entry of the map, in no guaranteed order.
\end{compactitem}

\Needspace*{3\baselineskip}
\textbf{Time Complexity}:
\begin{compactitem}
  \item $O(1)$ per call to the constructor, \inlinecode{size()}, and \inlinecode{empty()}.
  \item $O(1)$ amortized per call to \inlinecode{insert()}, \inlinecode{erase()}, \inlinecode{find()}, and \inlinecode{operator[]}.
  \item $O(n)$ per call to \inlinecode{walk()}, where $n$ is the number of entries in the map.
\end{compactitem}

\Needspace*{3\baselineskip}
\textbf{Space Complexity}:
\begin{compactitem}
  \item $O(n)$ for storage of the map elements.
  \item $O(n)$ auxiliary heap space for \inlinecode{insert()}.
  \item $O(1)$ auxiliary for all other operations.
\end{compactitem}

\begin{lstlisting}[language=C++]
#include <cstddef>
#include <list>
#include <vector>

template<class K, class V, class Hash>
class HashMap {
  struct HashMapEntry {
    K key;
    V value;

    HashMapEntry(const K &k, const V &v) : key(k), value(v) {}
  };

  std::vector<std::list<HashMapEntry>> table;
  int table_size, num_entries;

  void double_capacity_and_rehash() {
    std::vector<std::list<HashMapEntry>> old = std::move(table);
    table_size = 2 * table_size;
    table.assign(table_size, std::list<HashMapEntry>());
    num_entries = 0;
    for (auto &bucket : old) {
      for (auto &entry : bucket) {
        insert(entry.key, entry.value);
      }
    }
  }

 public:
  explicit HashMap(int size = 128) : table(size), table_size(size), num_entries(0) {}

  int size() const { return num_entries; }
  bool empty() const { return num_entries == 0; }

  bool insert(const K &k, const V &v) {
    if (find(k) != nullptr) {
      return false;
    }
    if (num_entries >= table_size) {
      double_capacity_and_rehash();
    }
    unsigned int i = Hash()(k) % table_size;
    table[i].emplace_back(k, v);
    num_entries++;
    return true;
  }

  bool erase(const K &k) {
    unsigned int i = Hash()(k) % table_size;
    auto it = table[i].begin();
    while (it != table[i].end() && !(it->key == k)) {
      ++it;
    }
    if (it == table[i].end()) {
      return false;
    }
    table[i].erase(it);
    num_entries--;
    return true;
  }

  V *find(const K &k) {
    unsigned int i = Hash()(k) % table_size;
    auto it = table[i].begin();
    while (it != table[i].end() && !(it->key == k)) {
      ++it;
    }
    if (it == table[i].end()) {
      return nullptr;
    }
    return &(it->value);
  }

  V &operator[](const K &k) {
    if (V *ptr = find(k); ptr != nullptr) {
      return *ptr;
    }
    if (num_entries >= table_size) {
      double_capacity_and_rehash();
    }
    unsigned int i = Hash()(k) % table_size;
    table[i].emplace_back(k, V());
    num_entries++;
    return table[i].back().value;
  }

  template<class Fn>
  void walk(Fn f) const {
    for (int i = 0; i < table_size; i++) {
      for (const auto &entry : table[i]) {
        f(entry.key, entry.value);
      }
    }
  }
};
\end{lstlisting}
\Needspace*{4\baselineskip}
\textbf{Example Usage}\par
\begin{lstlisting}[language=C++]
#include <cassert>
#include <iostream>
using namespace std;

struct ClassHash {
  unsigned int operator()(int k) { return ClassHash()(static_cast<unsigned int>(k)); }
  unsigned int operator()(long long k) { return ClassHash()(static_cast<unsigned long long>(k)); }

  // Knuth's one-to-one multiplicative method.
  unsigned int operator()(unsigned int k) {
    return k * 2654435761u;  // Or just return k.
  }

  // Jenkins's 64-bit hash.
  unsigned int operator()(unsigned long long k) {
    k += ~(k << 32);
    k ^= (k >> 22);
    k += ~(k << 13);
    k ^= (k >> 8);
    k += (k << 3);
    k ^= (k >> 15);
    k += ~(k << 27);
    k ^= (k >> 31);
    return k;
  }

  // Jenkins's one-at-a-time hash.
  unsigned int operator()(const std::string &k) {
    unsigned int hash = 0;
    for (char c : k) {
      hash += ((hash + static_cast<unsigned char>(c)) << 10);
      hash ^= (hash >> 6);
    }
    hash += (hash << 3);
    hash ^= (hash >> 11);
    return hash + (hash << 15);
  }
};

int main() {
  HashMap<string, char, ClassHash> m;
  m["foo"] = 'a';
  m.insert("bar", 'b');
  assert(m["foo"] == 'a');
  assert(m["bar"] == 'b');
  assert(m["baz"] == '\0');
  m["baz"] = 'c';
  m.walk([](const string &k, char v) { cout << v; });
  cout << endl;
  assert(m.erase("foo"));
  assert(m.size() == 2);
  assert(m["foo"] == '\0');
  assert(m.size() == 3);
  return 0;
}
\end{lstlisting}
\begin{exampleoutput}
cab
\end{exampleoutput}
\subsection{Skip List}
Maintain a map, that is, a collection of key-value pairs such that each possible key appears at most once in the collection. This implementation requires operators \inlinecode{\textless{}} and \inlinecode{==} to be defined on the key type. A skip list maintains a linked hierarchy of sorted subsequences with each successive subsequence skipping over fewer elements than the previous one.

\begin{compactitem}
  \item \inlinecode{SkipList()} constructs an empty map.
  \item \inlinecode{size()} returns the size of the map.
  \item \inlinecode{empty()} returns whether the map is empty.
  \item \inlinecode{insert(k, v)} adds an entry with key \inlinecode{k} and value \inlinecode{v} to the map, returning \inlinecode{true} if an new entry was added or \inlinecode{false} if the key already exists (in which case the map is unchanged and the old value associated with the key is preserved).
  \item \inlinecode{erase(k)} removes the entry with key \inlinecode{k} from the map, returning \inlinecode{true} if the removal was successful or \inlinecode{false} if the key to be removed was not found.
  \item \inlinecode{find(k)} returns a pointer to the value associated with key \inlinecode{k}, or \inlinecode{nullptr} if the key was not found.
  \item \inlinecode{operator[]} returns a reference to key \inlinecode{k}'s associated value (which may be modified), or if necessary, inserts and returns a new entry with the default constructed value if key \inlinecode{k} was not originally found.
  \item \inlinecode{walk(f)} calls the function \inlinecode{f(k, v)} on each entry of the map, in ascending order of keys.
\end{compactitem}

\Needspace*{3\baselineskip}
\textbf{Time Complexity}:
\begin{compactitem}
  \item $O(1)$ per call to the constructor, \inlinecode{size()}, and \inlinecode{empty()}.
  \item $O(\log n)$ on average per call to \inlinecode{insert()}, \inlinecode{erase()}, \inlinecode{find()}, and \inlinecode{operator[]}, where $n$ is the number of entries currently in the map.
  \item $O(n)$ per call to \inlinecode{walk()}.
\end{compactitem}

\Needspace*{3\baselineskip}
\textbf{Space Complexity}:
\begin{compactitem}
  \item $O(n)$ on average for storage of the map elements.
  \item $O(n)$ auxiliary heap space for \inlinecode{insert()} and \inlinecode{erase()}.
  \item $O(1)$ auxiliary for all other operations.
\end{compactitem}

\begin{lstlisting}[language=C++]
#include <random>
#include <vector>

template<class K, class V>
class SkipList {
  static const int MAX_LEVELS = 32;  // log2(max possible keys)

  struct Node {
    K key;
    V value;
    std::vector<Node *> next;

    Node(const K &k, const V &v, int levels) : key(k), value(v), next(levels, (Node *)nullptr) {}
  } *head;

  int num_nodes;

  static int random_level() {
    static std::mt19937 rng(std::random_device{}());
    static std::uniform_int_distribution<int> coin(0, 1);
    int level = 1;
    while (coin(rng) && level < MAX_LEVELS) {
      level++;
    }
    return level;
  }

  static int node_level(const std::vector<Node *> &v) {
    int i = 0;
    while (i < static_cast<int>(v.size()) && v[i] != nullptr) {
      i++;
    }
    return i + 1;
  }

 public:
  SkipList() : head(new Node(K(), V(), MAX_LEVELS)), num_nodes(0) {
    for (auto &ptr : head->next) {
      ptr = nullptr;
    }
  }

  ~SkipList() { delete head; }
  int size() const { return num_nodes; }
  bool empty() const { return num_nodes == 0; }

  bool insert(const K &k, const V &v) {
    std::vector<Node *> update(head->next);
    int curr_level = node_level(update);
    Node *n = head;
    for (int i = curr_level; i-- > 0;) {
      while (n->next[i] != nullptr && n->next[i]->key < k) {
        n = n->next[i];
      }
      update[i] = n;
    }
    n = n->next[0];
    if (n != nullptr && n->key == k) {
      return false;
    }
    int new_level = random_level();
    if (new_level > curr_level) {
      for (int i = curr_level; i < new_level; i++) {
        update[i] = head;
      }
    }
    n = new Node(k, v, new_level);
    for (int i = 0; i < new_level; i++) {
      n->next[i] = update[i]->next[i];
      update[i]->next[i] = n;
    }
    num_nodes++;
    return true;
  }

  bool erase(const K &k) {
    std::vector<Node *> update(head->next);
    Node *n = head;
    for (int i = node_level(update); i-- > 0;) {
      while (n->next[i] != nullptr && n->next[i]->key < k) {
        n = n->next[i];
      }
      update[i] = n;
    }
    n = n->next[0];
    if (n != nullptr && n->key == k) {
      for (int i = 0, levels = static_cast<int>(update.size()); i < levels; i++) {
        if (update[i]->next[i] != n) {
          break;
        }
        update[i]->next[i] = n->next[i];
      }
      delete n;
      num_nodes--;
      return true;
    }
    return false;
  }

  V *find(const K &k) const {
    Node *n = head;
    for (int i = node_level(n->next); i-- > 0;) {
      while (n->next[i] != nullptr && n->next[i]->key < k) {
        n = n->next[i];
      }
    }
    n = n->next[0];
    return (n != nullptr && n->key == k) ? &(n->value) : nullptr;
  }

  V &operator[](const K &k) {
    V *ptr = find(k);
    if (ptr != nullptr) {
      return *ptr;
    }
    insert(k, V());
    return *find(k);
  }

  template<class Fn>
  void walk(Fn f) const {
    Node *n = head->next[0];
    while (n != nullptr) {
      f(n->key, n->value);
      n = n->next[0];
    }
  }
};
\end{lstlisting}
\Needspace*{4\baselineskip}
\textbf{Example Usage}\par
\begin{lstlisting}[language=C++]
#include <cassert>
#include <iostream>
using namespace std;

int main() {
  SkipList<int, char> l;
  l.insert(2, 'b');
  l.insert(1, 'a');
  l.insert(3, 'c');
  l.insert(5, 'e');
  assert(l.insert(4, 'd'));
  assert(*l.find(4) == 'd');
  assert(!l.insert(4, 'd'));
  auto printch = [](int k, char v) { cout << v; };
  l.walk(printch);
  cout << endl;
  assert(l.erase(1));
  assert(!l.erase(1));
  assert(l.find(1) == nullptr);
  l.walk(printch);
  cout << endl;
  return 0;
}
\end{lstlisting}
\begin{exampleoutput}
abcde
bcde
\end{exampleoutput}

\section{Range Queries in One Dimension}
\setcounter{section}{4}
\setcounter{subsection}{0}
\subsection{Sparse Table (Range Minimum Query)}
Given a static array with indices from $0$ to $n - 1$, precompute a table that may later be used to perform range minimum queries on the array in constant time. This version is simplified to only work on integer arrays.

The dynamic programming state $dp[i][j]$ holds the index of the minimum value in the sub-array starting at $i$ and having length $2^j$. Each $dp[i][j]$ will always be equal to either $dp[i][j - 1]$ or $dp[i + 2^{j - 1}][j - 1]$, whichever of the indices corresponds to the smaller value in the array.

\Needspace*{3\baselineskip}
\textbf{Time Complexity}:
\begin{compactitem}
  \item $O(n \log n)$ per call to \inlinecode{build()}, where $n$ is the size of the array.
  \item $O(1)$ per call to \inlinecode{query()}.
\end{compactitem}

\Needspace*{3\baselineskip}
\textbf{Space Complexity}:
\begin{compactitem}
  \item $O(n \log n)$ for storage of the sparse table, where $n$ is the size of the array.
  \item $O(1)$ auxiliary for \inlinecode{query()}.
\end{compactitem}

\begin{lstlisting}[language=C++]
#include <vector>

std::vector<int> table;
std::vector<std::vector<int>> dp;

void build(const std::vector<int> &a) {
  int n = a.size();
  table.resize(n + 1);
  dp.assign(n, std::vector<int>());
  for (int i = 2; i <= n; i++) {
    table[i] = table[i >> 1] + 1;
  }
  for (int i = 0; i < n; i++) {
    dp[i].resize(table[n] + 1);
    dp[i][0] = i;
  }
  for (int j = 1; (1 << j) < n; j++) {
    for (int i = 0; i + (1 << j) <= n; i++) {
      int x = dp[i][j - 1];
      int y = dp[i + (1 << (j - 1))][j - 1];
      dp[i][j] = (a[x] < a[y]) ? x : y;
    }
  }
}

int query(const std::vector<int> &a, int lo, int hi) {
  int j = table[hi - lo];
  int x = dp[lo][j];
  int y = dp[hi - (1 << j) + 1][j];
  return (a[x] < a[y]) ? a[x] : a[y];
}
\end{lstlisting}
\Needspace*{4\baselineskip}
\textbf{Example Usage}\par
\begin{lstlisting}[language=C++]
#include <cassert>

int main() {
  std::vector<int> a{6, -2, 1, 8, 10};
  build(a);
  assert(query(a, 0, 3) == -2);
  return 0;
}
\end{lstlisting}
\subsection{Square Root Decomposition}
Maintain a fixed-size array (from 0 to \inlinecode{size() - 1}) while supporting dynamic queries of contiguous sub-arrays and dynamic updates of individual indices.

The query operation is defined by an associative aggregate function \inlinecode{combine(a, b)}. The default definition below assumes a numerical array type, supporting queries for the ``min'' of the target range. Another possible query operation is ``sum'', in which case \inlinecode{combine(a, b)} should return $a + b$.

The point update operation is defined by \inlinecode{apply\_delta(v, d)}, which returns the new value at a single updated index. The default definition below supports updates that ``set'' the chosen array index to a new value. Another possible update operation is ``increment'', in which case \inlinecode{apply\_delta(v, d)} should return $v + d$.

The operations supported by this data structure are identical to those of the point update segment tree found in this section.

\Needspace*{3\baselineskip}
\textbf{Time Complexity}:
\begin{compactitem}
  \item $O(n)$ per call to both constructors, where $n$ is the size of the array.
  \item $O(1)$ per call to \inlinecode{size()}.
  \item $O(\sqrt{n})$ per call to \inlinecode{at()}, \inlinecode{update()}, and \inlinecode{query()}.
\end{compactitem}

\Needspace*{3\baselineskip}
\textbf{Space Complexity}:
\begin{compactitem}
  \item $O(n)$ for storage of the array elements.
  \item $O(1)$ auxiliary for all operations.
\end{compactitem}

\begin{lstlisting}[language=C++]
#include <algorithm>
#include <cmath>
#include <vector>

template<class T>
class SqrtDecomposition {
  static T combine(const T &a, const T &b) { return std::min(a, b); }
  static T apply_delta(const T &v, const T &d) { return d; }

  int len, blocklen;
  std::vector<T> value, block;

  void init() {
    blocklen = std::max(1, static_cast<int>(sqrt(len)));
    int nblocks = (len + blocklen - 1) / blocklen;
    for (int i = 0; i < nblocks; i++) {
      T blockval = value[i * blocklen];
      int blockhi = std::min(len, (i + 1) * blocklen);
      for (int j = i * blocklen + 1; j < blockhi; j++) {
        blockval = combine(blockval, value[j]);
      }
      block.push_back(blockval);
    }
  }

 public:
  explicit SqrtDecomposition(int n, const T &v = T()) : len(n), value(n, v) { init(); }

  template<class It>
  SqrtDecomposition(It lo, It hi) : len(hi - lo), value(lo, hi) {
    init();
  }

  int size() const { return len; }
  T at(int i) const { return query(i, i); }

  T query(int lo, int hi) const {
    T res;
    int blocklo = ceil(static_cast<double>(lo) / blocklen), blockhi = (hi + 1) / blocklen - 1;
    if (blocklo > blockhi) {
      res = value[lo];
      for (int i = lo + 1; i <= hi; i++) {
        res = combine(res, value[i]);
      }
    } else {
      res = block[blocklo];
      for (int i = blocklo + 1; i <= blockhi; i++) {
        res = combine(res, block[i]);
      }
      for (int i = lo; i < blocklo * blocklen; i++) {
        res = combine(res, value[i]);
      }
      for (int i = (blockhi + 1) * blocklen; i <= hi; i++) {
        res = combine(res, value[i]);
      }
    }
    return res;
  }

  void update(int i, const T &d) {
    value[i] = apply_delta(value[i], d);
    int b = i / blocklen;
    int blockhi = std::min(len, (b + 1) * blocklen);
    block[b] = value[b * blocklen];
    for (int i = b * blocklen + 1; i < blockhi; i++) {
      block[b] = combine(block[b], value[i]);
    }
  }
};
\end{lstlisting}
\Needspace*{4\baselineskip}
\textbf{Example Usage}\par
\begin{lstlisting}[language=C++]
#include <cassert>
#include <iostream>
using namespace std;

int main() {
  vector<int> arr{6, -2, 1, 8, 10};
  SqrtDecomposition<int> sd(arr.begin(), arr.end());
  sd.update(2, 4);
  cout << "Values:";
  for (int i = 0; i < sd.size(); i++) {
    cout << " " << sd.at(i);
  }
  cout << endl;
  assert(sd.query(0, 3) == -2);
  return 0;
}
\end{lstlisting}
\begin{exampleoutput}
Values: 6 -2 4 8 10
\end{exampleoutput}
\subsection{Segment Tree (Point Update)}
Maintain a fixed-size array while supporting dynamic queries of contiguous sub-arrays and dynamic updates of individual indices.

The query operation is defined by an associative aggregate function \inlinecode{combine(a, b)}. The default code below assumes a numerical array type, defining queries for the ``min'' of the target range. Another possible query operation is ``sum'', in which case \inlinecode{combine(a, b)} should return $a + b$.

The point update operation is defined by \inlinecode{apply\_delta(v, d)}, which returns the new value at a single updated index. The default definition below supports updates that ``set'' the chosen array index to a new value. Another possible update operation is ``increment'', in which case \inlinecode{apply\_delta(v, d)} should return $v + d$.

\begin{compactitem}
  \item \inlinecode{SegTree(n, v)} constructs an array of size \inlinecode{n} with indices from 0 to \inlinecode{n - 1}, inclusive, and all values initialized to \inlinecode{v}.
  \item \inlinecode{SegTree(lo, hi)} constructs an array from two random-access iterators as a range \inlinecode{[lo, hi)}, initialized to the elements of the range in the same order.
  \item \inlinecode{size()} returns the size of the array.
  \item \inlinecode{at(i)} returns the value at index \inlinecode{i}.
  \item \inlinecode{query(lo, hi)} returns the result of \inlinecode{combine()} applied to all indices from \inlinecode{lo} to \inlinecode{hi}, inclusive. If \inlinecode{lo == hi}, then the single specified value is returned.
  \item \inlinecode{update(i, d)} assigns the value \inlinecode{v} at index \inlinecode{i} to \inlinecode{apply\_delta(v, d)}.
\end{compactitem}

\Needspace*{3\baselineskip}
\textbf{Time Complexity}:
\begin{compactitem}
  \item $O(n)$ per call to both constructors, where $n$ is the size of the array.
  \item $O(1)$ per call to \inlinecode{size()}.
  \item $O(\log n)$ per call to \inlinecode{at()}, \inlinecode{update()}, and \inlinecode{query()}.
\end{compactitem}

\Needspace*{3\baselineskip}
\textbf{Space Complexity}:
\begin{compactitem}
  \item $O(n)$ for storage of the array elements.
  \item $O(\log n)$ auxiliary stack space for \inlinecode{update()} and \inlinecode{query()}.
  \item $O(1)$ auxiliary for \inlinecode{size()}.
\end{compactitem}

\begin{lstlisting}[language=C++]
#include <algorithm>
#include <vector>

template<class T>
class SegTree {
  static T combine(const T &a, const T &b) { return std::min(a, b); }
  static T apply_delta(const T &v, const T &d) { return d; }

  int len;
  std::vector<T> value;

  void build(int i, int lo, int hi, const T &v) {
    if (lo == hi) {
      value[i] = v;
      return;
    }
    int mid = lo + (hi - lo) / 2;
    build(i * 2 + 1, lo, mid, v);
    build(i * 2 + 2, mid + 1, hi, v);
    value[i] = combine(value[i * 2 + 1], value[i * 2 + 2]);
  }

  template<class It>
  void build(int i, int lo, int hi, It arr) {
    if (lo == hi) {
      value[i] = *(arr + lo);
      return;
    }
    int mid = lo + (hi - lo) / 2;
    build(i * 2 + 1, lo, mid, arr);
    build(i * 2 + 2, mid + 1, hi, arr);
    value[i] = combine(value[i * 2 + 1], value[i * 2 + 2]);
  }

  T query(int i, int lo, int hi, int tgt_lo, int tgt_hi) const {
    if (lo == tgt_lo && hi == tgt_hi) {
      return value[i];
    }
    int mid = lo + (hi - lo) / 2;
    if (tgt_lo <= mid && mid < tgt_hi) {
      return combine(
          query(i * 2 + 1, lo, mid, tgt_lo, std::min(tgt_hi, mid)),
          query(i * 2 + 2, mid + 1, hi, std::max(tgt_lo, mid + 1), tgt_hi)
      );
    }
    if (tgt_lo <= mid) {
      return query(i * 2 + 1, lo, mid, tgt_lo, std::min(tgt_hi, mid));
    }
    return query(i * 2 + 2, mid + 1, hi, std::max(tgt_lo, mid + 1), tgt_hi);
  }

  void update(int i, int lo, int hi, int target, const T &d) {
    if (target < lo || target > hi) {
      return;
    }
    if (lo == hi) {
      value[i] = apply_delta(value[i], d);
      return;
    }
    int mid = lo + (hi - lo) / 2;
    update(i * 2 + 1, lo, mid, target, d);
    update(i * 2 + 2, mid + 1, hi, target, d);
    value[i] = combine(value[i * 2 + 1], value[i * 2 + 2]);
  }

 public:
  explicit SegTree(int n, const T &v = T()) : len(n), value(4 * len) {
    if (len > 0) {
      build(0, 0, len - 1, v);
    }
  }

  template<class It>
  SegTree(It lo, It hi) : len(hi - lo), value(4 * len) {
    if (len > 0) {
      build(0, 0, len - 1, lo);
    }
  }

  int size() const { return len; }
  T at(int i) const { return query(i, i); }
  T query(int lo, int hi) const { return query(0, 0, len - 1, lo, hi); }
  void update(int i, const T &d) { update(0, 0, len - 1, i, d); }
};
\end{lstlisting}
\Needspace*{4\baselineskip}
\textbf{Example Usage}\par
\begin{lstlisting}[language=C++]
#include <cassert>
#include <iostream>
using namespace std;

int main() {
  vector<int> arr{6, -2, 1, 8, 10};
  SegTree<int> t(arr.begin(), arr.end());
  t.update(2, 4);
  cout << "Values:";
  for (int i = 0; i < t.size(); i++) {
    cout << " " << t.at(i);
  }
  cout << endl;
  assert(t.query(0, 3) == -2);
  return 0;
}
\end{lstlisting}
\begin{exampleoutput}
Values: 6 -2 4 8 10
The minimum value in the range [0, 3] is -2.
\end{exampleoutput}
\subsection{Segment Tree (Range Update)}
Maintain a fixed-size array while supporting both dynamic queries and updates of contiguous subarrays via the lazy propagation technique.

The query operation is defined by an associative aggregate function \inlinecode{combine(a, b)}. The default code below assumes a numerical array type, defining queries for the ``min'' of the target range. Another possible query operation is ``sum'', in which case \inlinecode{combine(a, b)} should return $a + b$.

Range updates are defined by \inlinecode{apply\_delta(v, d, len)}, which applies an update delta \inlinecode{d} to an aggregate summary \inlinecode{v} representing \inlinecode{len} array values, and by \inlinecode{compose\_deltas(old, d)}, which combines a pending older delta with a newer delta in that order. These functions do not support arbitrary combinations: applying a delta to a combined segment must be equivalent to applying it to each child segment and then combining the results, and composed deltas must be equivalent to performing their updates sequentially. The default code below defines range assignment. For range increment, \inlinecode{compose\_deltas(old, d)} should return \inlinecode{old + d}; \inlinecode{apply\_delta(v, d, len)} should return \inlinecode{v + d} for range-min/range-max queries, and \inlinecode{v + d * len} for range-sum queries.

\begin{compactitem}
  \item \inlinecode{SegTree(n, v)} constructs an array of size \inlinecode{n} with indices from 0 to \inlinecode{n - 1}, inclusive, and all values initialized to \inlinecode{v}.
  \item \inlinecode{SegTree(lo, hi)} constructs an array from two random-access iterators as a range \inlinecode{[lo, hi)}, initialized to the elements of the range in the same order.
  \item \inlinecode{size()} returns the size of the array.
  \item \inlinecode{at(i)} returns the value at index \inlinecode{i}, where \inlinecode{i} is between 0 and \inlinecode{size() - 1}.
  \item \inlinecode{query(lo, hi)} returns the result of \inlinecode{combine()} applied to all indices from \inlinecode{lo} to \inlinecode{hi}, inclusive. If \inlinecode{lo == hi}, then the single specified value is returned.
  \item \inlinecode{update(i, d)} assigns the value \inlinecode{v} at index \inlinecode{i} to \inlinecode{apply\_delta(v, d)}.
  \item \inlinecode{update(lo, hi, d)} modifies the value at each array index from \inlinecode{lo} to \inlinecode{hi}, inclusive, by applying the delta \inlinecode{d} to each value.
\end{compactitem}

\Needspace*{3\baselineskip}
\textbf{Time Complexity}:
\begin{compactitem}
  \item $O(n)$ per call to both constructors, where $n$ is the size of the array.
  \item $O(1)$ per call to \inlinecode{size()}.
  \item $O(\log n)$ per call to \inlinecode{at()}, \inlinecode{update()}, and \inlinecode{query()}.
\end{compactitem}

\Needspace*{3\baselineskip}
\textbf{Space Complexity}:
\begin{compactitem}
  \item $O(n)$ for storage of the array elements.
  \item $O(\log n)$ auxiliary stack space for \inlinecode{update()} and \inlinecode{query()}.
  \item $O(1)$ auxiliary for \inlinecode{size()}.
\end{compactitem}

\begin{lstlisting}[language=C++]
#include <algorithm>
#include <vector>

template<class T>
class SegTree {
  static T combine(const T &a, const T &b) { return std::min(a, b); }
  static T apply_delta(const T &v, const T &d, long long len) { return d; }
  static T compose_deltas(const T &d1, const T &d2) { return d2; }

  int len;
  std::vector<T> value, delta;
  std::vector<bool> pending;

  void build(int i, int lo, int hi, const T &v) {
    if (lo == hi) {
      value[i] = v;
      return;
    }
    int mid = lo + (hi - lo) / 2;
    build(i * 2 + 1, lo, mid, v);
    build(i * 2 + 2, mid + 1, hi, v);
    value[i] = combine(value[i * 2 + 1], value[i * 2 + 2]);
  }

  template<class It>
  void build(int i, int lo, int hi, It arr) {
    if (lo == hi) {
      value[i] = *(arr + lo);
      return;
    }
    int mid = lo + (hi - lo) / 2;
    build(i * 2 + 1, lo, mid, arr);
    build(i * 2 + 2, mid + 1, hi, arr);
    value[i] = combine(value[i * 2 + 1], value[i * 2 + 2]);
  }

  void push_delta(int i, int lo, int hi) {
    if (pending[i]) {
      value[i] = apply_delta(value[i], delta[i], hi - lo + 1);
      if (lo != hi) {
        int l = 2 * i + 1, r = 2 * i + 2;
        delta[l] = pending[l] ? compose_deltas(delta[l], delta[i]) : delta[i];
        delta[r] = pending[r] ? compose_deltas(delta[r], delta[i]) : delta[i];
        pending[l] = pending[r] = true;
      }
      pending[i] = false;
    }
  }

  T query(int i, int lo, int hi, int tgt_lo, int tgt_hi) {
    push_delta(i, lo, hi);
    if (lo == tgt_lo && hi == tgt_hi) {
      return value[i];
    }
    int mid = lo + (hi - lo) / 2;
    if (tgt_lo <= mid && mid < tgt_hi) {
      return combine(
          query(i * 2 + 1, lo, mid, tgt_lo, std::min(tgt_hi, mid)),
          query(i * 2 + 2, mid + 1, hi, std::max(tgt_lo, mid + 1), tgt_hi)
      );
    }
    if (tgt_lo <= mid) {
      return query(i * 2 + 1, lo, mid, tgt_lo, std::min(tgt_hi, mid));
    }
    return query(i * 2 + 2, mid + 1, hi, std::max(tgt_lo, mid + 1), tgt_hi);
  }

  void update(int i, int lo, int hi, int tgt_lo, int tgt_hi, const T &d) {
    push_delta(i, lo, hi);
    if (hi < tgt_lo || lo > tgt_hi) {
      return;
    }
    if (tgt_lo <= lo && hi <= tgt_hi) {
      delta[i] = d;
      pending[i] = true;
      push_delta(i, lo, hi);
      return;
    }
    update(2 * i + 1, lo, lo + (hi - lo) / 2, tgt_lo, tgt_hi, d);
    update(2 * i + 2, lo + (hi - lo) / 2 + 1, hi, tgt_lo, tgt_hi, d);
    value[i] = combine(value[2 * i + 1], value[2 * i + 2]);
  }

 public:
  explicit SegTree(int n, const T &v = T())
      : len(n), value(4 * len), delta(4 * len), pending(4 * len, false) {
    if (len > 0) {
      build(0, 0, len - 1, v);
    }
  }

  template<class It>
  SegTree(It lo, It hi) : len(hi - lo), value(4 * len), delta(4 * len), pending(4 * len, false) {
    if (len > 0) {
      build(0, 0, len - 1, lo);
    }
  }

  int size() const { return len; }
  T at(int i) { return query(i, i); }
  T query(int lo, int hi) { return query(0, 0, len - 1, lo, hi); }
  void update(int i, const T &d) { update(0, 0, len - 1, i, i, d); }
  void update(int lo, int hi, const T &d) { update(0, 0, len - 1, lo, hi, d); }
};
\end{lstlisting}
\Needspace*{4\baselineskip}
\textbf{Example Usage}\par
\begin{lstlisting}[language=C++]
#include <cassert>
#include <iostream>
using namespace std;

int main() {
  vector<int> arr{6, -2, 1, 8, 10};
  SegTree<int> t(arr.begin(), arr.end());
  t.update(2, 4);
  cout << "Values:";
  for (int i = 0; i < t.size(); i++) {
    cout << " " << t.at(i);
  }
  cout << endl;
  assert(t.query(0, 3) == -2);
  t.update(0, 4, 5);
  t.update(3, 2);
  t.update(3, 1);
  cout << "Values:";
  for (int i = 0; i < t.size(); i++) {
    cout << " " << t.at(i);
  }
  cout << endl;
  assert(t.query(0, 3) == 1);
  return 0;
}
\end{lstlisting}
\begin{exampleoutput}
Values: 6 -2 4 8 10
Values: 5 5 5 1 5
\end{exampleoutput}
\subsection{Segment Tree (Compressed)}
Maintain an array over a large index range while supporting both dynamic queries and updates of contiguous subarrays via the lazy propagation technique. This implementation uses lazy initialization of nodes to conserve memory: only the nodes covering touched indices are ever allocated, so indices up to \inlinecode{N} are supported without preallocating the whole tree.

The query operation is defined by an associative aggregate function \inlinecode{combine(a, b)}. Since untouched nodes are implicit, \inlinecode{repeat\_value(v, len)} must return the aggregate summary of \inlinecode{len} copies of the initial value \inlinecode{v}. The default code below assumes a numerical array type, defining queries for the ``min'' of the target range. For range-sum queries, \inlinecode{combine(a, b)} should return $a + b$ and \inlinecode{repeat\_value(v, len)} should return $v \cdot len$.

Range updates are defined by \inlinecode{apply\_delta(v, d, len)}, which applies an update delta \inlinecode{d} to an aggregate summary \inlinecode{v} representing \inlinecode{len} array values, and by \inlinecode{compose\_deltas(old, d)}, which combines a pending older delta with a newer delta in that order. These functions do not support arbitrary combinations: applying a delta to a combined segment must be equivalent to applying it to each child segment and then combining the results, and composed deltas must be equivalent to performing their updates sequentially. The default code below defines range assignment. For range increment, \inlinecode{compose\_deltas(old, d)} should return \inlinecode{old + d}; \inlinecode{apply\_delta(v, d, len)} should return \inlinecode{v + d} for range-min/range-max queries, and \inlinecode{v + d * len} for range-sum queries.

\begin{compactitem}
  \item \inlinecode{SegTree(v)} constructs an array over indices 0 to \inlinecode{N}, inclusive, with every value implicitly initialized to \inlinecode{v}. Nodes are allocated lazily as indices are touched.
  \item \inlinecode{at(i)} returns the value at index \inlinecode{i}, where \inlinecode{i} is between 0 and \inlinecode{N}.
  \item \inlinecode{query(lo, hi)} returns the result of \inlinecode{combine()} applied to all indices from \inlinecode{lo} to \inlinecode{hi}, inclusive. If \inlinecode{lo == hi}, then the single specified value is returned.
  \item \inlinecode{update(i, d)} assigns the value \inlinecode{v} at index \inlinecode{i} to \inlinecode{apply\_delta(v, d)}.
  \item \inlinecode{update(lo, hi, d)} modifies the value at each array index from \inlinecode{lo} to \inlinecode{hi}, inclusive, by applying the delta \inlinecode{d} to each value.
\end{compactitem}

\Needspace*{3\baselineskip}
\textbf{Time Complexity}:
\begin{compactitem}
  \item $O(1)$ per call to the constructor.
  \item $O(\log N)$ per call to \inlinecode{at()}, \inlinecode{update()}, and \inlinecode{query()}.
\end{compactitem}

\Needspace*{3\baselineskip}
\textbf{Space Complexity}:
\begin{compactitem}
  \item $O(k \log\left(N\right))$ for storage after $k$ index updates.
  \item $O(\log N)$ auxiliary stack space for \inlinecode{update()} and \inlinecode{query()}.
\end{compactitem}

\begin{lstlisting}[language=C++]
#include <algorithm>
#include <cstddef>

template<class T>
class SegTree {
  static const int N = 1000000000;

  static T combine(const T &a, const T &b) { return std::min(a, b); }
  static T repeat_value(const T &v, long long len) { return v; }
  static T apply_delta(const T &v, const T &d, long long len) { return d; }
  static T compose_deltas(const T &d1, const T &d2) { return d2; }

  struct Node {
    T value, delta;
    bool pending;
    Node *left, *right;

    explicit Node(const T &v) : value(v), pending(false), left(nullptr), right(nullptr) {}
  } *root;

  T init;

  void update_delta(Node *&n, const T &d, long long len) {
    if (n == nullptr) {
      n = new Node(repeat_value(init, len));
    }
    n->delta = n->pending ? compose_deltas(n->delta, d) : d;
    n->pending = true;
  }

  void push_delta(Node *n, int lo, int hi) {
    if (n == nullptr) {
      return;
    }
    if (n->pending) {
      n->value = apply_delta(n->value, n->delta, hi - lo + 1);
      if (lo != hi) {
        int mid = lo + (hi - lo) / 2;
        update_delta(n->left, n->delta, mid - lo + 1);
        update_delta(n->right, n->delta, hi - mid);
      }
    }
    n->pending = false;
  }

  T query(Node *n, int lo, int hi, int tgt_lo, int tgt_hi) {
    if (n == nullptr) {
      return repeat_value(init, hi - lo + 1);
    }
    push_delta(n, lo, hi);
    if (lo == tgt_lo && hi == tgt_hi) {
      return n->value;
    }
    int mid = lo + (hi - lo) / 2;
    if (tgt_lo <= mid && mid < tgt_hi) {
      return combine(
          query(n->left, lo, mid, tgt_lo, std::min(tgt_hi, mid)),
          query(n->right, mid + 1, hi, std::max(tgt_lo, mid + 1), tgt_hi)
      );
    }
    if (tgt_lo <= mid) {
      return query(n->left, lo, mid, tgt_lo, std::min(tgt_hi, mid));
    }
    return query(n->right, mid + 1, hi, std::max(tgt_lo, mid + 1), tgt_hi);
  }

  void update(Node *&n, int lo, int hi, int tgt_lo, int tgt_hi, const T &d) {
    if (n == nullptr) {
      if (hi < tgt_lo || lo > tgt_hi) {
        return;
      }
      n = new Node(repeat_value(init, hi - lo + 1));
    } else {
      push_delta(n, lo, hi);
    }
    if (hi < tgt_lo || lo > tgt_hi) {
      return;
    }
    if (tgt_lo <= lo && hi <= tgt_hi) {
      n->delta = d;
      n->pending = true;
      push_delta(n, lo, hi);
      return;
    }
    int mid = lo + (hi - lo) / 2;
    update(n->left, lo, mid, tgt_lo, tgt_hi, d);
    update(n->right, mid + 1, hi, tgt_lo, tgt_hi, d);
    T left_value = (n->left != nullptr) ? n->left->value : repeat_value(init, mid - lo + 1);
    T right_value = (n->right != nullptr) ? n->right->value : repeat_value(init, hi - mid);
    n->value = combine(left_value, right_value);
  }

  void clean_up(Node *n) {
    if (n != nullptr) {
      clean_up(n->left);
      clean_up(n->right);
      delete n;
    }
  }

 public:
  explicit SegTree(const T &v = T()) : root(nullptr), init(v) {}

  ~SegTree() { clean_up(root); }
  T at(int i) { return query(i, i); }
  T query(int lo, int hi) { return query(root, 0, N, lo, hi); }
  void update(int i, const T &d) { return update(i, i, d); }
  void update(int lo, int hi, const T &d) { return update(root, 0, N, lo, hi, d); }
};
\end{lstlisting}
\Needspace*{4\baselineskip}
\textbf{Example Usage}\par
\begin{lstlisting}[language=C++]
#include <cassert>
#include <iostream>
using namespace std;

int main() {
  SegTree<int> t(0);
  t.update(0, 6);
  t.update(1, -2);
  t.update(2, 4);
  t.update(3, 8);
  t.update(4, 10);
  cout << "Values:";
  for (int i = 0; i < 5; i++) {
    cout << " " << t.at(i);
  }
  cout << endl;
  assert(t.query(0, 3) == -2);
  t.update(0, 4, 5);
  t.update(3, 2);
  t.update(3, 1);
  cout << "Values:";
  for (int i = 0; i < 5; i++) {
    cout << " " << t.at(i);
  }
  cout << endl;
  assert(t.query(0, 3) == 1);
  return 0;
}
\end{lstlisting}
\begin{exampleoutput}
Values: 6 -2 4 8 10
Values: 5 5 5 1 5
\end{exampleoutput}
\subsection{Segment Tree (Persistent)}
Maintains immutable versions of a segment tree under point updates. Each update creates a new root by copying only the nodes on the path from the root to the updated leaf, while all unchanged subtrees are shared with previous versions.

The implementation below stores sums over a fixed index range \inlinecode{[0, n - 1]}. Persistent segment trees are useful for offline order-statistics queries, rollback-like histories, and functional dynamic programming states.

\begin{compactitem}
  \item \inlinecode{PersistentSegTree(a)} constructs version 0 from array \inlinecode{a}.
  \item \inlinecode{update(version, index, value)} returns a new version where \inlinecode{a[index]} is set to \inlinecode{value}, leaving the old version unchanged.
  \item \inlinecode{query(version, lo, hi)} returns the sum of the inclusive range \inlinecode{[lo, hi]} in the specified version.
  \item \inlinecode{versions()} returns the number of stored roots.
\end{compactitem}

\Needspace*{3\baselineskip}
\textbf{Time Complexity}:
\begin{compactitem}
  \item $O(n)$ for construction.
  \item $O(\log n)$ per call to \inlinecode{update(version, index, value)} and \inlinecode{query(version, lo, hi)}.
\end{compactitem}

\Needspace*{3\baselineskip}
\textbf{Space Complexity}:
\begin{compactitem}
  \item $O(n + q \log n)$ for $q$ updates.
  \item $O(\log n)$ auxiliary stack space per update and query.
\end{compactitem}

\begin{lstlisting}[language=C++]
#include <vector>

class PersistentSegTree {
  struct node {
    long long sum;
    int left, right;

    node(long long sum = 0, int left = -1, int right = -1) : sum(sum), left(left), right(right) {}
  };

  int n;
  std::vector<node> tree;
  std::vector<int> root;

  int build(const std::vector<int> &a, int lo, int hi) {
    if (lo == hi) {
      tree.emplace_back(a[lo]);
      return static_cast<int>(tree.size()) - 1;
    }
    int mid = lo + (hi - lo) / 2;
    int left = build(a, lo, mid);
    int right = build(a, mid + 1, hi);
    tree.emplace_back(tree[left].sum + tree[right].sum, left, right);
    return static_cast<int>(tree.size()) - 1;
  }

  int update(int cur, int lo, int hi, int index, int value) {
    if (lo == hi) {
      tree.emplace_back(value);
      return static_cast<int>(tree.size()) - 1;
    }
    int mid = lo + (hi - lo) / 2;
    int left = tree[cur].left, right = tree[cur].right;
    if (index <= mid) {
      left = update(left, lo, mid, index, value);
    } else {
      right = update(right, mid + 1, hi, index, value);
    }
    tree.emplace_back(tree[left].sum + tree[right].sum, left, right);
    return static_cast<int>(tree.size()) - 1;
  }

  long long query(int cur, int lo, int hi, int qlo, int qhi) const {
    if (qlo <= lo && hi <= qhi) {
      return tree[cur].sum;
    }
    int mid = lo + (hi - lo) / 2;
    long long res = 0;
    if (qlo <= mid) {
      res += query(tree[cur].left, lo, mid, qlo, qhi);
    }
    if (mid < qhi) {
      res += query(tree[cur].right, mid + 1, hi, qlo, qhi);
    }
    return res;
  }

 public:
  explicit PersistentSegTree(const std::vector<int> &a) : n(a.size()) {
    if (n > 0) {
      root.push_back(build(a, 0, n - 1));
    }
  }

  int versions() const { return root.size(); }

  int update(int version, int index, int value) {
    root.push_back(update(root[version], 0, n - 1, index, value));
    return static_cast<int>(root.size()) - 1;
  }

  long long query(int version, int lo, int hi) const {
    return query(root[version], 0, n - 1, lo, hi);
  }
};
\end{lstlisting}
\Needspace*{4\baselineskip}
\textbf{Example Usage}\par
\begin{lstlisting}[language=C++]
#include <cassert>
using namespace std;

int main() {
  vector<int> a{1, 2, 3, 4};
  PersistentSegTree t(a);
  int v1 = t.update(0, 1, 10);
  int v2 = t.update(v1, 3, -1);
  assert(t.query(0, 0, 3) == 10);
  assert(t.query(v1, 0, 3) == 18);
  assert(t.query(v2, 1, 3) == 12);
  assert(t.versions() == 3);
  return 0;
}
\end{lstlisting}
\subsection{Cartesian Tree}
Builds the min Cartesian tree of an array in linear time. A Cartesian tree is a binary tree whose inorder traversal is the original array order and whose heap property is determined by array values. For a min Cartesian tree, each parent has value less than or equal to its children.

Cartesian trees connect arrays, stacks, range minimum queries, treaps, and largest rectangle style problems. Equal values are broken by position, so earlier equal values become ancestors of later equal values.

\begin{compactitem}
  \item \inlinecode{CartesianTree(a)} constructs the tree for array \inlinecode{a}.
  \item \inlinecode{root} stores the root index.
  \item \inlinecode{parent[i]}, \inlinecode{left[i]}, and \inlinecode{right[i]} store neighboring node indices, or $-1$ if absent.
\end{compactitem}

\textbf{Time Complexity}: $O(n)$ construction time, where $n$ is the array size.

\textbf{Space Complexity}: $O(n)$ for tree arrays and the monotone stack.

\begin{lstlisting}[language=C++]
#include <vector>

struct CartesianTree {
  int root;
  std::vector<int> parent, left, right;

  explicit CartesianTree(const std::vector<int> &a)
      : root(-1), parent(a.size(), -1), left(a.size(), -1), right(a.size(), -1) {
    std::vector<int> st;
    for (int i = 0, n = static_cast<int>(a.size()); i < n; i++) {
      int last = -1;
      while (!st.empty() && a[i] < a[st.back()]) {
        last = st.back();
        st.pop_back();
      }
      if (!st.empty()) {
        right[st.back()] = i;
        parent[i] = st.back();
      }
      if (last != -1) {
        left[i] = last;
        parent[last] = i;
      }
      st.push_back(i);
    }
    root = st.empty() ? -1 : st[0];
  }
};
\end{lstlisting}
\Needspace*{4\baselineskip}
\textbf{Example Usage}\par
\begin{lstlisting}[language=C++]
#include <cassert>
using namespace std;

void inorder(const CartesianTree &t, int u, vector<int> *order) {
  if (u == -1) {
    return;
  }
  inorder(t, t.left[u], order);
  order->push_back(u);
  inorder(t, t.right[u], order);
}

int main() {
  vector<int> a{3, 2, 6, 1, 9};
  CartesianTree t(a);
  assert(t.root == 3);
  assert(t.parent[1] == 3);
  assert(t.parent[4] == 3);

  vector<int> order;
  inorder(t, t.root, &order);
  for (int i = 0; i < 5; i++) {
    assert(order[i] == i);
  }
  return 0;
}
\end{lstlisting}
\subsection{Implicit Treap}
Maintain a dynamically-sized array using a balanced binary search tree while supporting both dynamic queries and updates of contiguous subarrays via the lazy propagation technique. A treap maintains a balanced binary tree structure by preserving the heap property on the randomly generated priority values of nodes, thereby making insertions and deletions run in $O(\log n)$ with high probability.

The query operation is defined by an associative aggregate function \inlinecode{combine(a, b)}. The default code below assumes a numerical array type, defining queries for the ``min'' of the target range. Another possible query operation is ``sum'', in which case \inlinecode{combine(a, b)} should return $a + b$.

Range updates are defined by \inlinecode{apply\_delta(v, d, len)}, which applies an update delta \inlinecode{d} to an aggregate summary \inlinecode{v} representing \inlinecode{len} array values, and by \inlinecode{compose\_deltas(old, d)}, which combines a pending older delta with a newer delta in that order. These functions do not support arbitrary combinations: applying a delta to a combined segment must be equivalent to applying it to each child segment and then combining the results, and composed deltas must be equivalent to performing their updates sequentially. The default code below defines range assignment. For range increment, \inlinecode{compose\_deltas(old, d)} should return \inlinecode{old + d}; \inlinecode{apply\_delta(v, d, len)} should return \inlinecode{v + d} for range-min/range-max queries, and \inlinecode{v + d * len} for range-sum queries.

This data structure shares every operation of one-dimensional segment trees in this section, with the additional operations \inlinecode{empty()}, \inlinecode{insert()}, \inlinecode{erase()}, \inlinecode{push\_back()}, and \inlinecode{pop\_back()} analogous to those of \inlinecode{std::vector} (here, \inlinecode{insert()} and \inlinecode{erase()} both take an index instead of an iterator).

\Needspace*{3\baselineskip}
\textbf{Time Complexity}:
\begin{compactitem}
  \item $O(n)$ per call to both constructors, where $n$ is the size of the array.
  \item $O(1)$ per call to \inlinecode{size()} and \inlinecode{empty()}.
  \item $O(\log n)$ on average per call to all other operations.
\end{compactitem}

\Needspace*{3\baselineskip}
\textbf{Space Complexity}:
\begin{compactitem}
  \item $O(n)$ for storage of the array elements.
  \item $O(1)$ auxiliary for \inlinecode{size()} and \inlinecode{empty()}.
  \item $O(\log n)$ auxiliary stack space for all other operations.
\end{compactitem}

\begin{lstlisting}[language=C++]
#include <cstdlib>

template<class T>
class ImplicitTreap {
  static T combine(const T &a, const T &b) { return a < b ? a : b; }
  static T apply_delta(const T &v, const T &d, long long len) { return d; }
  static T compose_deltas(const T &d1, const T &d2) { return d2; }

  struct Node {
    static inline int rand32() { return (rand() & 0x7fff) | ((rand() & 0x7fff) << 15); }

    T value, subtree_value, delta;
    bool pending;
    int size, priority;
    Node *left, *right;

    explicit Node(const T &v)
        : value(v),
          subtree_value(v),
          pending(false),
          size(1),
          priority(rand32()),
          left(nullptr),
          right(nullptr) {}
  } *root;

  static int size(Node *n) { return (n == nullptr) ? 0 : n->size; }

  static void update_value(Node *n) {
    if (n == nullptr) {
      return;
    }
    n->subtree_value = n->value;
    if (n->left != nullptr) {
      n->subtree_value = combine(n->subtree_value, n->left->subtree_value);
    }
    if (n->right != nullptr) {
      n->subtree_value = combine(n->subtree_value, n->right->subtree_value);
    }
    n->size = 1 + size(n->left) + size(n->right);
  }

  static void update_delta(Node *n, const T &d) {
    if (n != nullptr) {
      n->value = apply_delta(n->value, d, 1);
      n->subtree_value = apply_delta(n->subtree_value, d, n->size);
      n->delta = n->pending ? compose_deltas(n->delta, d) : d;
      n->pending = true;
    }
  }

  static void push_delta(Node *n) {
    if (n == nullptr || !n->pending) {
      return;
    }
    if (n->size > 1) {
      update_delta(n->left, n->delta);
      update_delta(n->right, n->delta);
    }
    n->pending = false;
  }

  static void merge(Node *&n, Node *left, Node *right) {
    push_delta(left);
    push_delta(right);
    if (left == nullptr) {
      n = right;
    } else if (right == nullptr) {
      n = left;
    } else if (left->priority < right->priority) {
      merge(left->right, left->right, right);
      n = left;
    } else {
      merge(right->left, left, right->left);
      n = right;
    }
    update_value(n);
  }

  static void split(Node *n, Node *&left, Node *&right, int i) {
    push_delta(n);
    if (n == nullptr) {
      left = right = nullptr;
    } else if (i <= size(n->left)) {
      split(n->left, left, n->left, i);
      right = n;
    } else {
      split(n->right, n->right, right, i - size(n->left) - 1);
      left = n;
    }
    update_value(n);
  }

  static void insert(Node *&n, Node *new_node, int i) {
    push_delta(n);
    if (n == nullptr) {
      n = new_node;
    } else if (new_node->priority < n->priority) {
      split(n, new_node->left, new_node->right, i);
      n = new_node;
    } else if (i <= size(n->left)) {
      insert(n->left, new_node, i);
    } else {
      insert(n->right, new_node, i - size(n->left) - 1);
    }
    update_value(n);
  }

  static void erase(Node *&n, int i) {
    push_delta(n);
    if (i == size(n->left)) {
      Node *left = n->left, *right = n->right;
      delete n;
      merge(n, left, right);
    } else if (i < size(n->left)) {
      erase(n->left, i);
    } else {
      erase(n->right, i - size(n->left) - 1);
    }
    update_value(n);
  }

  static Node *select(Node *n, int i) {
    push_delta(n);
    if (i < size(n->left)) {
      return select(n->left, i);
    }
    if (i > size(n->left)) {
      return select(n->right, i - size(n->left) - 1);
    }
    return n;
  }

  void clean_up(Node *&n) {
    if (n != nullptr) {
      clean_up(n->left);
      clean_up(n->right);
      delete n;
    }
  }

 public:
  explicit ImplicitTreap(int n = 0, const T &v = T()) : root(nullptr) {
    for (int i = 0; i < n; i++) {
      push_back(v);
    }
  }

  template<class It>
  ImplicitTreap(It lo, It hi) : root(nullptr) {
    for (; lo != hi; ++lo) {
      push_back(*lo);
    }
  }

  ~ImplicitTreap() { clean_up(root); }
  int size() const { return size(root); }
  bool empty() const { return root == nullptr; }
  void insert(int i, const T &v) { insert(root, new Node(v), i); }
  void erase(int i) { erase(root, i); }
  void push_back(const T &v) { insert(size(), v); }
  void pop_back() { erase(size() - 1); }
  T at(int i) const { return select(root, i)->value; }

  T query(int lo, int hi) {
    Node *l1, *r1, *l2, *r2, *t;
    split(root, l1, r1, hi + 1);
    split(l1, l2, r2, lo);
    T res = r2->subtree_value;
    merge(t, l2, r2);
    merge(root, t, r1);
    return res;
  }

  void update(int lo, int hi, const T &d) {
    Node *l1, *r1, *l2, *r2, *t;
    split(root, l1, r1, hi + 1);
    split(l1, l2, r2, lo);
    update_delta(r2, d);
    merge(t, l2, r2);
    merge(root, t, r1);
  }

  void update(int i, const T &d) { update(i, i, d); }
};
\end{lstlisting}
\Needspace*{4\baselineskip}
\textbf{Example Usage}\par
\begin{lstlisting}[language=C++]
#include <iostream>
#include <vector>
using namespace std;

void print(ImplicitTreap<int> &t) {
  cout << "Values:";
  for (int i = 0; i < t.size(); i++) {
    cout << " " << t.at(i);
  }
  cout << " (min: " << t.query(0, t.size() - 1) << ")" << endl;
}

int main() {
  vector<int> arr{99, -2, 1, 8, 10};
  ImplicitTreap<int> t(arr.begin(), arr.end());
  t.push_back(11);
  t.push_back(12);
  t.pop_back();
  print(t);
  t.insert(0, 90);
  t.erase(1);
  print(t);
  t.update(0, 1, 2);
  print(t);
  return 0;
}
\end{lstlisting}
\begin{exampleoutput}
Values: 99 -2 1 8 10 11 (min: -2)
Values: 90 -2 1 8 10 11 (min: -2)
Values: 2 2 1 8 10 11 (min: 1)
\end{exampleoutput}

\section{Range Queries in Two Dimensions}
\setcounter{section}{5}
\setcounter{subsection}{0}
\subsection{Quadtree (Point Update)}
Maintain a two-dimensional array while supporting dynamic queries of rectangular sub-arrays and dynamic updates of individual indices. This implementation uses lazy initialization of nodes to conserve memory while supporting large indices.

The query operation is defined by a commutative associative aggregate function \inlinecode{combine(a, b)}. Because untouched regions are implicit, \inlinecode{repeat\_value(v, area)} must return the aggregate summary of \inlinecode{area} copies of the initial value \inlinecode{v}. The default code below assumes a numerical array type, defining queries for the ``min'' of the target range. For rectangle-sum queries, \inlinecode{combine(a, b)} should return $a + b$ and \inlinecode{repeat\_value(v, area)} should return $v \cdot area$.

The point update operation is defined by \inlinecode{apply\_delta(v, d)}, which returns the new value at one updated cell. The default code below defines updates that ``set'' the chosen cell to a new value. Another possible update operation is ``increment'', in which case \inlinecode{apply\_delta(v, d)} should return $v + d$.

\begin{compactitem}
  \item \inlinecode{Quadtree(v)} constructs a two-dimensional array with rows from 0 to \inlinecode{R} and columns from 0 to \inlinecode{C}, inclusive. All values are implicitly initialized to \inlinecode{v}.
  \item \inlinecode{at(r, c)} returns the value at row \inlinecode{r}, column \inlinecode{c}.
  \item \inlinecode{query(r1, c1, r2, c2)} returns the result of \inlinecode{combine()} applied to every value in the rectangular region consisting of rows from \inlinecode{r1} to \inlinecode{r2}, inclusive, and columns from \inlinecode{c1} to \inlinecode{c2}, inclusive.
  \item \inlinecode{update(r, c, d)} assigns the value \inlinecode{v} at (\inlinecode{r}, \inlinecode{c}) to \inlinecode{apply\_delta(v, d)}.
\end{compactitem}

\Needspace*{3\baselineskip}
\textbf{Time Complexity}:
\begin{compactitem}
  \item $O(1)$ per call to the constructor.
  \item $O(\max\left(R, C\right))$ per call to \inlinecode{at()}, \inlinecode{update()}, and \inlinecode{query()}.
\end{compactitem}

\Needspace*{3\baselineskip}
\textbf{Space Complexity}:
\begin{compactitem}
  \item $O(n \log\left(\max\left(R, C\right)\right))$ for storage of the quadtree nodes after $n$ point updates.
  \item $O(\log\left(\max\left(R, C\right)\right))$ auxiliary stack space for \inlinecode{update()}, \inlinecode{query()}, and \inlinecode{at()}.
\end{compactitem}

\begin{lstlisting}[language=C++]
#include <algorithm>
#include <cstddef>

template<class T>
class Quadtree {
  static const int R = 1000000000;
  static const int C = 1000000000;

  static T combine(const T &a, const T &b) { return std::min(a, b); }
  static T repeat_value(const T &v, long long area) { return v; }
  static T apply_delta(const T &v, const T &d) { return d; }

  struct Node {
    T value;
    Node *child[4];

    explicit Node(const T &v) : value(v) {
      for (auto &c : child) {
        c = nullptr;
      }
    }
  };

  Node *root;
  T init;

  // Helper variables for query().
  int tgt_r1, tgt_c1, tgt_r2, tgt_c2;
  T res;
  bool found;

  static long long area(int r1, int c1, int r2, int c2) {
    return static_cast<long long>(r2 - r1 + 1) * (c2 - c1 + 1);
  }

  void append_result(const T &v) {
    res = found ? combine(res, v) : v;
    found = true;
  }

  T child_value(Node *n, int r1, int c1, int r2, int c2) {
    return (n != nullptr) ? n->value : repeat_value(init, area(r1, c1, r2, c2));
  }

  void query(Node *n, int r1, int c1, int r2, int c2) {
    if (tgt_r2 < r1 || tgt_r1 > r2 || tgt_c2 < c1 || tgt_c1 > c2) {
      return;
    }
    if (n == nullptr) {
      int rlen = std::min(r2, tgt_r2) - std::max(r1, tgt_r1) + 1;
      int clen = std::min(c2, tgt_c2) - std::max(c1, tgt_c1) + 1;
      append_result(repeat_value(init, static_cast<long long>(rlen) * clen));
      return;
    }
    if (tgt_r1 <= r1 && r2 <= tgt_r2 && tgt_c1 <= c1 && c2 <= tgt_c2) {
      append_result(n->value);
      return;
    }
    int rmid = r1 + (r2 - r1) / 2, cmid = c1 + (c2 - c1) / 2;
    query(n->child[0], r1, c1, rmid, cmid);
    query(n->child[1], rmid + 1, c1, r2, cmid);
    query(n->child[2], r1, cmid + 1, rmid, c2);
    query(n->child[3], rmid + 1, cmid + 1, r2, c2);
  }

  // Helper variables for update().
  int tgt_r, tgt_c;
  T delta;

  void update(Node *&n, int r1, int c1, int r2, int c2) {
    if (tgt_r < r1 || tgt_r > r2 || tgt_c < c1 || tgt_c > c2) {
      return;
    }
    if (n == nullptr) {
      n = new Node(repeat_value(init, area(r1, c1, r2, c2)));
    }
    if (r1 == r2 && c1 == c2) {
      n->value = apply_delta(n->value, delta);
      return;
    }
    int rmid = r1 + (r2 - r1) / 2, cmid = c1 + (c2 - c1) / 2;
    update(n->child[0], r1, c1, rmid, cmid);
    update(n->child[1], rmid + 1, c1, r2, cmid);
    update(n->child[2], r1, cmid + 1, rmid, c2);
    update(n->child[3], rmid + 1, cmid + 1, r2, c2);
    n->value = combine(
        combine(
            child_value(n->child[0], r1, c1, rmid, cmid),
            child_value(n->child[1], rmid + 1, c1, r2, cmid)
        ),
        combine(
            child_value(n->child[2], r1, cmid + 1, rmid, c2),
            child_value(n->child[3], rmid + 1, cmid + 1, r2, c2)
        )
    );
  }

  static void clean_up(Node *n) {
    if (n != nullptr) {
      for (auto *c : n->child) {
        clean_up(c);
      }
      delete n;
    }
  }

 public:
  explicit Quadtree(const T &v = T()) : root(nullptr), init(v) {}

  ~Quadtree() { clean_up(root); }
  T at(int r, int c) { return query(r, c, r, c); }

  T query(int r1, int c1, int r2, int c2) {
    tgt_r1 = r1;
    tgt_c1 = c1;
    tgt_r2 = r2;
    tgt_c2 = c2;
    found = false;
    query(root, 0, 0, R, C);
    return found ? res : repeat_value(init, area(r1, c1, r2, c2));
  }

  void update(int r, int c, const T &d) {
    tgt_r = r;
    tgt_c = c;
    delta = d;
    update(root, 0, 0, R, C);
  }
};
\end{lstlisting}
\Needspace*{4\baselineskip}
\textbf{Example Usage}\par
\begin{lstlisting}[language=C++]
#include <cassert>
#include <iostream>
using namespace std;

int main() {
  Quadtree<int> t(0);
  t.update(0, 0, 7);
  t.update(0, 1, 6);
  t.update(1, 0, 5);
  t.update(1, 1, 4);
  t.update(2, 1, 1);
  t.update(2, 2, 9);
  cout << "Values:" << endl;
  for (int i = 0; i < 3; i++) {
    for (int j = 0; j < 3; j++) {
      cout << t.at(i, j) << " ";
    }
    cout << endl;
  }
  assert(t.query(0, 0, 0, 1) == 6);
  assert(t.query(0, 0, 1, 0) == 5);
  assert(t.query(1, 1, 2, 2) == 0);
  assert(t.query(0, 0, 1000000000, 1000000000) == 0);
  t.update(500000000, 500000000, -100);
  assert(t.query(0, 0, 1000000000, 1000000000) == -100);
  return 0;
}
\end{lstlisting}
\begin{exampleoutput}
Values:
7 6 0
5 4 0
0 1 9
\end{exampleoutput}
\subsection{Quadtree (Range Update)}
Maintain a two-dimensional array while supporting both dynamic queries and updates of rectangular sub-arrays via the lazy propagation technique. This implementation uses lazy initialization of nodes to conserve memory while supporting large indices.

The query operation is defined by a commutative associative aggregate function \inlinecode{combine(a, b)}. Because untouched regions are implicit, \inlinecode{repeat\_value(v, area)} must return the aggregate summary of \inlinecode{area} copies of the initial value \inlinecode{v}. The default code below assumes a numerical array type, defining queries for the ``min'' of the target range. For rectangle-sum queries, \inlinecode{combine(a, b)} should return $a + b$ and \inlinecode{repeat\_value(v, area)} should return $v \cdot area$.

Range updates are defined by \inlinecode{apply\_delta(v, d, area)}, which applies an update delta \inlinecode{d} to an aggregate summary \inlinecode{v} representing \inlinecode{area} cells, and by \inlinecode{compose\_deltas(old, d)}, which combines a pending older delta with a newer delta in that order. These functions do not support arbitrary combinations: applying a delta to a combined region must be equivalent to applying it to each child region and then combining the results, and composed deltas must be equivalent to performing their updates sequentially. The default code below defines rectangle assignment. For rectangle increment, \inlinecode{compose\_deltas(old, d)} should return \inlinecode{old + d}; \inlinecode{apply\_delta(v, d, area)} should return \inlinecode{v + d} for range-min/range-max queries, and \inlinecode{v + d * area} for range-sum queries.

\begin{compactitem}
  \item \inlinecode{Quadtree(v)} constructs a two-dimensional array with rows from 0 to \inlinecode{R} and columns from 0 to \inlinecode{C}, inclusive. All values are implicitly initialized to \inlinecode{v}.
  \item \inlinecode{at(r, c)} returns the value at row \inlinecode{r}, column \inlinecode{c}.
  \item \inlinecode{query(r1, c1, r2, c2)} returns the result of \inlinecode{combine()} applied to every value in the rectangular region consisting of rows from \inlinecode{r1} to \inlinecode{r2} and columns from \inlinecode{c1} to \inlinecode{c2}, inclusive.
  \item \inlinecode{update(r, c, d)} assigns the value \inlinecode{v} at (\inlinecode{r}, \inlinecode{c}) to \inlinecode{apply\_delta(v, d)}.
  \item \inlinecode{update(r1, c1, r2, c2)} modifies the value at each index of the rectangular region consisting of rows from \inlinecode{r1} to \inlinecode{r2} and columns from \inlinecode{c1} to \inlinecode{c2}, inclusive, by applying the delta \inlinecode{d} to each value.
\end{compactitem}

\Needspace*{3\baselineskip}
\textbf{Time Complexity}:
\begin{compactitem}
  \item $O(1)$ per call to the constructor.
  \item $O(\max\left(R, C\right))$ per call to \inlinecode{at()}, \inlinecode{update()}, and \inlinecode{query()}.
\end{compactitem}

\Needspace*{3\baselineskip}
\textbf{Space Complexity}:
\begin{compactitem}
  \item $O(n \log\left(\max\left(R, C\right)\right))$ for storage of the quadtree nodes after $n$ point or rectangle updates.
  \item $O(\log\left(\max\left(R, C\right)\right))$ auxiliary stack space for \inlinecode{update()}, \inlinecode{query()}, and \inlinecode{at()}.
\end{compactitem}

\begin{lstlisting}[language=C++]
#include <algorithm>
#include <cstddef>

template<class T>
class Quadtree {
  static const int R = 1000000000;
  static const int C = 1000000000;

  static T combine(const T &a, const T &b) { return std::min(a, b); }
  static T repeat_value(const T &v, long long area) { return v; }
  static T apply_delta(const T &v, const T &d, long long area) { return d; }
  static T compose_deltas(const T &d1, const T &d2) { return d2; }

  struct Node {
    T value, delta;
    bool pending;
    Node *child[4];

    explicit Node(const T &v) : value(v), pending(false) {
      for (auto &c : child) {
        c = nullptr;
      }
    }
  } *root;

  T init;

  static long long area(int r1, int c1, int r2, int c2) {
    return static_cast<long long>(r2 - r1 + 1) * (c2 - c1 + 1);
  }

  void update_delta(Node *&n, const T &d, long long area) {
    if (n == nullptr) {
      n = new Node(repeat_value(init, area));
    }
    n->delta = n->pending ? compose_deltas(n->delta, d) : d;
    n->pending = true;
  }

  void push_delta(Node *n, int r1, int c1, int r2, int c2) {
    if (n->pending) {
      int rmid = r1 + (r2 - r1) / 2, cmid = c1 + (c2 - c1) / 2;
      long long cur_area = area(r1, c1, r2, c2);
      n->value = apply_delta(n->value, n->delta, cur_area);
      if (cur_area > 1) {
        int rlen = r2 - r1 + 1, clen = c2 - c1 + 1;
        int rlen1 = rmid - r1 + 1, rlen2 = rlen - rlen1;
        int clen1 = cmid - c1 + 1, clen2 = clen - clen1;
        update_delta(n->child[0], n->delta, static_cast<long long>(rlen1) * clen1);
        update_delta(n->child[1], n->delta, static_cast<long long>(rlen2) * clen1);
        update_delta(n->child[2], n->delta, static_cast<long long>(rlen1) * clen2);
        update_delta(n->child[3], n->delta, static_cast<long long>(rlen2) * clen2);
      }
      n->pending = false;
    }
  }

  // Helper variables for query() and update().
  int tgt_r1, tgt_c1, tgt_r2, tgt_c2;
  T res, delta;
  bool found;

  void append_result(const T &v) {
    res = found ? combine(res, v) : v;
    found = true;
  }

  T child_value(Node *n, int r1, int c1, int r2, int c2) {
    return (n != nullptr) ? n->value : repeat_value(init, area(r1, c1, r2, c2));
  }

  void query(Node *n, int r1, int c1, int r2, int c2) {
    if (tgt_r2 < r1 || tgt_r1 > r2 || tgt_c2 < c1 || tgt_c1 > c2) {
      return;
    }
    if (n == nullptr) {
      int rlen = std::min(r2, tgt_r2) - std::max(r1, tgt_r1) + 1;
      int clen = std::min(c2, tgt_c2) - std::max(c1, tgt_c1) + 1;
      append_result(repeat_value(init, static_cast<long long>(rlen) * clen));
      return;
    }
    push_delta(n, r1, c1, r2, c2);
    if (tgt_r1 <= r1 && r2 <= tgt_r2 && tgt_c1 <= c1 && c2 <= tgt_c2) {
      append_result(n->value);
      return;
    }
    int rmid = r1 + (r2 - r1) / 2, cmid = c1 + (c2 - c1) / 2;
    query(n->child[0], r1, c1, rmid, cmid);
    query(n->child[1], rmid + 1, c1, r2, cmid);
    query(n->child[2], r1, cmid + 1, rmid, c2);
    query(n->child[3], rmid + 1, cmid + 1, r2, c2);
  }

  void update(Node *&n, int r1, int c1, int r2, int c2) {
    if (n == nullptr) {
      if (tgt_r2 < r1 || tgt_r1 > r2 || tgt_c2 < c1 || tgt_c1 > c2) {
        return;
      }
      n = new Node(repeat_value(init, area(r1, c1, r2, c2)));
    } else {
      push_delta(n, r1, c1, r2, c2);
    }
    if (tgt_r2 < r1 || tgt_r1 > r2 || tgt_c2 < c1 || tgt_c1 > c2) {
      return;
    }
    if (tgt_r1 <= r1 && r2 <= tgt_r2 && tgt_c1 <= c1 && c2 <= tgt_c2) {
      n->delta = delta;
      n->pending = true;
      push_delta(n, r1, c1, r2, c2);
      return;
    }
    int rmid = r1 + (r2 - r1) / 2, cmid = c1 + (c2 - c1) / 2;
    update(n->child[0], r1, c1, rmid, cmid);
    update(n->child[1], rmid + 1, c1, r2, cmid);
    update(n->child[2], r1, cmid + 1, rmid, c2);
    update(n->child[3], rmid + 1, cmid + 1, r2, c2);
    n->value = combine(
        combine(
            child_value(n->child[0], r1, c1, rmid, cmid),
            child_value(n->child[1], rmid + 1, c1, r2, cmid)
        ),
        combine(
            child_value(n->child[2], r1, cmid + 1, rmid, c2),
            child_value(n->child[3], rmid + 1, cmid + 1, r2, c2)
        )
    );
  }

  static void clean_up(Node *n) {
    if (n != nullptr) {
      for (auto *c : n->child) {
        clean_up(c);
      }
      delete n;
    }
  }

 public:
  explicit Quadtree(const T &v = T()) : root(nullptr), init(v) {}

  ~Quadtree() { clean_up(root); }
  T at(int r, int c) { return query(r, c, r, c); }

  T query(int r1, int c1, int r2, int c2) {
    tgt_r1 = r1;
    tgt_c1 = c1;
    tgt_r2 = r2;
    tgt_c2 = c2;
    found = false;
    query(root, 0, 0, R, C);
    return found ? res : repeat_value(init, area(r1, c1, r2, c2));
  }

  void update(int r, int c, const T &d) { update(r, c, r, c, d); }

  void update(int r1, int c1, int r2, int c2, const T &d) {
    tgt_r1 = r1;
    tgt_c1 = c1;
    tgt_r2 = r2;
    tgt_c2 = c2;
    delta = d;
    update(root, 0, 0, R, C);
  }
};
\end{lstlisting}
\Needspace*{4\baselineskip}
\textbf{Example Usage}\par
\begin{lstlisting}[language=C++]
#include <cassert>
#include <iostream>
using namespace std;

int main() {
  Quadtree<int> t(0);
  t.update(0, 0, 7);
  t.update(0, 1, 6);
  t.update(1, 0, 5);
  t.update(1, 1, 4);
  t.update(2, 1, 1);
  t.update(2, 2, 9);
  cout << "Values:" << endl;
  for (int i = 0; i < 3; i++) {
    for (int j = 0; j < 3; j++) {
      cout << t.at(i, j) << " ";
    }
    cout << endl;
  }
  assert(t.query(0, 0, 0, 1) == 6);
  assert(t.query(0, 0, 1, 0) == 5);
  assert(t.query(1, 1, 2, 2) == 0);
  assert(t.query(0, 0, 1000000000, 1000000000) == 0);
  t.update(500000000, 500000000, -100);
  assert(t.query(0, 0, 1000000000, 1000000000) == -100);
  return 0;
}
\end{lstlisting}
\begin{exampleoutput}
Values:
7 6 0
5 4 0
0 1 9
\end{exampleoutput}
\subsection{2D Segment Tree}
Maintain a two-dimensional array while supporting dynamic queries of rectangular sub-arrays and dynamic updates of individual indices. This implementation uses lazy initialization of nodes to conserve memory while supporting large indices.

The query operation is defined by a commutative associative aggregate function \inlinecode{combine(a, b)}. Because untouched regions are implicit, \inlinecode{repeat\_value(v, area)} must return the aggregate summary of \inlinecode{area} copies of the initial value \inlinecode{v}. The default code below assumes a numerical array type, defining queries for the ``min'' of the target range. For rectangle-sum queries, \inlinecode{combine(a, b)} should return $a + b$ and \inlinecode{repeat\_value(v, area)} should return $v \cdot area$.

The point update operation is defined by \inlinecode{apply\_delta(v, d)}, which returns the new value at one updated cell. The default code below defines updates that ``set'' the chosen cell to a new value. Another possible update operation is ``increment'', in which case \inlinecode{apply\_delta(v, d)} should return $v + d$.

\begin{compactitem}
  \item \inlinecode{SegTree2D(v)} constructs a two-dimensional array with rows from 0 to \inlinecode{R} and columns from 0 to \inlinecode{C}, inclusive. All values are implicitly initialized to \inlinecode{v}.
  \item \inlinecode{at(r, c)} returns the value at row \inlinecode{r}, column \inlinecode{c}.
  \item \inlinecode{query(r1, c1, r2, c2)} returns the result of \inlinecode{combine()} applied to every value in the rectangular region consisting of rows from \inlinecode{r1} to \inlinecode{r2}, inclusive, and columns from \inlinecode{c1} to \inlinecode{c2}, inclusive.
  \item \inlinecode{update(r, c, d)} assigns the value \inlinecode{v} at (\inlinecode{r}, \inlinecode{c}) to \inlinecode{apply\_delta(v, d)}.
\end{compactitem}

\Needspace*{3\baselineskip}
\textbf{Time Complexity}:
\begin{compactitem}
  \item $O(1)$ per call to the constructor.
  \item $O(\log\left(R\right) \cdot \log\left(C\right))$ per call to \inlinecode{at()}, \inlinecode{update()}, and \inlinecode{query()}.
\end{compactitem}

\Needspace*{3\baselineskip}
\textbf{Space Complexity}:
\begin{compactitem}
  \item $O(n \log\left(R\right) \log\left(C\right))$ for storage after $n$ point updates.
  \item $O(\log\left(R\right) + \log\left(C\right))$ auxiliary stack space for \inlinecode{update()}, \inlinecode{query()}, and \inlinecode{at()}.
\end{compactitem}

\begin{lstlisting}[language=C++]
#include <algorithm>
#include <cstddef>

template<class T>
class SegTree2D {
  static const int R = 1000000000;
  static const int C = 1000000000;

  static T combine(const T &a, const T &b) { return std::min(a, b); }
  static T repeat_value(const T &v, long long area) { return v; }
  static T apply_delta(const T &v, const T &d) { return d; }

  struct InnerNode {
    T value;
    int low, high;
    InnerNode *left, *right;

    InnerNode(int lo, int hi, const T &v)
        : value(v), low(lo), high(hi), left(nullptr), right(nullptr) {}
  };

  struct OuterNode {
    InnerNode root;
    int low, high;
    OuterNode *left, *right;

    OuterNode(int lo, int hi, const T &v)
        : root(0, C, v), low(lo), high(hi), left(nullptr), right(nullptr) {}
  } *root;

  T init;

  // Helper variables for query() and update().
  int tgt_r1, tgt_c1, tgt_r2, tgt_c2;
  long long width;

  static long long length(int lo, int hi) { return hi - lo + 1LL; }

  void append_result(T &res, bool &found, const T &v) {
    res = found ? combine(res, v) : v;
    found = true;
  }

  T query(InnerNode *n, int qlo, int qhi, long long rows) {
    if (n == nullptr) {
      return repeat_value(init, rows * length(qlo, qhi));
    }
    int lo = n->low, hi = n->high, mid = lo + (hi - lo) / 2;
    T res;
    bool found = false;
    if (qlo < lo) {
      int missing_hi = std::min(qhi, lo - 1);
      append_result(res, found, repeat_value(init, rows * length(qlo, missing_hi)));
    }
    int overlap_lo = std::max(qlo, lo), overlap_hi = std::min(qhi, hi);
    if (overlap_lo <= overlap_hi) {
      if (overlap_lo == lo && overlap_hi == hi) {
        append_result(res, found, n->value);
      } else {
        if (overlap_lo <= mid) {
          append_result(res, found, query(n->left, overlap_lo, std::min(overlap_hi, mid), rows));
        }
        if (mid < overlap_hi) {
          append_result(
              res, found, query(n->right, std::max(overlap_lo, mid + 1), overlap_hi, rows)
          );
        }
      }
    }
    if (hi < qhi) {
      int missing_lo = std::max(qlo, hi + 1);
      append_result(res, found, repeat_value(init, rows * length(missing_lo, qhi)));
    }
    return res;
  }

  T query(OuterNode *n, int qlo, int qhi) {
    if (n == nullptr) {
      return repeat_value(init, width * length(qlo, qhi));
    }
    int lo = n->low, hi = n->high, mid = lo + (hi - lo) / 2;
    T res;
    bool found = false;
    if (qlo < lo) {
      int missing_hi = std::min(qhi, lo - 1);
      append_result(res, found, repeat_value(init, width * length(qlo, missing_hi)));
    }
    int overlap_lo = std::max(qlo, lo), overlap_hi = std::min(qhi, hi);
    if (overlap_lo <= overlap_hi) {
      if (overlap_lo == lo && overlap_hi == hi) {
        append_result(res, found, query(&(n->root), tgt_c1, tgt_c2, length(lo, hi)));
      } else {
        if (overlap_lo <= mid) {
          append_result(res, found, query(n->left, overlap_lo, std::min(overlap_hi, mid)));
        }
        if (mid < overlap_hi) {
          append_result(res, found, query(n->right, std::max(overlap_lo, mid + 1), overlap_hi));
        }
      }
    }
    if (hi < qhi) {
      int missing_lo = std::max(qlo, hi + 1);
      append_result(res, found, repeat_value(init, width * length(missing_lo, qhi)));
    }
    return res;
  }

  void update(InnerNode *n, int c, const T &d, bool leaf_row, long long rows) {
    int lo = n->low, hi = n->high, mid = lo + (hi - lo) / 2;
    if (lo == hi) {
      if (leaf_row) {
        n->value = apply_delta(n->value, d);
      } else {
        n->value = d;
      }
      return;
    }
    InnerNode *&target = (c <= mid) ? n->left : n->right;
    if (target == nullptr) {
      target = new InnerNode(c, c, repeat_value(init, rows));
    }
    if (target->low <= c && c <= target->high) {
      update(target, c, d, leaf_row, rows);
    } else {
      int split_lo = lo, split_hi = hi, split_mid = mid;
      do {
        if (c <= split_mid) {
          split_hi = split_mid;
        } else {
          split_lo = split_mid + 1;
        }
        split_mid = split_lo + (split_hi - split_lo) / 2;
      } while ((c <= split_mid) == (target->low <= split_mid));
      InnerNode *tmp =
          new InnerNode(split_lo, split_hi, repeat_value(init, rows * length(split_lo, split_hi)));
      if (target->low <= split_mid) {
        tmp->left = target;
      } else {
        tmp->right = target;
      }
      target = tmp;
      update(tmp, c, d, leaf_row, rows);
    }
    T left_value =
        (n->left != nullptr) ? n->left->value : repeat_value(init, rows * length(lo, mid));
    T right_value =
        (n->right != nullptr) ? n->right->value : repeat_value(init, rows * length(mid + 1, hi));
    n->value = combine(left_value, right_value);
  }

  void update(OuterNode *n, int r, int c, const T &d) {
    int lo = n->low, hi = n->high, mid = lo + (hi - lo) / 2;
    long long rows = length(lo, hi);
    if (lo == hi) {
      update(&(n->root), c, d, true, 1);
      return;
    }
    if (r <= mid) {
      if (n->left == nullptr) {
        n->left = new OuterNode(lo, mid, repeat_value(init, length(lo, mid) * length(0, C)));
      }
      update(n->left, r, c, d);
    } else {
      if (n->right == nullptr) {
        n->right =
            new OuterNode(mid + 1, hi, repeat_value(init, length(mid + 1, hi) * length(0, C)));
      }
      update(n->right, r, c, d);
    }
    T value = repeat_value(init, rows);
    if (n->left != nullptr || n->right != nullptr) {
      tgt_c1 = tgt_c2 = c;
      T left_value = (n->left != nullptr) ? query(&(n->left->root), c, c, length(lo, mid))
                                          : repeat_value(init, length(lo, mid));
      T right_value = (n->right != nullptr) ? query(&(n->right->root), c, c, length(mid + 1, hi))
                                            : repeat_value(init, length(mid + 1, hi));
      value = combine(left_value, right_value);
    }
    update(&(n->root), c, value, false, rows);
  }

  static void clean_up(InnerNode *n) {
    if (n != nullptr) {
      clean_up(n->left);
      clean_up(n->right);
      delete n;
    }
  }

  static void clean_up(OuterNode *n) {
    if (n != nullptr) {
      clean_up(n->root.left);
      clean_up(n->root.right);
      clean_up(n->left);
      clean_up(n->right);
      delete n;
    }
  }

 public:
  explicit SegTree2D(const T &v = T())
      : root(new OuterNode(0, R, repeat_value(v, length(0, R) * length(0, C)))), init(v) {}

  ~SegTree2D() { clean_up(root); }
  T at(int r, int c) { return query(r, c, r, c); }

  T query(int r1, int c1, int r2, int c2) {
    tgt_r1 = r1;
    tgt_c1 = c1;
    tgt_r2 = r2;
    tgt_c2 = c2;
    width = length(c1, c2);
    return query(root, r1, r2);
  }

  void update(int r, int c, const T &d) { update(root, r, c, d); }
};
\end{lstlisting}
\Needspace*{4\baselineskip}
\textbf{Example Usage}\par
\begin{lstlisting}[language=C++]
#include <cassert>
#include <iostream>
using namespace std;

int main() {
  SegTree2D<int> t(0);
  t.update(0, 0, 7);
  t.update(0, 1, 6);
  t.update(1, 0, 5);
  t.update(1, 1, 4);
  t.update(2, 1, 1);
  t.update(2, 2, 9);
  cout << "Values:" << endl;
  for (int i = 0; i < 3; i++) {
    for (int j = 0; j < 3; j++) {
      cout << t.at(i, j) << " ";
    }
    cout << endl;
  }
  assert(t.query(0, 0, 0, 1) == 6);
  assert(t.query(0, 0, 1, 0) == 5);
  assert(t.query(1, 1, 2, 2) == 0);
  assert(t.query(0, 0, 1000000000, 1000000000) == 0);
  t.update(500000000, 500000000, -100);
  assert(t.query(0, 0, 1000000000, 1000000000) == -100);
  return 0;
}
\end{lstlisting}
\begin{exampleoutput}
Values:
7 6 0
5 4 0
0 1 9
\end{exampleoutput}
\subsection{2D Range Tree}
Maintain a set of two-dimensional points while supporting queries for all points that fall inside given rectangular regions. This implementation uses \inlinecode{std::pair} to represent points, requiring operators \inlinecode{\textless{}} and \inlinecode{==} to be defined on the numeric template type.

\begin{compactitem}
  \item \inlinecode{RangeTree(lo, hi)} constructs a set from two random-access iterators to \inlinecode{std::pair} as a range \inlinecode{[lo, hi)} of points.
  \item \inlinecode{query(x1, y1, x2, y2, f)} calls the function \inlinecode{f(i, p)} on each point in the set that falls into the rectangular region consisting of rows from \inlinecode{x1} to \inlinecode{x2}, inclusive, and columns from \inlinecode{y1} to \inlinecode{y2}, inclusive. The first argument to \inlinecode{f} is the zero-based index of the point in the original range given to the constructor. The second argument is the point itself as an \inlinecode{std::pair}.
\end{compactitem}

\Needspace*{3\baselineskip}
\textbf{Time Complexity}:
\begin{compactitem}
  \item $O(n \log n)$ per call to the constructor, where $n$ is the number of points.
  \item $O(\log^{2}(n) + m)$ per call to \inlinecode{query()}, where $m$ is the number of points that are reported by the query.
\end{compactitem}

\Needspace*{3\baselineskip}
\textbf{Space Complexity}:
\begin{compactitem}
  \item $O(n \log n)$ for storage of the points.
  \item $O(\log^{2}(n))$ auxiliary stack space for \inlinecode{query()}.
\end{compactitem}

\begin{lstlisting}[language=C++]
#include <algorithm>
#include <iterator>
#include <utility>
#include <vector>

template<class T>
class RangeTree {
  using point = std::pair<T, T>;
  using colindex = std::pair<int, T>;

  std::vector<point> points;
  std::vector<std::vector<colindex>> columns;

  static inline bool comp1(const colindex &a, const colindex &b) { return a.second < b.second; }
  static inline bool comp2(const colindex &a, const T &v) { return a.second < v; }

  void build(int n, int lo, int hi) {
    if (points[lo].first == points[hi].first) {
      for (int i = lo; i <= hi; i++) {
        columns[n].emplace_back(i, points[i].second);
      }
      return;
    }
    int l = n * 2 + 1, r = n * 2 + 2, mid = lo + (hi - lo) / 2;
    build(l, lo, mid);
    build(r, mid + 1, hi);
    columns[n].resize(columns[l].size() + columns[r].size());
    std::merge(
        columns[l].begin(), columns[l].end(), columns[r].begin(), columns[r].end(),
        columns[n].begin(), comp1
    );
  }

  // Helper variables for query().
  T x1, y1, x2, y2;

  template<class Fn>
  void query(int n, int lo, int hi, Fn f) {
    if (points[hi].first < x1 || x2 < points[lo].first) {
      return;
    }
    if (!(points[lo].first < x1 || x2 < points[hi].first)) {
      if (!columns[n].empty() && !(y2 < y1)) {
        auto it = std::lower_bound(columns[n].begin(), columns[n].end(), y1, comp2);
        for (; it != columns[n].end() && it->second <= y2; ++it) {
          f(it->first, points[it->first]);
        }
      }
    } else if (lo != hi) {
      int mid = lo + (hi - lo) / 2;
      query(n * 2 + 1, lo, mid, f);
      query(n * 2 + 2, mid + 1, hi, f);
    }
  }

 public:
  template<class It>
  RangeTree(It lo, It hi) : points(lo, hi) {
    int n = std::distance(lo, hi);
    columns.resize(4 * n + 1);
    std::sort(points.begin(), points.end());
    build(0, 0, n - 1);
  }

  template<class Fn>
  void query(const T &x1, const T &y1, const T &x2, const T &y2, Fn f) {
    this->x1 = x1;
    this->y1 = y1;
    this->x2 = x2;
    this->y2 = y2;
    query(0, 0, points.size() - 1, f);
  }
};
\end{lstlisting}
\Needspace*{4\baselineskip}
\textbf{Example Usage}\par
\begin{lstlisting}[language=C++]
#include <iostream>
using namespace std;

void print(int i, const pair<int, int> &p) {
  cout << "(" << p.first << ", " << p.second << ") ";
}

int main() {
  vector<pair<int, int>> v{{1, 4},  {5, 4},  {2, 2},   {3, 1},   {6, -5},
                           {5, -1}, {3, -3}, {-1, -2}, {-1, -1}, {2, -1}};
  RangeTree<int> t(v.begin(), v.end());
  t.query(-1, -1, 2, 5, print);
  cout << endl;
  t.query(1, 1, 4, 8, print);
  cout << endl;
  return 0;
}
\end{lstlisting}
\begin{exampleoutput}
(-1, -1) (2, -1) (2, 2) (1, 4)
(1, 4) (2, 2) (3, 1)
\end{exampleoutput}
\subsection{K-d Tree (2D Range Query)}
Maintain a set of two-dimensional points while supporting queries for all points that fall inside given rectangular regions. This implementation uses \inlinecode{std::pair} to represent points, requiring operators \inlinecode{\textless{}} and \inlinecode{==} to be defined on the numeric template type.

\begin{compactitem}
  \item \inlinecode{KDTree(lo, hi)} constructs a set from two random-access iterators to \inlinecode{std::pair} as a range \inlinecode{[lo, hi)} of points.
  \item \inlinecode{query(x1, y1, x2, y2, f)} calls the function \inlinecode{f(i, p)} on each point in the set that falls into the rectangular region consisting of rows from \inlinecode{x1} to \inlinecode{x2}, inclusive, and columns from \inlinecode{y1} to \inlinecode{y2}, inclusive. The first argument to \inlinecode{f} is the zero-based index of the point in the original range given to the constructor. The second argument is the point itself as an \inlinecode{std::pair}.
\end{compactitem}

\Needspace*{3\baselineskip}
\textbf{Time Complexity}:
\begin{compactitem}
  \item $O(n \log n)$ per call to the constructor, where $n$ is the number of points.
  \item $O(\log\left(n\right) + m)$ on average per call to \inlinecode{query()}, where $m$ is the number of points that are reported by the query.
\end{compactitem}

\Needspace*{3\baselineskip}
\textbf{Space Complexity}:
\begin{compactitem}
  \item $O(n)$ for storage of the points.
  \item $O(\log n)$ auxiliary stack space for \inlinecode{query()}.
\end{compactitem}

\begin{lstlisting}[language=C++]
#include <algorithm>
#include <utility>
#include <vector>

template<class T>
class KDTree {
  using point = std::pair<T, T>;

  static inline bool comp1(const point &a, const point &b) { return a.first < b.first; }
  static inline bool comp2(const point &a, const point &b) { return a.second < b.second; }

  std::vector<point> tree, minp, maxp;
  std::vector<int> l_index, h_index;

  void build(int lo, int hi, bool div_x) {
    if (lo >= hi) {
      return;
    }
    int mid = lo + (hi - lo) / 2;
    std::nth_element(
        tree.begin() + lo, tree.begin() + mid, tree.begin() + hi, div_x ? comp1 : comp2
    );
    l_index[mid] = lo;
    h_index[mid] = hi;
    minp[mid].first = maxp[mid].first = tree[lo].first;
    minp[mid].second = maxp[mid].second = tree[lo].second;
    for (int i = lo + 1; i < hi; i++) {
      minp[mid].first = std::min(minp[mid].first, tree[i].first);
      minp[mid].second = std::min(minp[mid].second, tree[i].second);
      maxp[mid].first = std::max(maxp[mid].first, tree[i].first);
      maxp[mid].second = std::max(maxp[mid].second, tree[i].second);
    }
    build(lo, mid, !div_x);
    build(mid + 1, hi, !div_x);
  }

  // Helper variables for query().
  T x1, y1, x2, y2;

  template<class Fn>
  void query(int lo, int hi, Fn f) {
    if (lo >= hi) {
      return;
    }
    int mid = lo + (hi - lo) / 2;
    T ax = minp[mid].first, ay = minp[mid].second;
    T bx = maxp[mid].first, by = maxp[mid].second;
    if (x2 < ax || bx < x1 || y2 < ay || by < y1) {
      return;
    }
    if (!(ax < x1 || x2 < bx || ay < y1 || y2 < by)) {
      for (int i = l_index[mid]; i < h_index[mid]; i++) {
        f(tree[i]);
      }
      return;
    }
    query(lo, mid, f);
    query(mid + 1, hi, f);
    if (tree[mid].first < x1 || x2 < tree[mid].first || tree[mid].second < y1 ||
        y2 < tree[mid].second) {
      return;
    }
    f(tree[mid]);
  }

 public:
  template<class It>
  KDTree(It lo, It hi) : tree(lo, hi) {
    int n = std::distance(lo, hi);
    l_index.resize(n);
    h_index.resize(n);
    minp.resize(n);
    maxp.resize(n);
    build(0, n, true);
  }

  template<class Fn>
  void query(const T &x1, const T &y1, const T &x2, const T &y2, Fn f) {
    this->x1 = x1;
    this->y1 = y1;
    this->x2 = x2;
    this->y2 = y2;
    query(0, tree.size(), f);
  }
};
\end{lstlisting}
\Needspace*{4\baselineskip}
\textbf{Example Usage}\par
\begin{lstlisting}[language=C++]
#include <iostream>
using namespace std;

void print(const pair<int, int> &p) {
  cout << "(" << p.first << ", " << p.second << ") ";
}

int main() {
  vector<pair<int, int>> v{{1, 4},  {5, 4},  {2, 2},   {3, 1},   {6, -5},
                           {5, -1}, {3, -3}, {-1, -2}, {-1, -1}, {2, -1}};
  KDTree<int> t(v.begin(), v.end());
  t.query(-1, -1, 2, 5, print);
  cout << endl;
  t.query(1, 1, 4, 8, print);
  cout << endl;
  return 0;
}
\end{lstlisting}
\begin{exampleoutput}
(2, -1) (1, 4) (2, 2) (-1, -1)
(1, 4) (2, 2) (3, 1)
\end{exampleoutput}
\subsection{K-d Tree (Nearest Neighbor)}
Maintain a set of two-dimensional points while supporting queries for the closest point in the set to a given query point. This implementation uses \inlinecode{std::pair} to represent points, requiring operators \inlinecode{\textless{}}, \inlinecode{==}, \inlinecode{-}, and \inlinecode{long double} casting to be defined on the numeric template type.

\begin{compactitem}
  \item \inlinecode{KDTree(lo, hi)} constructs a set from two random-access iterators to \inlinecode{std::pair} as a range \inlinecode{[lo, hi)} of points.
  \item \inlinecode{nearest(x, y, can\_equal)} returns a point in the set that is closest to (\inlinecode{x}, \inlinecode{y}) by Euclidean distance. This may be equal to (\inlinecode{x}, \inlinecode{y}) only if \inlinecode{can\_equal} is \inlinecode{true}.
\end{compactitem}

\Needspace*{3\baselineskip}
\textbf{Time Complexity}:
\begin{compactitem}
  \item $O(n \log n)$ per call to the constructor, where $n$ is the number of points.
  \item $O(\log n)$ on average per call to \inlinecode{nearest()}.
\end{compactitem}

\Needspace*{3\baselineskip}
\textbf{Space Complexity}:
\begin{compactitem}
  \item $O(n)$ for storage of the points.
  \item $O(\log n)$ auxiliary stack space for \inlinecode{nearest()}.
\end{compactitem}

\begin{lstlisting}[language=C++]
#include <algorithm>
#include <cfloat>
#include <stdexcept>
#include <utility>
#include <vector>

template<class T>
class KDTree {
  using point = std::pair<T, T>;

  static inline bool comp1(const point &a, const point &b) { return a.first < b.first; }
  static inline bool comp2(const point &a, const point &b) { return a.second < b.second; }

  std::vector<point> tree;
  std::vector<bool> div_x;

  void build(int lo, int hi) {
    if (lo >= hi) {
      return;
    }
    int mid = lo + (hi - lo) / 2;
    T minx, maxx, miny, maxy;
    minx = maxx = tree[lo].first;
    miny = maxy = tree[lo].second;
    for (int i = lo + 1; i < hi; i++) {
      minx = std::min(minx, tree[i].first);
      miny = std::min(miny, tree[i].second);
      maxx = std::max(maxx, tree[i].first);
      maxy = std::max(maxy, tree[i].second);
    }
    div_x[mid] = !((maxx - minx) < (maxy - miny));
    std::nth_element(
        tree.begin() + lo, tree.begin() + mid, tree.begin() + hi, div_x[mid] ? comp1 : comp2
    );
    if (lo + 1 == hi) {
      return;
    }
    build(lo, mid);
    build(mid + 1, hi);
  }

  // Helper variables for nearest().
  long double min_dist;
  int id;

  void nearest(int lo, int hi, const T &x, const T &y, bool can_equal) {
    if (lo >= hi) {
      return;
    }
    int mid = lo + (hi - lo) / 2;
    T dx = x - tree[mid].first, dy = y - tree[mid].second;
    long double d = dx * static_cast<long double>(dx) + dy * static_cast<long double>(dy);
    if (d < min_dist && (can_equal || d != 0)) {
      min_dist = d;
      id = mid;
    }
    if (lo + 1 == hi) {
      return;
    }
    d = static_cast<long double>((div_x[mid] ? dx : dy));
    int l1 = lo, r1 = mid, l2 = mid + 1, r2 = hi;
    if (d > 0) {
      std::swap(l1, l2);
      std::swap(r1, r2);
    }
    nearest(l1, r1, x, y, can_equal);
    if (d * static_cast<long double>(d) < min_dist) {
      nearest(l2, r2, x, y, can_equal);
    }
  }

 public:
  template<class It>
  KDTree(It lo, It hi) : tree(lo, hi) {
    int n = std::distance(lo, hi);
    if (n <= 1) {
      throw std::runtime_error("K-d tree must be have at least 2 points.");
    }
    div_x.resize(n);
    build(0, n);
  }

  point nearest(const T &x, const T &y, bool can_equal = true) {
    min_dist = LDBL_MAX;
    nearest(0, tree.size(), x, y, can_equal);
    return tree[id];
  }
};
\end{lstlisting}
\Needspace*{4\baselineskip}
\textbf{Example Usage}\par
\begin{lstlisting}[language=C++]
#include <cassert>
using namespace std;

int main() {
  vector<pair<int, int>> p{{0, 2}, {0, 3}, {-1, 0}};
  KDTree<int> t(p.begin(), p.end());
  assert(t.nearest(0, 2, true) == (pair<int, int>{0, 2}));
  assert(t.nearest(0, 2, false) == (pair<int, int>{0, 3}));
  assert(t.nearest(0, 0) == (pair<int, int>{-1, 0}));
  assert(t.nearest(-10000, 0) == (pair<int, int>{-1, 0}));
  return 0;
}
\end{lstlisting}

\section{Fenwick Trees}
\setcounter{section}{6}
\setcounter{subsection}{0}
\subsection{Fenwick Tree (Simple)}
Maintain an array of numerical type, allowing for updates of individual indices (point update) and queries for the sum of contiguous sub-arrays (range queries). This implementation assumes that the array is 1-based (i.e. has valid indices from 1 to \inlinecode{n}, inclusive).

\begin{compactitem}
  \item \inlinecode{initialize(n)} resets the dat structure.
  \item \inlinecode{vals[i]} stores the value at index \inlinecode{i}.
  \item \inlinecode{add(i, x)} adds \inlinecode{x} to the value at index \inlinecode{i}.
  \item \inlinecode{set(i, x)} assigns the value at index \inlinecode{i} to \inlinecode{x}.
  \item \inlinecode{sum(hi)} returns the sum of all values at indices from 1 to \inlinecode{hi}, inclusive.
  \item \inlinecode{sum(lo, hi)} returns the sum of all values at indices from \inlinecode{lo} to \inlinecode{hi}, inclusive.
\end{compactitem}

\Needspace*{3\baselineskip}
\textbf{Time Complexity}:
\begin{compactitem}
  \item $O(n)$ per call to \inlinecode{initialize(n)}, where $n$ is the size of the array.
  \item $O(\log n)$ per call to all other operations.
\end{compactitem}

\Needspace*{3\baselineskip}
\textbf{Space Complexity}:
\begin{compactitem}
  \item $O(n)$ for storage of the array elements.
  \item $O(1)$ auxiliary for all operations.
\end{compactitem}

\begin{lstlisting}[language=C++]
#include <vector>

std::vector<int> vals, tree;

void initialize(int n) {
  vals.assign(n + 2, 0);
  tree.assign(n + 2, 0);
}

void add(int i, int x) {
  vals[i] += x;
  for (; i < static_cast<int>(tree.size()); i += i & -i) {
    tree[i] += x;
  }
}

void set(int i, int x) {
  add(i, x - vals[i]);
}

int sum(int hi) {
  int res = 0;
  for (; hi > 0; hi -= hi & -hi) {
    res += tree[hi];
  }
  return res;
}

int sum(int lo, int hi) {
  return sum(hi) - sum(lo - 1);
}
\end{lstlisting}
\Needspace*{4\baselineskip}
\textbf{Example Usage}\par
\begin{lstlisting}[language=C++]
#include <cassert>
#include <iostream>
using namespace std;

int main() {
  vector<int> v{10, 1, 2, 3, 4};
  int n = 5;
  initialize(n);
  for (int i = 1; i < n; i++) {
    set(i, v[i - 1]);
  }
  add(1, -5);
  cout << "Values: ";
  for (int i = 1; i < n; i++) {
    cout << vals[i] << " ";
  }
  cout << endl;
  assert(sum(2, 4) == 6);
  return 0;
}
\end{lstlisting}
\begin{exampleoutput}
Values: 5 1 2 3 4
\end{exampleoutput}
\subsection{Fenwick Tree (Range Update, Point Query)}
Maintain an array of numerical type, allowing for contiguous sub-arrays to be simultaneously incremented by arbitrary values (range update) and values at individual indices to be queried (point query). This implementation assumes that the array is 0-based (i.e. has valid indices from 0 to \inlinecode{size() - 1}, inclusive).

\begin{compactitem}
  \item \inlinecode{size()} returns the size of the array.
  \item \inlinecode{at(i)} returns the value at index \inlinecode{i}.
  \item \inlinecode{add(i, x)} adds \inlinecode{x} to the value at index \inlinecode{i}.
  \item \inlinecode{add(lo, hi, x)} adds \inlinecode{x} to the values at all indices from \inlinecode{lo} to \inlinecode{hi}, inclusive.
\end{compactitem}

\Needspace*{3\baselineskip}
\textbf{Time Complexity}:
\begin{compactitem}
  \item $O(n)$ per call to the constructor, where $n$ is the size of the array.
  \item $O(1)$ per call to \inlinecode{size()}.
  \item $O(\log n)$ per call to \inlinecode{at()} and both \inlinecode{add()} functions.
\end{compactitem}

\Needspace*{3\baselineskip}
\textbf{Space Complexity}:
\begin{compactitem}
  \item $O(n)$ for storage of the array elements.
  \item $O(1)$ auxiliary for all operations.
\end{compactitem}

\begin{lstlisting}[language=C++]
#include <vector>

template<class T>
class FenwickTree {
  int len;
  std::vector<T> tree;

 public:
  explicit FenwickTree(int n) : len(n), tree(n + 2) {}

  int size() const { return len; }

  T at(int i) const {
    T res = 0;
    for (i++; i > 0; i -= i & -i) {
      res += tree[i];
    }
    return res;
  }

  void add(int i, const T &x) {
    for (i++; i <= len + 1; i += i & -i) {
      tree[i] += x;
    }
  }

  void add(int lo, int hi, const T &x) {
    add(lo, x);
    add(hi + 1, -x);
  }
};
\end{lstlisting}
\Needspace*{4\baselineskip}
\textbf{Example Usage}\par
\begin{lstlisting}[language=C++]
#include <iostream>
using namespace std;

int main() {
  FenwickTree<int> t(5);
  t.add(0, 1, 5);
  t.add(1, 2, 5);
  t.add(2, 4, 10);
  cout << "Values: ";
  for (int i = 0; i < t.size(); i++) {
    cout << t.at(i) << " ";
  }
  cout << endl;
  return 0;
}
\end{lstlisting}
\begin{exampleoutput}
Values: 5 10 15 10 10
\end{exampleoutput}
\subsection{Fenwick Tree (Point Update, Range Query)}
Maintain an array of numerical type, allowing for updates of individual indices (point update) and queries for the sum of contiguous sub-arrays (range queries). This implementation assumes that the array is 0-based (i.e. has valid indices from 0 to \inlinecode{size() - 1}, inclusive).

\begin{compactitem}
  \item \inlinecode{size()} returns the size of the array.
  \item \inlinecode{at(i)} returns the value at index \inlinecode{i}.
  \item \inlinecode{add(i, x)} adds \inlinecode{x} to the value at index \inlinecode{i}.
  \item \inlinecode{set(i, x)} assigns the value at index \inlinecode{i} to \inlinecode{x}.
  \item \inlinecode{sum(hi)} returns the sum of all values at indices from 0 to \inlinecode{hi}, inclusive.
  \item \inlinecode{sum(lo, hi)} returns the sum of all values at indices from \inlinecode{lo} to \inlinecode{hi}, inclusive.
\end{compactitem}

\Needspace*{3\baselineskip}
\textbf{Time Complexity}:
\begin{compactitem}
  \item $O(n)$ per call to the constructor, where $n$ is the size of the array.
  \item $O(1)$ per call to \inlinecode{size()} and \inlinecode{at()}.
  \item $O(\log n)$ per call to \inlinecode{add()}, \inlinecode{set()}, and both \inlinecode{sum()} functions.
\end{compactitem}

\Needspace*{3\baselineskip}
\textbf{Space Complexity}:
\begin{compactitem}
  \item $O(n)$ for storage of the array elements.
  \item $O(1)$ auxiliary for all operations.
\end{compactitem}

\begin{lstlisting}[language=C++]
#include <vector>

template<class T>
class FenwickTree {
  int len;
  std::vector<T> vals, tree;

 public:
  explicit FenwickTree(int n) : len(n), vals(n + 1), tree(n + 1) {}

  int size() const { return len; }
  T at(int i) const { return vals[i + 1]; }

  void add(int i, const T &x) {
    vals[++i] += x;
    for (; i <= len; i += i & -i) {
      tree[i] += x;
    }
  }

  void set(int i, const T &x) {
    T inc = x - vals[i + 1];
    add(i, inc);
  }

  T sum(int hi) {
    T res = 0;
    for (hi++; hi > 0; hi -= hi & -hi) {
      res += tree[hi];
    }
    return res;
  }

  T sum(int lo, int hi) { return sum(hi) - sum(lo - 1); }
};
\end{lstlisting}
\Needspace*{4\baselineskip}
\textbf{Example Usage}\par
\begin{lstlisting}[language=C++]
#include <cassert>
#include <iostream>
using namespace std;

int main() {
  vector<int> a{10, 1, 2, 3, 4};
  FenwickTree<int> t(5);
  for (int i = 0; i < static_cast<int>(a.size()); i++) {
    t.set(i, a[i]);
  }
  t.add(0, -5);
  cout << "Values: ";
  for (int i = 0; i < t.size(); i++) {
    cout << t.at(i) << " ";
  }
  cout << endl;
  assert(t.sum(1, 3) == 6);
  return 0;
}
\end{lstlisting}
\begin{exampleoutput}
Values: 5 1 2 3 4
\end{exampleoutput}
\subsection{Fenwick Tree (Range Update, Range Query)}
Maintain an array of numerical type, allowing for contiguous sub-arrays to be simultaneously incremented by arbitrary values (range update) and queries for the sum of contiguous sub-arrays (range query). This implementation assumes that the array is 0-based (i.e. has valid indices from 0 to \inlinecode{size() - 1}, inclusive).

\begin{compactitem}
  \item \inlinecode{size()} returns the size of the array.
  \item \inlinecode{at(i)} returns the value at index \inlinecode{i}.
  \item \inlinecode{add(i, x)} increments the value at index \inlinecode{i} by \inlinecode{x}.
  \item \inlinecode{add(lo, hi, x)} adds \inlinecode{x} to the values at all indices from \inlinecode{lo} to \inlinecode{hi}, inclusive.
  \item \inlinecode{set(i, x)} assigns the value at index \inlinecode{i} to \inlinecode{x}.
  \item \inlinecode{sum(hi)} returns the sum of all values at indices from 0 to \inlinecode{hi}, inclusive.
  \item \inlinecode{sum(lo, hi)} returns the sum of all values at indices from \inlinecode{lo} to \inlinecode{hi}, inclusive.
\end{compactitem}

\Needspace*{3\baselineskip}
\textbf{Time Complexity}:
\begin{compactitem}
  \item $O(n)$ per call to the constructor, where $n$ is the size of the array.
  \item $O(1)$ per call to \inlinecode{size()}.
  \item $O(\log n)$ per call to all other operations.
\end{compactitem}

\Needspace*{3\baselineskip}
\textbf{Space Complexity}:
\begin{compactitem}
  \item $O(n)$ for storage of the array elements.
  \item $O(1)$ auxiliary for all operations.
\end{compactitem}

\begin{lstlisting}[language=C++]
#include <vector>

template<class T>
class FenwickTree {
  int len;
  std::vector<T> t1, t2;

  T sum(const std::vector<T> &tree, int i) {
    T res = 0;
    for (; i != 0; i -= i & -i) {
      res += tree[i];
    }
    return res;
  }

  void add(std::vector<T> &tree, int i, const T &x) {
    for (; i <= len + 1; i += i & -i) {
      tree[i] += x;
    }
  }

 public:
  explicit FenwickTree(int n) : len(n), t1(n + 2), t2(n + 2) {}

  int size() const { return len; }

  void add(int lo, int hi, const T &x) {
    lo++;
    hi++;
    add(t1, lo, x);
    add(t1, hi + 1, -x);
    add(t2, lo, x * (lo - 1));
    add(t2, hi + 1, -x * hi);
  }

  void add(int i, const T &x) { return add(i, i, x); }
  void set(int i, const T &x) { add(i, x - at(i)); }

  T sum(int hi) {
    hi++;
    return hi * sum(t1, hi) - sum(t2, hi);
  }

  T sum(int lo, int hi) { return sum(hi) - sum(lo - 1); }
  T at(int i) { return sum(i, i); }
};
\end{lstlisting}
\Needspace*{4\baselineskip}
\textbf{Example Usage}\par
\begin{lstlisting}[language=C++]
#include <cassert>
#include <iostream>
using namespace std;

int main() {
  vector<int> a{10, 1, 2, 3, 4};
  FenwickTree<int> t(5);
  for (int i = 0; i < t.size(); i++) {
    t.set(i, a[i]);
  }
  t.add(0, 2, 5);
  t.set(3, -5);
  cout << "Values: ";
  for (int i = 0; i < t.size(); i++) {
    cout << t.at(i) << " ";
  }
  cout << endl;
  assert(t.sum(0, 4) == 27);
  return 0;
}
\end{lstlisting}
\begin{exampleoutput}
Values: 15 6 7 -5 4
\end{exampleoutput}
\subsection{Fenwick Tree (Compressed)}
Maintain an array of numerical type, allowing for contiguous sub-arrays to be simultaneously incremented by arbitrary values (range update) and queries for the sum of contiguous sub-arrays (range query). This implementation uses \inlinecode{std::unordered\_map} for coordinate compression, allowing for large indices to be accessed with efficient space complexity. That is, all array indices from 0 to \inlinecode{N}, inclusive, are accessible.

\begin{compactitem}
  \item \inlinecode{at(i)} returns the value at index \inlinecode{i}.
  \item \inlinecode{add(i, x)} adds \inlinecode{x} to the value at index \inlinecode{i}.
  \item \inlinecode{add(lo, hi, x)} adds \inlinecode{x} to the values at all indices from \inlinecode{lo} to \inlinecode{hi}, inclusive.
  \item \inlinecode{set(i, x)} assigns the value at index \inlinecode{i} to \inlinecode{x}.
  \item \inlinecode{sum(hi)} returns the sum of all values at indices from 0 to \inlinecode{hi}, inclusive.
  \item \inlinecode{sum(lo, hi)} returns the sum of all values at indices from \inlinecode{lo} to \inlinecode{hi}, inclusive.
\end{compactitem}

\textbf{Time Complexity}: $O(\log n)$ per call to all member functions, where $n$ is the number of distinct indices that have been accessed.

\Needspace*{3\baselineskip}
\textbf{Space Complexity}:
\begin{compactitem}
  \item $O(n \log n)$ for storage of the array elements, where $n$ is the number of distinct indices that have been accessed across all of the operations so far.
  \item $O(1)$ auxiliary for all operations.
\end{compactitem}

\begin{lstlisting}[language=C++]
#include <unordered_map>

template<class T>
class FenwickTree {
  static const int N = 1000000001;
  std::unordered_map<int, T> tmul, tadd;

  void add_helper(int at, int mul, T add) {
    for (int i = at; i <= N; i |= i + 1) {
      tmul[i] += mul;
      tadd[i] += add;
    }
  }

 public:
  void add(int lo, int hi, const T &x) {
    add_helper(lo, x, -x * (lo - 1));
    add_helper(hi, -x, x * hi);
  }

  void add(int i, const T &x) { add(i, i, x); }
  void set(int i, const T &x) { add(i, x - at(i)); }

  T sum(int hi) {
    T mul = 0, add = 0;
    for (int i = hi; i >= 0; i = (i & (i + 1)) - 1) {
      if (auto it = tmul.find(i); it != tmul.end()) {
        mul += it->second;
      }
      if (auto it = tadd.find(i); it != tadd.end()) {
        add += it->second;
      }
    }
    return mul * hi + add;
  }

  T sum(int lo, int hi) { return sum(hi) - sum(lo - 1); }
  T at(int i) { return sum(i, i); }
};
\end{lstlisting}
\Needspace*{4\baselineskip}
\textbf{Example Usage}\par
\begin{lstlisting}[language=C++]
#include <cassert>
#include <iostream>
#include <vector>
using namespace std;

int main() {
  vector<int> a{10, 1, 2, 3, 4};
  FenwickTree<int> t;
  for (int i = 0; i < static_cast<int>(a.size()); i++) {
    t.set(i, a[i]);
  }
  t.add(0, 2, 5);
  t.set(3, -5);
  cout << "Values: ";
  for (int i = 0; i < 5; i++) {
    cout << t.at(i) << " ";
  }
  cout << endl;
  assert(t.sum(0, 4) == 27);
  t.add(500000001, 500000010, 3);
  t.add(500000011, 500000015, 5);
  t.set(500000000, 10);
  assert(t.sum(0, 1000000000) == 92);
  return 0;
}
\end{lstlisting}
\begin{exampleoutput}
Values: 15 6 7 -5 4
\end{exampleoutput}
\subsection{Fenwick Tree (Order Statistic)}
Maintains frequencies over a 1-based integer domain and finds order statistics by binary lifting on a Fenwick tree. This is useful for dynamic multisets of bounded integer values, coordinate-compressed kth-element queries, and online rank queries.

\begin{compactitem}
  \item \inlinecode{initialize(n)} resets all frequencies.
  \item \inlinecode{add(i, delta)} adds \inlinecode{delta} to the frequency of value/index \inlinecode{i}.
  \item \inlinecode{sum(i)} returns the total frequency over indices \inlinecode{[1, i]}.
  \item \inlinecode{kth(k)} returns the smallest index \inlinecode{i} such that \inlinecode{sum(i) \textgreater{}= k}. The argument \inlinecode{k} is 1-based and must satisfy \inlinecode{1 \textless{}= k \textless{}= sum(n)}.
\end{compactitem}

\Needspace*{3\baselineskip}
\textbf{Time Complexity}:
\begin{compactitem}
  \item $O(n)$ per call to \inlinecode{initialize()}.
  \item $O(\log n)$ per call to \inlinecode{add(i, delta)}, \inlinecode{sum(i)}, and \inlinecode{kth(k)}.
\end{compactitem}

\textbf{Space Complexity}: $O(n)$ for the Fenwick tree.

\begin{lstlisting}[language=C++]
#include <algorithm>
#include <vector>

std::vector<int> tree;

void initialize(int n) {
  tree.assign(n + 1, 0);
}

void add(int i, int delta) {
  for (; i < tree.size(); i += i & -i) {
    tree[i] += delta;
  }
}

int sum(int i) {
  int res = 0;
  for (; i > 0; i -= i & -i) {
    res += tree[i];
  }
  return res;
}

int kth(int k) {
  int idx = 0, bits = 1;
  while (bits * 2 < tree.size()) {
    bits *= 2;
  }
  for (; bits > 0; bits >>= 1) {
    int next = idx + bits;
    if (next < tree.size() && tree[next] < k) {
      idx = next;
      k -= tree[next];
    }
  }
  return idx + 1;
}
\end{lstlisting}
\Needspace*{4\baselineskip}
\textbf{Example Usage}\par
\begin{lstlisting}[language=C++]
#include <cassert>

int main() {
  initialize(100);
  add(2, 1);
  add(4, 3);
  add(7, 1);
  assert(sum(4) == 4);
  assert(kth(1) == 2);
  assert(kth(2) == 4);
  assert(kth(4) == 4);
  assert(kth(5) == 7);
  return 0;
}
\end{lstlisting}
\subsection{2D Fenwick Tree (Simple)}
Maintain a 2D array of numerical type, allowing for updates of individual cells in the matrix (point update) and queries for the sum of rectangular sub-matrices (range query). This implementation assumes that array dimensions are 1-based (i.e. rows have valid indices from 1 to \inlinecode{R}, inclusive, and columns have valid indices from 1 to \inlinecode{C}, inclusive).

\begin{compactitem}
  \item \inlinecode{initialize(R, C)} resets the data structure.
  \item \inlinecode{vals[r][c]} stores the value at index (\inlinecode{r}, \inlinecode{c}).
  \item \inlinecode{add(r, c, x)} adds \inlinecode{x} to the value at index (\inlinecode{r}, \inlinecode{c}).
  \item \inlinecode{set(r, c, x)} assigns \inlinecode{x} to the value at index (\inlinecode{r}, \inlinecode{c}).
  \item \inlinecode{sum(r, c)} returns the sum of the rectangle with upper-left corner (1, 1) and lower-right corner (\inlinecode{r}, \inlinecode{c}).
  \item \inlinecode{sum(r1, c1, r2, c2)} returns the sum of the rectangle with upper-left corner (\inlinecode{r1}, \inlinecode{c1}) and lower-right corner (\inlinecode{r2}, \inlinecode{c2}).
\end{compactitem}

\Needspace*{3\baselineskip}
\textbf{Time Complexity}:
\begin{compactitem}
  \item $O(R \cdot C)$ per call to \inlinecode{initialize(R, C)}.
  \item $O(\log\left(R\right) \cdot \log\left(C\right))$ per call to all other operations.
\end{compactitem}

\Needspace*{3\baselineskip}
\textbf{Space Complexity}:
\begin{compactitem}
  \item $O(R \cdot C)$ for storage of the array elements.
  \item $O(1)$ auxiliary for all operations.
\end{compactitem}

\begin{lstlisting}[language=C++]
#include <algorithm>
#include <vector>

std::vector<std::vector<int>> vals, tree;

void initialize(int R, int C) {
  vals.assign(R + 2, std::vector<int>(C + 2, 0));
  tree.assign(R + 2, std::vector<int>(C + 2, 0));
}

void add(int r, int c, int x) {
  vals[r][c] += x;
  for (int i = r; i < static_cast<int>(tree.size()); i += i & -i) {
    for (int j = c; j < static_cast<int>(tree[0].size()); j += j & -j) {
      tree[i][j] += x;
    }
  }
}

void set(int r, int c, int x) {
  add(r, c, x - vals[r][c]);
}

int sum(int r, int c) {
  int res = 0;
  for (int i = r; i > 0; i -= i & -i) {
    for (int j = c; j > 0; j -= j & -j) {
      res += tree[i][j];
    }
  }
  return res;
}

int sum(int r1, int c1, int r2, int c2) {
  return sum(r2, c2) + sum(r1 - 1, c1 - 1) - sum(r1 - 1, c2) - sum(r2, c1 - 1);
}
\end{lstlisting}
\Needspace*{4\baselineskip}
\textbf{Example Usage}\par
\begin{lstlisting}[language=C++]
#include <cassert>
#include <iostream>
using namespace std;

int main() {
  initialize(3, 3);
  set(1, 1, 5);
  set(1, 2, 6);
  set(2, 1, 7);
  add(3, 3, 9);
  add(2, 1, -4);
  cout << "Values:" << endl;
  for (int i = 1; i <= 3; i++) {
    for (int j = 1; j <= 3; j++) {
      cout << vals[i][j] << " ";
    }
    cout << endl;
  }
  assert(sum(1, 1, 1, 2) == 11);
  assert(sum(1, 1, 2, 1) == 8);
  assert(sum(1, 1, 3, 3) == 23);
  return 0;
}
\end{lstlisting}
\begin{exampleoutput}
Values:
5 6 0
3 0 0
0 0 9
\end{exampleoutput}
\subsection{2D Fenwick Tree (Compressed)}
Maintain a 2D array of numerical type, allowing for rectangular sub-matrices to be simultaneously incremented by arbitrary values (range update) and queries for the sum of rectangular sub-matrices (range query). This implementation uses \inlinecode{std::unordered\_map} for coordinate compression, allowing for large indices to be accessed with efficient space complexity. That is, rows have valid indices from 0 to \inlinecode{R}, inclusive, and columns have valid indices from 0 to \inlinecode{C}, inclusive.

\begin{compactitem}
  \item \inlinecode{add(r, c, x)} adds \inlinecode{x} to the value at index (\inlinecode{r}, \inlinecode{c}).
  \item \inlinecode{add(r1, c1, r2, c2, x)} adds \inlinecode{x} to all indices in the rectangle with upper-left corner (\inlinecode{r1}, \inlinecode{c1}) and lower-right corner (\inlinecode{r2}, \inlinecode{c2}).
  \item \inlinecode{set(r, c, x)} assigns \inlinecode{x} to the value at index (\inlinecode{r}, \inlinecode{c}).
  \item \inlinecode{sum(r, c)} returns the sum of the rectangle with upper-left corner (0, 0) and lower-right corner (\inlinecode{r}, \inlinecode{c}).
  \item \inlinecode{sum(r1, c1, r2, c2)} returns the sum of the rectangle with upper-left corner (\inlinecode{r1}, \inlinecode{c1}) and lower-right corner (\inlinecode{r2}, \inlinecode{c2}).
  \item \inlinecode{at(r, c)} returns the value at index (\inlinecode{r}, \inlinecode{c}).
\end{compactitem}

\textbf{Time Complexity}: $O(\log\left(R\right) \cdot \log\left(C\right))$ per call to all member functions.

\Needspace*{3\baselineskip}
\textbf{Space Complexity}:
\begin{compactitem}
  \item $O(n \cdot \log\left(R\right) \cdot \log\left(C\right))$ for storage of the array elements, where $n$ is the number of distinct indices that have been accessed across all of the operations so far.
  \item $O(1)$ auxiliary for all operations.
\end{compactitem}

\begin{lstlisting}[language=C++]
#include <unordered_map>
#include <utility>

template<class T>
class FenwickTree2D {
  static const int R = 1000000001;
  static const int C = 1000000001;

  std::unordered_map<long long, T> t1, t2, t3, t4;

  template<class Map>
  static T get(const Map &tree, int r, int c) {
    auto it = tree.find(static_cast<long long>(r) * C + c);
    return it == tree.end() ? T() : it->second;
  }

  template<class Map>
  static void add(Map &tree, int r, int c, const T &x) {
    for (int i = r + 1; i <= R; i += i & -i) {
      for (int j = c + 1; j <= C; j += j & -j) {
        tree[static_cast<long long>(i) * C + j] += x;
      }
    }
  }

  void add_helper(int r, int c, const T &x) {
    add(t1, 0, 0, x);
    add(t1, 0, c, -x);
    add(t2, 0, c, x * c);
    add(t1, r, 0, -x);
    add(t3, r, 0, x * r);
    add(t1, r, c, x);
    add(t2, r, c, -x * c);
    add(t3, r, c, -x * r);
    add(t4, r, c, x * r * c);
  }

 public:
  void add(int r1, int c1, int r2, int c2, const T &x) {
    add_helper(r2 + 1, c2 + 1, x);
    add_helper(r1, c2 + 1, -x);
    add_helper(r2 + 1, c1, -x);
    add_helper(r1, c1, x);
  }

  void add(int r, int c, const T &x) { add(r, c, r, c, x); }
  void set(int r, int c, const T &x) { add(r, c, x - at(r, c)); }

  T sum(int r, int c) {
    r++;
    c++;
    T s1 = 0, s2 = 0, s3 = 0, s4 = 0;
    for (int i = r; i > 0; i -= i & -i) {
      for (int j = c; j > 0; j -= j & -j) {
        s1 += get(t1, i, j);
        s2 += get(t2, i, j);
        s3 += get(t3, i, j);
        s4 += get(t4, i, j);
      }
    }
    return s1 * r * c + s2 * r + s3 * c + s4;
  }

  T sum(int r1, int c1, int r2, int c2) {
    return sum(r2, c2) + sum(r1 - 1, c1 - 1) - sum(r1 - 1, c2) - sum(r2, c1 - 1);
  }

  T at(int r, int c) { return sum(r, c, r, c); }
};
\end{lstlisting}
\Needspace*{4\baselineskip}
\textbf{Example Usage}\par
\begin{lstlisting}[language=C++]
#include <cassert>
#include <iostream>
using namespace std;

int main() {
  FenwickTree2D<int> t;
  t.set(0, 0, 5);
  t.set(0, 1, 6);
  t.set(1, 0, 7);
  t.add(2, 2, 9);
  t.add(1, 0, -4);
  t.add(1, 1, 2, 2, 5);
  cout << "Values:" << endl;
  for (int i = 0; i < 3; i++) {
    for (int j = 0; j < 3; j++) {
      cout << t.at(i, j) << " ";
    }
    cout << endl;
  }
  assert(t.sum(0, 0, 0, 1) == 11);
  assert(t.sum(0, 0, 1, 0) == 8);
  assert(t.sum(1, 1, 2, 2) == 29);
  t.set(500000000, 500000000, 100);
  assert(t.sum(0, 0, 1000000000, 1000000000) == 143);
  return 0;
}
\end{lstlisting}
\begin{exampleoutput}
Values:
5 6 0
3 5 5
0 5 14
\end{exampleoutput}

\section{Disjoint Sets and Tree Structures}
\setcounter{section}{7}
\setcounter{subsection}{0}
\subsection{Disjoint Set Forest (Simple)}
Maintain a set of elements partitioned into non-overlapping subsets. Each partition is assigned a unique representative known as the parent, or root. The following implements two well-known optimizations known as union-by-rank and path compression. This version is simplified to only work on integer elements.

\begin{compactitem}
  \item \inlinecode{initialize()} resets the data structure.
  \item \inlinecode{make\_set(u)} creates a new partition consisting of the single element \inlinecode{u}, which must not have been previously added to the data structure.
  \item \inlinecode{find\_root(u)} returns the unique representative of the partition containing \inlinecode{u}.
  \item \inlinecode{is\_united(u, v)} returns whether elements \inlinecode{u} and \inlinecode{v} belong to the same partition.
  \item \inlinecode{unite(u, v)} replaces the partitions containing \inlinecode{u} and \inlinecode{v} with a single new partition consisting of the union of elements in the original partitions.
\end{compactitem}

A precondition to the last three operations is that \inlinecode{make\_set()} must have been previously called on their arguments.

\Needspace*{3\baselineskip}
\textbf{Time Complexity}:
\begin{compactitem}
  \item $O(1)$ per call to \inlinecode{initialize()} and \inlinecode{make\_set()}.
  \item $O(\alpha\left(n\right))$ per call to \inlinecode{find\_root()}, \inlinecode{is\_united()}, and \inlinecode{unite()}, where $n$ is the number of elements that has been added via \inlinecode{make\_set()} so far, and $\alpha(n)$ is the extremely slow growing inverse of the Ackermann function (effectively a very small constant for all practical values of $n$).
\end{compactitem}

\Needspace*{3\baselineskip}
\textbf{Space Complexity}:
\begin{compactitem}
  \item $O(n)$ for storage of the disjoint set forest elements.
  \item $O(1)$ auxiliary for all operations.
\end{compactitem}

\begin{lstlisting}[language=C++]
#include <vector>

int num_sets;
std::vector<int> root, rank;

void initialize() {
  num_sets = 0;
  root.clear();
  rank.clear();
}

void make_set(int u) {
  if (u >= static_cast<int>(root.size())) {
    root.resize(u + 1);
    rank.resize(u + 1);
  }
  root[u] = u;
  rank[u] = 0;
  num_sets++;
}

int find_root(int u) {
  if (root[u] != u) {
    root[u] = find_root(root[u]);
  }
  return root[u];
}

bool is_united(int u, int v) {
  return find_root(u) == find_root(v);
}

void unite(int u, int v) {
  int ru = find_root(u), rv = find_root(v);
  if (ru == rv) {
    return;
  }
  num_sets--;
  if (rank[ru] < rank[rv]) {
    root[ru] = rv;
  } else {
    root[rv] = ru;
    if (rank[ru] == rank[rv]) {
      rank[ru]++;
    }
  }
}
\end{lstlisting}
\Needspace*{4\baselineskip}
\textbf{Example Usage}\par
\begin{lstlisting}[language=C++]
#include <cassert>

int main() {
  initialize();
  for (char c = 'a'; c <= 'g'; c++) {
    make_set(c);
  }
  unite('a', 'b');
  unite('b', 'f');
  unite('d', 'e');
  unite('d', 'g');
  assert(num_sets == 3);
  assert(is_united('a', 'b'));
  assert(!is_united('a', 'c'));
  assert(!is_united('b', 'g'));
  assert(is_united('e', 'g'));
  return 0;
}
\end{lstlisting}
\subsection{Disjoint Set Forest (Compressed)}
Maintain a set of elements partitioned into non-overlapping subsets using a collection of trees. Each partition is assigned a unique representative known as the parent, or root. The following implements two well-known optimizations known as union-by-rank and path compression. This version uses an \inlinecode{std::unordered\_map} for storage and coordinate compression (thus, element types must meet the requirements of key types for \inlinecode{std::unordered\_map}).

\begin{compactitem}
  \item \inlinecode{make\_set(u)} creates a new partition consisting of the single element \inlinecode{u}, which must not have been previously added to the data structure.
  \item \inlinecode{is\_united(u, v)} returns whether elements \inlinecode{u} and \inlinecode{v} belong to the same partition.
  \item \inlinecode{unite(u, v)} replaces the partitions containing \inlinecode{u} and \inlinecode{v} with a single new partition consisting of the union of elements in the original partitions.
  \item \inlinecode{get\_all\_sets()} returns all current partitions as a vector of vectors.
\end{compactitem}

A precondition to the last three operations is that \inlinecode{make\_set()} must have been previously called on their arguments.

\Needspace*{3\baselineskip}
\textbf{Time Complexity}:
\begin{compactitem}
  \item $O(1)$ per call to the constructor.
  \item $O(\log n)$ per call to \inlinecode{make\_set()}, where $n$ is the number of elements that have been added via \inlinecode{make\_set()} so far.
  \item $O(\alpha\left(n\right) \log n)$ per call to \inlinecode{is\_united()} and \inlinecode{unite()}, where $n$ is the number of elements that have been added via \inlinecode{make\_set()} so far, and $\alpha(n)$ is the extremely slow growing inverse of the Ackermann function (effectively a very small constant for all practical values of $n$).
  \item $O(n)$ per call to \inlinecode{get\_all\_sets()}.
\end{compactitem}

\Needspace*{3\baselineskip}
\textbf{Space Complexity}:
\begin{compactitem}
  \item $O(n)$ for storage of the disjoint set forest elements.
  \item $O(n)$ auxiliary heap space for \inlinecode{get\_all\_sets()}.
  \item $O(1)$ auxiliary for all other operations.
\end{compactitem}

\begin{lstlisting}[language=C++]
#include <unordered_map>
#include <vector>

template<class T>
class DisjointSetForest {
  int num_elements, num_sets;
  std::unordered_map<T, int> id;
  std::vector<int> root, rank;

  int find_root(int u) {
    if (root[u] != u) {
      root[u] = find_root(root[u]);
    }
    return root[u];
  }

 public:
  DisjointSetForest() : num_elements(0), num_sets(0) {}

  int size() const { return num_elements; }
  int sets() const { return num_sets; }

  void make_set(const T &u) {
    if (id.find(u) != id.end()) {
      return;
    }
    id[u] = num_elements;
    root.push_back(num_elements++);
    rank.push_back(0);
    num_sets++;
  }

  bool is_united(const T &u, const T &v) { return find_root(id[u]) == find_root(id[v]); }

  void unite(const T &u, const T &v) {
    int ru = find_root(id[u]), rv = find_root(id[v]);
    if (ru == rv) {
      return;
    }
    num_sets--;
    if (rank[ru] < rank[rv]) {
      root[ru] = rv;
    } else {
      root[rv] = ru;
      if (rank[ru] == rank[rv]) {
        rank[ru]++;
      }
    }
  }

  std::vector<std::vector<T>> get_all_sets() {
    std::unordered_map<int, std::vector<T>> tmp;
    for (auto &[key, val] : id) {
      tmp[find_root(val)].push_back(key);
    }
    std::vector<std::vector<T>> res;
    for (auto &[key, val] : tmp) {
      res.push_back(val);
    }
    return res;
  }
};
\end{lstlisting}
\Needspace*{4\baselineskip}
\textbf{Example Usage}\par
\begin{lstlisting}[language=C++]
#include <cassert>
#include <iostream>
using namespace std;

int main() {
  DisjointSetForest<char> dsf;
  for (char c = 'a'; c <= 'g'; c++) {
    dsf.make_set(c);
  }
  dsf.unite('a', 'b');
  dsf.unite('b', 'f');
  dsf.unite('d', 'e');
  dsf.unite('d', 'g');
  assert(dsf.size() == 7);
  assert(dsf.sets() == 3);
  auto s = dsf.get_all_sets();
  bool first_set = true;
  for (const auto &set : s) {
    cout << (first_set ? "[" : ", [");
    first_set = false;
    bool first_value = true;
    for (char value : set) {
      cout << (first_value ? "" : ", ") << value;
      first_value = false;
    }
    cout << "]";
  }
  cout << endl;
  return 0;
}
\end{lstlisting}
\begin{exampleoutput}
[a, b, f], [c], [d, e, g]
\end{exampleoutput}
\subsection{Disjoint Set Forest (Rollback)}
Maintains disjoint sets while supporting rollback to previous snapshots. Rollback DSU is useful for offline dynamic connectivity, divide-and-conquer over time, and backtracking search where unions must be undone.

Path compression is intentionally not used, because it is hard to undo. Union by size keeps tree height logarithmic, and every successful union records enough information to restore the previous state.

\begin{compactitem}
  \item \inlinecode{initialize(n)} creates singleton sets over elements \inlinecode{0} through \inlinecode{n - 1}.
  \item \inlinecode{find\_root(u)} returns the representative of the set containing \inlinecode{u}.
  \item \inlinecode{is\_united(u, v)} returns whether \inlinecode{u} and \inlinecode{v} are in the same set.
  \item \inlinecode{unite(u, v)} merges two sets and returns whether a merge occurred.
  \item \inlinecode{snapshot()} returns a token representing the current history size.
  \item \inlinecode{rollback(s)} undoes all changes made after snapshot \inlinecode{s}.
\end{compactitem}

\Needspace*{3\baselineskip}
\textbf{Time Complexity}:
\begin{compactitem}
  \item $O(n)$ per call to \inlinecode{initialize(n)}.
  \item $O(\log n)$ worst-case per call to \inlinecode{find\_root(u)}, \inlinecode{is\_united(u, v)}, and \inlinecode{unite(u, v)}.
  \item $O(1)$ per undone union during \inlinecode{rollback(s)}.
\end{compactitem}

\textbf{Space Complexity}: $O(n + q)$ for $q$ successful unions stored in the rollback history.

\begin{lstlisting}[language=C++]
#include <numeric>
#include <vector>

class RollbackDSU {
  struct Change {
    int child, parent, parent_size;

    Change(int child = -1, int parent = -1, int parent_size = 0)
        : child(child), parent(parent), parent_size(parent_size) {}
  };

  std::vector<int> root, size;
  std::vector<Change> history;
  int sets;

 public:
  RollbackDSU() : sets(0) {}

  explicit RollbackDSU(int n) { initialize(n); }

  void initialize(int n) {
    root.resize(n);
    size.assign(n, 1);
    history.clear();
    sets = n;
    std::iota(root.begin(), root.end(), 0);
  }

  int count_sets() const { return sets; }

  int find_root(int u) const {
    while (root[u] != u) {
      u = root[u];
    }
    return u;
  }

  bool is_united(int u, int v) const { return find_root(u) == find_root(v); }

  bool unite(int u, int v) {
    u = find_root(u);
    v = find_root(v);
    if (u == v) {
      return false;
    }
    if (size[u] > size[v]) {
      int tmp = u;
      u = v;
      v = tmp;
    }
    history.emplace_back(u, v, size[v]);
    root[u] = v;
    size[v] += size[u];
    sets--;
    return true;
  }

  int snapshot() const { return history.size(); }

  void rollback(int snapshot) {
    while (static_cast<int>(history.size()) > snapshot) {
      Change c = history.back();
      history.pop_back();
      root[c.child] = c.child;
      size[c.parent] = c.parent_size;
      sets++;
    }
  }
};
\end{lstlisting}
\Needspace*{4\baselineskip}
\textbf{Example Usage}\par
\begin{lstlisting}[language=C++]
#include <cassert>

int main() {
  RollbackDSU dsu(5);
  dsu.unite(0, 1);
  int s = dsu.snapshot();
  dsu.unite(1, 2);
  dsu.unite(3, 4);
  assert(dsu.is_united(0, 2));
  assert(dsu.count_sets() == 2);
  dsu.rollback(s);
  assert(dsu.is_united(0, 1));
  assert(!dsu.is_united(0, 2));
  assert(!dsu.is_united(3, 4));
  assert(dsu.count_sets() == 4);
  return 0;
}
\end{lstlisting}
\subsection{Lowest Common Ancestor (Sparse Table)}
Given a tree, determine the lowest common ancestor of any two nodes in the tree. The lowest common ancestor of two nodes $u$ and $v$ is the node that has the longest distance from the root while having both $u$ and $v$ as its descendant. A nodes is considered to be a descendant of itself. \inlinecode{build()} applies to a global, pre-populated adjacency list \inlinecode{adj} which must only consist of nodes numbered with integers between 0 (inclusive) and the total number of nodes (exclusive), as passed in the function argument.

\Needspace*{3\baselineskip}
\textbf{Time Complexity}:
\begin{compactitem}
  \item $O(n \log n)$ per call to \inlinecode{build()}, where $n$ is the number of nodes.
  \item $O(\log n)$ per call to \inlinecode{lca()}.
\end{compactitem}

\Needspace*{3\baselineskip}
\textbf{Space Complexity}:
\begin{compactitem}
  \item $O(n \log n)$ to store the sparse table, where $n$ is the number of nodes.
  \item $O(n)$ auxiliary stack space for \inlinecode{build()}.
  \item $O(1)$ auxiliary for \inlinecode{lca()}.
\end{compactitem}

\begin{lstlisting}[language=C++]
#include <vector>

std::vector<std::vector<int>> adj, dp;
int len, timer;
std::vector<int> tin, tout;

void dfs(int u, int p) {
  tin[u] = timer++;
  dp[u][0] = p;
  for (int i = 1; i < len; i++) {
    dp[u][i] = dp[dp[u][i - 1]][i - 1];
  }
  for (int v : adj[u]) {
    if (v != p) {
      dfs(v, u);
    }
  }
  tout[u] = timer++;
}

void build(int nodes, int root = 0) {
  len = 1;
  while ((1 << len) <= nodes) {
    len++;
  }
  dp.assign(nodes, std::vector<int>(len));
  tin.assign(nodes, 0);
  tout.assign(nodes, 0);
  timer = 0;
  dfs(root, root);
}

bool is_ancestor(int parent, int child) {
  return (tin[parent] <= tin[child]) && (tout[child] <= tout[parent]);
}

int lca(int u, int v) {
  if (is_ancestor(u, v)) {
    return u;
  }
  if (is_ancestor(v, u)) {
    return v;
  }
  for (int i = len - 1; i >= 0; i--) {
    if (!is_ancestor(dp[u][i], v)) {
      u = dp[u][i];
    }
  }
  return dp[u][0];
}
\end{lstlisting}
\Needspace*{4\baselineskip}
\textbf{Example Usage}\par
\begin{lstlisting}[language=C++]
#include <cassert>
using namespace std;

int main() {
  int nodes = 5;
  adj.resize(nodes);
  adj[0].push_back(1);
  adj[1].push_back(0);
  adj[1].push_back(2);
  adj[2].push_back(1);
  adj[3].push_back(1);
  adj[1].push_back(3);
  adj[0].push_back(4);
  adj[4].push_back(0);
  build(nodes, 0);
  assert(lca(3, 2) == 1);
  assert(lca(2, 4) == 0);
  return 0;
}
\end{lstlisting}
\subsection{Lowest Common Ancestor (Segment Tree)}
Given a tree, determine the lowest common ancestor of any two nodes in the tree. The lowest common ancestor of two nodes $u$ and $v$ is the node that has the longest distance from the root while having both $u$ and $v$ as its descendant. A nodes is considered to be a descendant of itself. \inlinecode{build()} applies to a global, pre-populated adjacency list \inlinecode{adj} which must only consist of nodes numbered with integers between 0 (inclusive) and the total number of nodes (exclusive), as passed in the function argument.

\Needspace*{3\baselineskip}
\textbf{Time Complexity}:
\begin{compactitem}
  \item $O(n \log n)$ per call to \inlinecode{build()}, where $n$ is the number of nodes.
  \item $O(\log n)$ per call to \inlinecode{lca()}.
\end{compactitem}

\Needspace*{3\baselineskip}
\textbf{Space Complexity}:
\begin{compactitem}
  \item $O(n)$ for storage of the segment tree, where $n$ is the number of nodes.
  \item $O(n \log n)$ auxiliary stack space for \inlinecode{build()}.
  \item $O(\log n)$ auxiliary stack space for \inlinecode{lca()}.
\end{compactitem}

\begin{lstlisting}[language=C++]
#include <algorithm>
#include <vector>

std::vector<std::vector<int>> adj;
int len, counter;
std::vector<int> depth, dfs_order, first, minpos;

void dfs(int u, int d) {
  depth[u] = d;
  dfs_order[counter++] = u;
  for (int v : adj[u]) {
    if (depth[v] == -1) {
      dfs(v, d + 1);
      dfs_order[counter++] = u;
    }
  }
}

void build(int n, int lo, int hi) {
  if (lo == hi) {
    minpos[n] = dfs_order[lo];
    return;
  }
  int lchild = 2 * n + 1, rchild = 2 * n + 2, mid = lo + (hi - lo) / 2;
  build(lchild, lo, mid);
  build(rchild, mid + 1, hi);
  minpos[n] = depth[minpos[lchild]] < depth[minpos[rchild]] ? minpos[lchild] : minpos[rchild];
}

void build(int nodes, int root) {
  depth.assign(nodes, -1);
  dfs_order.assign(2 * nodes, 0);
  first.assign(nodes, -1);
  minpos.assign(8 * nodes, 0);
  len = 2 * nodes - 1;
  counter = 0;
  dfs(root, 0);
  build(0, 0, len - 1);
  for (int i = 0; i < len; i++) {
    if (first[dfs_order[i]] == -1) {
      first[dfs_order[i]] = i;
    }
  }
}

int get_minpos(int a, int b, int n, int lo, int hi) {
  if (a == lo && b == hi) {
    return minpos[n];
  }
  int mid = lo + (hi - lo) / 2;
  if (a <= mid && mid < b) {
    int p1 = get_minpos(a, std::min(b, mid), 2 * n + 1, lo, mid);
    int p2 = get_minpos(std::max(a, mid + 1), b, 2 * n + 2, mid + 1, hi);
    return depth[p1] < depth[p2] ? p1 : p2;
  }
  if (a <= mid) {
    return get_minpos(a, std::min(b, mid), 2 * n + 1, lo, mid);
  }
  return get_minpos(std::max(a, mid + 1), b, 2 * n + 2, mid + 1, hi);
}

int lca(int u, int v) {
  return get_minpos(std::min(first[u], first[v]), std::max(first[u], first[v]), 0, 0, len - 1);
}
\end{lstlisting}
\Needspace*{4\baselineskip}
\textbf{Example Usage}\par
\begin{lstlisting}[language=C++]
#include <cassert>
using namespace std;

int main() {
  int nodes = 5;
  adj.resize(nodes);
  adj[0].push_back(1);
  adj[1].push_back(0);
  adj[1].push_back(2);
  adj[2].push_back(1);
  adj[3].push_back(1);
  adj[1].push_back(3);
  adj[0].push_back(4);
  adj[4].push_back(0);
  build(nodes, 0);
  assert(lca(3, 2) == 1);
  assert(lca(2, 4) == 0);
  return 0;
}
\end{lstlisting}
\subsection{Heavy Light Decomposition}
Maintain a tree with values associated with either its edges or nodes, while supporting both dynamic queries and dynamic updates of all values on a given path between two nodes in the tree. Heavy-light decomposition partitions the nodes of the tree into disjoint paths where all nodes have degree two, except the endpoints of a path which has degree one.

The query operation is defined by an associative \inlinecode{combine()} function which satisfies \inlinecode{combine(x, combine(y, z)) = combine(combine(x, y), z)} for all values \inlinecode{x}, \inlinecode{y}, and \inlinecode{z} in the tree. The default code below assumes a numerical tree type, defining queries for the ``min'' of the target range. Another possible query operation is ``sum'', in which case \inlinecode{combine(a, b)} should return $a + b$. For direction-independent path queries, \inlinecode{combine()} should also be commutative; otherwise, store enough information in each aggregate to combine paths in the required order.

The update operation is defined by \inlinecode{apply\_delta()} and \inlinecode{compose\_deltas()}. A delta must act on an aggregate summary of a path of length \inlinecode{len}: \inlinecode{apply\_delta(v, d, len)} returns the aggregate after applying update \inlinecode{d} to every element represented by aggregate \inlinecode{v}. Pending deltas are combined in chronological order by \inlinecode{compose\_deltas(old, d)}, meaning ``apply \inlinecode{old}, then apply \inlinecode{d}''. These hooks do not support arbitrary query/update pairings; the delta operation must distribute over \inlinecode{combine()}, and composed deltas must have the same effect as applying their updates sequentially. The default code below defines updates that ``set'' a path's edges or nodes to a new value. For range increment updates, \inlinecode{apply\_delta(v, d, len)} would return \inlinecode{v + d} for min/max queries, or \inlinecode{v + d * len} for sum queries, and \inlinecode{compose\_deltas(old, d)} would return \inlinecode{old + d}.

\begin{compactitem}
  \item \inlinecode{HeavyLight(adj, v)} constructs a new heavy light decomposition on a tree defined by the adjacency list \inlinecode{adj}, with all values initialized to \inlinecode{v}. The adjacency list must consist of only the integers from 0 to \inlinecode{adj.size() - 1}, inclusive. No duplicate edges should exist, and the graph must be connected.
  \item \inlinecode{query(u, v)} returns the result of \inlinecode{combine()} applied to all values on the path from node \inlinecode{u} to node \inlinecode{v}.
  \item \inlinecode{update(u, v, d)} modifies all values on the path from node \inlinecode{u} to node \inlinecode{v} by respectively applying the delta \inlinecode{d}.
\end{compactitem}

\Needspace*{3\baselineskip}
\textbf{Time Complexity}:
\begin{compactitem}
  \item $O(n)$ per call to the constructor, where $n$ is the number of nodes.
  \item $O(\log^{2} n)$ per call to \inlinecode{query()} and \inlinecode{update()};
\end{compactitem}

\Needspace*{3\baselineskip}
\textbf{Space Complexity}:
\begin{compactitem}
  \item $O(n)$ for storage of the decomposition.
  \item $O(n)$ auxiliary stack space for the constructor.
  \item $O(1)$ auxiliary for \inlinecode{query()} and \inlinecode{update()}.
\end{compactitem}

\begin{lstlisting}[language=C++]
#include <algorithm>
#include <stdexcept>
#include <vector>

template<class T>
class HeavyLight {
  // Set this to true to store values on edges, false to store values on nodes.
  static const bool VALUES_ON_EDGES = true;

  static T combine(const T &a, const T &b) { return std::min(a, b); }
  static T apply_delta(const T &v, const T &d, int len) { return d; }
  static T compose_deltas(const T &old, const T &d) { return d; }

  int counter, paths;
  std::vector<std::vector<T>> value, delta;
  std::vector<std::vector<bool>> pending;
  std::vector<std::vector<int>> len;
  std::vector<int> size, parent, tin, tout, path, pathlen, pathpos, pathroot;
  const std::vector<std::vector<int>> &adj;

  void dfs(int u, int p) {
    tin[u] = counter++;
    parent[u] = p;
    size[u] = 1;
    for (int v : adj[u]) {
      if (v != p) {
        dfs(v, u);
        size[u] += size[v];
      }
    }
    tout[u] = counter++;
  }

  int new_path(int u) {
    pathroot[paths] = u;
    return paths++;
  }

  void build_paths(int u, int path) {
    this->path[u] = path;
    pathpos[u] = pathlen[path]++;
    for (int v : adj[u]) {
      if (v != parent[u]) {
        build_paths(v, (2 * size[v] >= size[u]) ? path : new_path(v));
      }
    }
  }

  inline T applied_value(int path, int i) {
    return pending[path][i] ? apply_delta(value[path][i], delta[path][i], len[path][i])
                            : value[path][i];
  }

  void push_delta(int path, int i) {
    int d = 0;
    while ((i >> d) > 0) {
      d++;
    }
    for (d -= 2; d >= 0; d--) {
      int l = (i >> d), r = (l ^ 1), n = l / 2;
      if (pending[path][n]) {
        value[path][n] = applied_value(path, n);
        delta[path][l] =
            pending[path][l] ? compose_deltas(delta[path][l], delta[path][n]) : delta[path][n];
        delta[path][r] =
            pending[path][r] ? compose_deltas(delta[path][r], delta[path][n]) : delta[path][n];
        pending[path][l] = pending[path][r] = true;
        pending[path][n] = false;
      }
    }
  }

  bool query(int path, int u, int v, T *res) {
    push_delta(path, u += value[path].size() / 2);
    push_delta(path, v += value[path].size() / 2);
    bool found = false;
    for (; u <= v; u = (u + 1) / 2, v = (v - 1) / 2) {
      if ((u & 1) != 0) {
        T value = applied_value(path, u);
        *res = found ? combine(*res, value) : value;
        found = true;
      }
      if ((v & 1) == 0) {
        T value = applied_value(path, v);
        *res = found ? combine(*res, value) : value;
        found = true;
      }
    }
    return found;
  }

  void update(int path, int u, int v, const T &d) {
    push_delta(path, u += value[path].size() / 2);
    push_delta(path, v += value[path].size() / 2);
    int tu = -1, tv = -1;
    for (; u <= v; u = (u + 1) / 2, v = (v - 1) / 2) {
      if ((u & 1) != 0) {
        delta[path][u] = pending[path][u] ? compose_deltas(delta[path][u], d) : d;
        pending[path][u] = true;
        if (tu == -1) {
          tu = u;
        }
      }
      if ((v & 1) == 0) {
        delta[path][v] = pending[path][v] ? compose_deltas(delta[path][v], d) : d;
        pending[path][v] = true;
        if (tv == -1) {
          tv = v;
        }
      }
    }
    for (int i = tu; i > 1; i /= 2) {
      value[path][i / 2] = combine(applied_value(path, i), applied_value(path, i ^ 1));
    }
    for (int i = tv; i > 1; i /= 2) {
      value[path][i / 2] = combine(applied_value(path, i), applied_value(path, i ^ 1));
    }
  }

  inline bool is_ancestor(int parent, int child) {
    return (tin[parent] <= tin[child]) && (tout[child] <= tout[parent]);
  }

 public:
  explicit HeavyLight(const std::vector<std::vector<int>> &adj, const T &v = T())
      : counter(0),
        paths(0),
        size(adj.size()),
        parent(adj.size()),
        tin(adj.size()),
        tout(adj.size()),
        path(adj.size()),
        pathlen(adj.size()),
        pathpos(adj.size()),
        pathroot(adj.size()),
        adj(adj) {
    dfs(0, -1);
    build_paths(0, new_path(0));
    value.resize(paths);
    delta.resize(paths);
    pending.resize(paths);
    len.resize(paths);
    for (int i = 0; i < paths; i++) {
      int m = pathlen[i];
      value[i].assign(2 * m, v);
      delta[i].resize(2 * m);
      pending[i].assign(2 * m, false);
      len[i].assign(2 * m, 1);
      for (int j = 2 * m - 1; j > 1; j -= 2) {
        value[i][j / 2] = combine(value[i][j], value[i][j ^ 1]);
        len[i][j / 2] = len[i][j] + len[i][j ^ 1];
      }
    }
  }

  T query(int u, int v) {
    if (VALUES_ON_EDGES && u == v) {
      throw std::runtime_error("No edge between u and v to be queried.");
    }
    bool found = false;
    T res = T(), value;
    int root;
    while (!is_ancestor(root = pathroot[path[u]], v)) {
      if (query(path[u], 0, pathpos[u], &value)) {
        res = found ? combine(res, value) : value;
        found = true;
      }
      u = parent[root];
    }
    while (!is_ancestor(root = pathroot[path[v]], u)) {
      if (query(path[v], 0, pathpos[v], &value)) {
        res = found ? combine(res, value) : value;
        found = true;
      }
      v = parent[root];
    }
    if (query(
            path[u], std::min(pathpos[u], pathpos[v]) + static_cast<int>(VALUES_ON_EDGES),
            std::max(pathpos[u], pathpos[v]), &value
        )) {
      res = found ? combine(res, value) : value;
      found = true;
    }
    if (!found) {
      throw std::runtime_error("Unexpected error: No values found.");
    }
    return res;
  }

  void update(int u, int v, const T &d) {
    if (VALUES_ON_EDGES && u == v) {
      return;
    }
    int root;
    while (!is_ancestor(root = pathroot[path[u]], v)) {
      update(path[u], 0, pathpos[u], d);
      u = parent[root];
    }
    while (!is_ancestor(root = pathroot[path[v]], u)) {
      update(path[v], 0, pathpos[v], d);
      v = parent[root];
    }
    update(
        path[u], std::min(pathpos[u], pathpos[v]) + static_cast<int>(VALUES_ON_EDGES),
        std::max(pathpos[u], pathpos[v]), d
    );
  }
};
\end{lstlisting}
\Needspace*{4\baselineskip}
\textbf{Example Usage}\par
\begin{lstlisting}[language=C++]
#include <cassert>
using namespace std;

int main() {
  //     w=40      w=20      w=10
  // 0---------1---------2---------3
  //                     |
  //                     ----------4
  //                         w=30
  vector<vector<int>> adj(5);
  adj[0].push_back(1);
  adj[1].push_back(0);
  adj[1].push_back(2);
  adj[2].push_back(1);
  adj[2].push_back(3);
  adj[3].push_back(2);
  adj[2].push_back(4);
  adj[4].push_back(2);
  HeavyLight<int> hld(adj, 0);
  hld.update(0, 1, 40);
  hld.update(1, 2, 20);
  hld.update(2, 3, 10);
  hld.update(2, 4, 30);
  assert(hld.query(0, 3) == 10);
  assert(hld.query(2, 4) == 30);
  hld.update(3, 4, 5);
  assert(hld.query(1, 4) == 5);
  return 0;
}
\end{lstlisting}
\subsection{Link-Cut Tree}
Maintain a forest of trees with values associated with its nodes, while supporting both dynamic queries and dynamic updates of all values on any path between two nodes in a given tree. In addition, support testing of whether two nodes are connected in the forest, as well as the merging and spliting of trees by adding or removing specific edges. Link/cut forests divide each of its trees into vertex-disjoint paths, each represented by a splay tree.

The query operation is defined by an associative \inlinecode{combine()} function which satisfies \inlinecode{combine(x, combine(y, z)) = combine(combine(x, y), z)} for all values \inlinecode{x}, \inlinecode{y}, and \inlinecode{z} in the forest. The default code below assumes a numerical forest type, defining queries for the ``min'' of the target range. Another possible query operation is ``sum'', in which case \inlinecode{combine(a, b)} should return $a + b$. For direction-independent path queries, \inlinecode{combine()} should also be commutative; otherwise, store enough information in each aggregate to combine paths in the required order.

The update operation is defined by \inlinecode{apply\_delta()} and \inlinecode{compose\_deltas()}. A delta must act on an aggregate summary of a path of length \inlinecode{len}: \inlinecode{apply\_delta(v, d, len)} returns the aggregate after applying update \inlinecode{d} to every element represented by aggregate \inlinecode{v}. Pending deltas are combined in chronological order by \inlinecode{compose\_deltas(old, d)}, meaning ``apply \inlinecode{old}, then apply \inlinecode{d}''. These hooks do not support arbitrary query/update pairings; the delta operation must distribute over \inlinecode{combine()}, and composed deltas must have the same effect as applying their updates sequentially. The default code below defines updates that ``set'' a path's nodes to a new value. For range increment updates, \inlinecode{apply\_delta(v, d, len)} would return \inlinecode{v + d} for min/max queries, or \inlinecode{v + d * len} for sum queries, and \inlinecode{compose\_deltas(old, d)} would return \inlinecode{old + d}.

\begin{compactitem}
  \item \inlinecode{LinkCutForest()} constructs an empty forest with no trees.
  \item \inlinecode{size()} returns the number of nodes in the forest.
  \item \inlinecode{trees()} returns the number of trees in the forest.
  \item \inlinecode{make\_root(i, v)} creates a new tree in the forest consisting of a single node labeled with the integer \inlinecode{i} and value initialized to \inlinecode{v}.
  \item \inlinecode{is\_connected(a, b)} returns whether nodes \inlinecode{a} and \inlinecode{b} are connected.
  \item \inlinecode{link(a, b)} adds an edge between the nodes \inlinecode{a} and \inlinecode{b}, both of which must exist and not be connected.
  \item \inlinecode{cut(a, b)} removes the edge between the nodes \inlinecode{a} and \inlinecode{b}, both of which must exist and be connected.
  \item \inlinecode{query(a, b)} returns the result of \inlinecode{combine()} applied to all values on the path from the node \inlinecode{a} to node \inlinecode{b}.
  \item \inlinecode{update(a, b, d)} modifies all the values on the path from node \inlinecode{a} to node \inlinecode{b} by respectively applying the delta \inlinecode{d}.
\end{compactitem}

\Needspace*{3\baselineskip}
\textbf{Time Complexity}:
\begin{compactitem}
  \item $O(1)$ per call to the constructor, \inlinecode{size()}, and \inlinecode{trees()}.
  \item $O(\log n)$ amortized per call to all other operations, where $n$ is the number of nodes.
\end{compactitem}

\Needspace*{3\baselineskip}
\textbf{Space Complexity}:
\begin{compactitem}
  \item $O(n)$ for storage of the forest, where $n$ is the number of nodes.
  \item $O(1)$ auxiliary for all operations.
\end{compactitem}

\begin{lstlisting}[language=C++]
#include <algorithm>
#include <cstddef>
#include <stdexcept>
#include <unordered_map>

template<class T>
class LinkCutForest {
  static T combine(const T &a, const T &b) { return std::min(a, b); }
  static T apply_delta(const T &v, const T &d, int len) { return d; }
  static T compose_deltas(const T &old, const T &d) { return d; }

  struct Node {
    T value, subtree_value, delta;
    int size;
    bool rev, pending;
    Node *left, *right, *parent;

    explicit Node(const T &v)
        : value(v),
          subtree_value(v),
          size(1),
          rev(false),
          pending(false),
          left(nullptr),
          right(nullptr),
          parent(nullptr) {}

    inline bool is_root() const {
      return parent == nullptr || (parent->left != this && parent->right != this);
    }

    inline T get_subtree_value() const {
      return pending ? apply_delta(subtree_value, delta, size) : subtree_value;
    }

    void push() {
      if (rev) {
        rev = false;
        std::swap(left, right);
        if (left != nullptr) {
          left->rev = !left->rev;
        }
        if (right != nullptr) {
          right->rev = !right->rev;
        }
      }
      if (pending) {
        value = apply_delta(value, delta, 1);
        subtree_value = apply_delta(subtree_value, delta, size);
        if (left != nullptr) {
          left->delta = left->pending ? compose_deltas(left->delta, delta) : delta;
          left->pending = true;
        }
        if (right != nullptr) {
          right->delta = right->pending ? compose_deltas(right->delta, delta) : delta;
          right->pending = true;
        }
        pending = false;
      }
    }

    void update() {
      size = 1;
      subtree_value = value;
      if (left != nullptr) {
        subtree_value = combine(subtree_value, left->get_subtree_value());
        size += left->size;
      }
      if (right != nullptr) {
        subtree_value = combine(subtree_value, right->get_subtree_value());
        size += right->size;
      }
    }
  };

  int num_trees;
  std::unordered_map<int, Node *> nodes;

  static void connect(Node *child, Node *parent, bool is_left) {
    if (child != nullptr) {
      child->parent = parent;
    }
    if (is_left) {
      parent->left = child;
    } else {
      parent->right = child;
    }
  }

  static void rotate(Node *n) {
    Node *parent = n->parent, *grandparent = parent->parent;
    bool parent_is_root = parent->is_root(), is_left = (n == parent->left);
    connect(is_left ? n->right : n->left, parent, is_left);
    connect(parent, n, !is_left);
    if (parent_is_root) {
      if (n != nullptr) {
        n->parent = grandparent;
      }
    } else {
      connect(n, grandparent, parent == grandparent->left);
    }
    parent->update();
  }

  static void splay(Node *n) {
    while (!n->is_root()) {
      Node *parent = n->parent, *grandparent = parent->parent;
      if (!parent->is_root()) {
        grandparent->push();
      }
      parent->push();
      n->push();
      if (!parent->is_root()) {
        if ((n == parent->left) == (parent == grandparent->left)) {
          rotate(parent);
        } else {
          rotate(n);
        }
      }
      rotate(n);
    }
    n->push();
    n->update();
  }

  static Node *expose(Node *n) {
    Node *prev = nullptr;
    for (Node *curr = n; curr != nullptr; curr = curr->parent) {
      splay(curr);
      curr->left = prev;
      prev = curr;
    }
    splay(n);
    return prev;
  }

  // Helper variables.
  Node *u, *v;

  void get_uv(int a, int b) {
    auto it1 = nodes.find(a);
    auto it2 = nodes.find(b);
    if (it1 == nodes.end() || it2 == nodes.end()) {
      throw std::runtime_error("Queried node ID does not exist in forest.");
    }
    u = it1->second;
    v = it2->second;
  }

 public:
  LinkCutForest() : num_trees(0) {}

  ~LinkCutForest() {
    for (auto &[key, node] : nodes) {
      delete node;
    }
  }

  int size() const { return nodes.size(); }
  int trees() const { return num_trees; }

  void make_root(int i, const T &v = T()) {
    if (auto it = nodes.find(i); it != nodes.end()) {
      throw std::runtime_error("Cannot make a root with an existing ID.");
    }
    Node *n = new Node(v);
    expose(n);
    n->rev = !n->rev;
    nodes[i] = n;
    num_trees++;
  }

  bool is_connected(int a, int b) {
    get_uv(a, b);
    if (a == b) {
      return true;
    }
    expose(u);
    expose(v);
    return u->parent != nullptr;
  }

  void link(int a, int b) {
    if (is_connected(a, b)) {
      throw std::runtime_error("Cannot link nodes that are already connected.");
    }
    get_uv(a, b);
    expose(u);
    u->rev = !u->rev;
    u->parent = v;
    num_trees--;
  }

  void cut(int a, int b) {
    get_uv(a, b);
    expose(u);
    u->rev = !u->rev;
    expose(v);
    if (v->right != u || u->left != nullptr) {
      throw std::runtime_error("Cannot cut edge that does not exist.");
    }
    v->right->parent = nullptr;
    v->right = nullptr;
    num_trees++;
  }

  T query(int a, int b) {
    if (!is_connected(a, b)) {
      throw std::runtime_error("Cannot query nodes that are not connected.");
    }
    get_uv(a, b);
    expose(u);
    u->rev = !u->rev;
    expose(v);
    return v->get_subtree_value();
  }

  void update(int a, int b, const T &d) {
    if (!is_connected(a, b)) {
      throw std::runtime_error("Cannot update nodes that are not connected.");
    }
    get_uv(a, b);
    expose(u);
    u->rev = !u->rev;
    expose(v);
    v->delta = v->pending ? compose_deltas(v->delta, d) : d;
    v->pending = true;
  }
};
\end{lstlisting}
\Needspace*{4\baselineskip}
\textbf{Example Usage}\par
\begin{lstlisting}[language=C++]
#include <cassert>
using namespace std;

int main() {
  LinkCutForest<int> lcf;
  lcf.make_root(0, 10);
  lcf.make_root(1, 40);
  lcf.make_root(2, 20);
  lcf.make_root(3, 10);
  lcf.make_root(4, 30);
  assert(lcf.size() == 5);
  assert(lcf.trees() == 5);
  lcf.link(0, 1);
  lcf.link(1, 2);
  lcf.link(2, 3);
  lcf.link(2, 4);
  assert(lcf.trees() == 1);

  // v=10      v=40      v=20      v=10
  //  0---------1---------2---------3
  //                      |
  //                      ----------4
  //                               v=30
  assert(lcf.query(1, 4) == 20);
  lcf.update(1, 1, 100);
  lcf.update(2, 4, 100);

  // v=10     v=100     v=100      v=10
  //  0---------1---------2---------3
  //                      |
  //                      ----------4
  //                              v=100
  assert(lcf.query(4, 4) == 100);
  assert(lcf.query(0, 4) == 10);
  assert(lcf.query(3, 4) == 10);
  lcf.cut(1, 2);

  // v=10     v=100     v=100      v=0
  //  0---------1         2---------3
  //                      |
  //                      ----------4
  //                              v=100
  assert(lcf.trees() == 2);
  assert(!lcf.is_connected(1, 2));
  assert(!lcf.is_connected(0, 4));
  assert(lcf.is_connected(2, 3));
  return 0;
}
\end{lstlisting}
